\documentclass[12pt,a4paper]{article}

\usepackage{fontspec}
\usepackage{xeCJK}
\usepackage{polyglossia}
\setmainlanguage{russian}
\setotherlanguage{english}
\defaultfontfeatures{Ligatures=TeX}
\IfFontExistsTF{Times New Roman}{%
  \setmainfont{Times New Roman}%
  \setsansfont{Arial}%
  \setmonofont{Courier New}%
}{%
  \setmainfont{Tinos}%
  \setsansfont{Liberation Sans}%
  \setmonofont{DejaVu Sans Mono}%
}
\setCJKmainfont{Noto Serif CJK SC}
\XeTeXlinebreaklocale "zh"
\XeTeXlinebreakskip=0pt plus 1pt

\usepackage{amsmath,amssymb,amsthm,mathtools}
\usepackage{geometry}
\geometry{left=20mm,right=20mm,top=20mm,bottom=22mm}
\usepackage{enumitem}
\setlist{itemsep=0.15em,topsep=0.25em,leftmargin=1.4em}
\usepackage{array,booktabs,longtable}
\usepackage{microtype}
\usepackage{titlesec}
\usepackage{fancyhdr}
\usepackage{xparse}
\usepackage{hyperref}
\hypersetup{colorlinks=true,linkcolor=black,urlcolor=blue}

\pagestyle{fancy}
\setlength{\headheight}{15pt}
\fancyhf{}
\lhead{Государственный экзамен 2026}
\rhead{Билеты 1--29: русский и китайский}
\cfoot{\thepage}

\setlength{\parindent}{0pt}
\setlength{\parskip}{0.45em}
\allowdisplaybreaks
\sloppy
\setcounter{secnumdepth}{0}
\setcounter{tocdepth}{1}
\titleformat{\section}{\Large\bfseries}{\thesection.}{0.5em}{}
\titleformat{\subsection}{\normalsize\normalfont}{\thesubsection.}{0.5em}{}
\titleformat{\subsubsection}{\normalsize\normalfont}{\thesubsubsection.}{0.5em}{}

\newcommand{\R}{\mathbb{R}}
\newcommand{\eps}{\varepsilon}
\newcommand{\T}{\top}
\newcommand{\dd}{\,d}
\newcommand{\norm}[1]{\left\lVert #1\right\rVert}
\newcommand{\inner}[2]{\left\langle #1,#2\right\rangle}
\newcommand{\sng}{\mathrm{sng}}
\DeclareMathOperator*{\argmin}{arg\,min}
\DeclareMathOperator*{\argmax}{arg\,max}
\newcommand{\Lie}{\operatorname{Lie}}
\newcommand{\Vect}{\operatorname{Vec}}
\newcommand{\rank}{\operatorname{rank}}

\newcommand{\ad}{\operatorname{ad}}
\newcommand{\sgn}{\operatorname{sgn}}
\newcommand{\PV}{\mathrm{PV}}
\newcommand{\CV}{\mathrm{CV}}
\newcommand{\Argmax}{\operatorname*{Arg\,max}}
\newcommand{\memory}{}
\newcommand{\official}[1]{}
\newcommand{\fullanswer}{}
\newcommand{\cnexplain}{}
\newcommand{\corrections}{}
\NewDocumentCommand{\ticket}{m m g}{%
  \clearpage
  \section*{Билет #1. #2}%
  \addcontentsline{toc}{section}{Билет #1. #2}%
  \IfNoValueF{#3}{{\large\bfseries 第#1题签。#3}\par\medskip}%
}
\NewDocumentCommand{\qpart}{m m g}{\par\bigskip\noindent{\Large\bfseries #1. #2}\par\IfNoValueF{#3}{{\Large\bfseries #1. #3}\par}\nobreak\medskip}
\newcommand{\plainheading}[2]{\par\bigskip\noindent{\large #1}\par\noindent{\large #2}\par\medskip}
\newcommand{\zh}[1]{\par\noindent #1\par}
\newcommand{\blocktitle}[3]{%
  \clearpage
  \thispagestyle{fancy}%
  \begin{center}
    {\LARGE\bfseries Блок #1. #2}\\[0.35em]
    {\large #3}
  \end{center}
  \addcontentsline{toc}{section}{Блок #1. #2}
  \vspace{0.5em}
}

\begin{document}
\begin{titlepage}
\centering
\vspace*{3cm}
{\Huge\bfseries Материалы к государственному экзамену\\[0.4em]кафедры оптимального управления ВМК МГУ}\\[1.2cm]
{\Large Билеты 1--29}\\[0.4cm]
{\Large 1--29 号题签中俄对照}\\[1.8cm]
{\large 2026}
\vfill
\end{titlepage}

\tableofcontents

\blocktitle{A}{Обратные задачи и метод динамической регуляризации}{Билеты 1--3}
\zh{A板块：反问题与动态正则化方法。第1--3题签。}

Общее замечание. Билеты 1--3 относятся к обратным задачам восстановления неизвестных характеристик динамической системы по наблюдениям её движения. Общая схема:
\zh{总说明。第1--3题签讨论的是反问题：根据对动力系统运动的观测来恢复系统的未知特性。整体思路如下：}
\begin{center}
\begin{minipage}{0.92\textwidth}
\centering
некорректная обратная задача $\Rightarrow$ нормальное управление $\Rightarrow$ дискретные наблюдения\\
$\Rightarrow$ траектория-поводырь $\Rightarrow$ регуляризованная минимизация\\
$\Rightarrow$ сходимость к нормальному управлению.
\end{minipage}
\end{center}
\zh{不适定反问题 $\Rightarrow$ 正规控制 $\Rightarrow$ 离散观测 $\Rightarrow$ 引导轨线 $\Rightarrow$ 正则化极小化 $\Rightarrow$ 收敛到正规控制。}

\ticket{1}{Обратные задачи, связанные с динамическими системами, описываемыми обыкновенными дифференциальными уравнениями. Метод динамической регуляризации для задачи восстановления характеристик систем обыкновенных дифференциальных уравнений.}{与常微分方程描述的动力系统有关的反问题。用于恢复常微分方程系统特性的动态正则化方法。}

\qpart{1}{Обратные задачи для систем ОДУ}{常微分方程系统的反问题}

Рассмотрим управляемую систему обыкновенных дифференциальных уравнений
\zh{考虑如下由常微分方程描述的受控系统：}
\begin{equation}
    \dot y(t)=f(t,y(t))u(t)+g(t,y(t)),\qquad y(0)=y_0,
    \qquad 0\le t\le T,
\end{equation}
где
\zh{其中}
\[
    y(t)\in \R^n,\qquad u(t)\in P\subset \R^r.
\]
Функции $f,g$, начальное состояние $y_0$, множество $P$ и горизонт $T$ считаются известными. Множество допустимых управлений имеет вид
\zh{函数 $f,g$、初始状态 $y_0$、集合 $P$ 以及时间区间长度 $T$ 均视为已知。可行控制集合写成}
\begin{equation}
    U=\{u(\cdot)\in L_2^r(0,T): u(t)\in P\ \text{для почти всех }t\in[0,T]\}.
\end{equation}
\zh{也就是说，$U$ 由几乎处处取值于 $P$ 的平方可积控制函数组成。}

Прямая задача состоит в нахождении траектории $y(\cdot)$ по заданному управлению $u(\cdot)$. Обратная задача состоит в восстановлении управления $u(\cdot)$ по наблюдаемой траектории $y(\cdot)$.
\zh{正问题是：给定控制 $u(\cdot)$，求相应的轨线 $y(\cdot)$。反问题是：根据观测到的轨线 $y(\cdot)$，恢复产生该轨线的控制 $u(\cdot)$。}

Обратная задача, вообще говоря, является некорректной. Во-первых, управление может быть неединственным. Например, если
\zh{一般来说，反问题是不适定的。首先，控制可能不唯一。例如，如果}
\[
    \dot y(t)=u_1(t)-u_2(t),\qquad y(0)=1,
\]
то наблюдаемой траектории $y(t)\equiv 1$ соответствует любое управление, для которого
\zh{那么对于观测轨线 $y(t)\equiv 1$，任何满足下式的控制都可以产生同一条轨线：}
\[
    u_1(t)=u_2(t)\quad \text{почти всюду}.
\]
Во-вторых, восстановление управления по зашумлённым наблюдениям может быть неустойчивым.
\zh{其次，当观测中含有噪声时，根据观测恢复控制的过程可能是不稳定的。}

Чтобы устранить неединственность, вводят понятие нормального управления. Пусть $y^r(\cdot)$ --- реальная наблюдаемая траектория, порождённая некоторым допустимым управлением. Определим множество
\zh{为了消除非唯一性，引入正规控制的概念。设 $y^r(\cdot)$ 是由某个可行控制产生的真实观测轨线。定义集合}
\begin{equation}
    U_* = \{u(\cdot)\in U: y(t,u)=y^r(t),\ 0\le t\le T\}.
\end{equation}
Нормальным управлением называется управление минимальной $L_2$-нормы:
\zh{正规控制就是在所有能够产生观测轨线的控制中 $L_2$ 范数最小的那个控制：}
\begin{equation}
    u_* = \operatorname*{arg\,min}_{u\in U_*}\norm{u}_{L_2^r(0,T)}.
\end{equation}
При стандартных условиях, например при липшицевости $f,g$, компактности $P$, а также выпуклости $P$, прямая задача корректна, а нормальное управление существует и единственно.
\zh{在标准条件下，例如 $f,g$ 满足 Lipschitz 条件，$P$ 为紧集且凸集时，正问题是适定的，并且正规控制存在且唯一。}

\qpart{2}{Метод динамической регуляризации для восстановления характеристик систем ОДУ}{恢复常微分方程系统特性的动态正则化方法}

На практике точная траектория неизвестна. Имеются только дискретные наблюдения
\zh{在实际问题中，精确轨线通常未知，只有离散时刻上的观测数据：}
\[
    0=t_0<t_1<\cdots<t_N=T,
\]
\begin{equation}
    |\tilde y_i-y^r(t_i)|\le \delta,
    \qquad i=0,\ldots,N.
\end{equation}
Обозначим шаг сетки
\zh{记网格步长为}
\begin{equation}
    h=\max_{0\le i\le N-1}(t_{i+1}-t_i).
\end{equation}
Требуется по данным $f,g,T,P,\{t_i\},\{\tilde y_i\},\delta$ построить управление $\tilde u_{\delta,h}(\cdot)$ такое, чтобы при согласованном стремлении $\delta,h\to0$ оно сходилось к нормальному управлению $u_*$. 
\zh{要求根据数据 $f,g,T,P,\{t_i\},\{\tilde y_i\},\delta$ 构造控制 $\tilde u_{\delta,h}(\cdot)$，使得当 $\delta,h\to0$ 且各参数协调趋于零时，该控制收敛到正规控制 $u_*$.}

Метод динамической регуляризации строит кусочно-постоянное управление. На каждом промежутке $[t_i,t_{i+1})$ полагают
\zh{动态正则化方法构造分段常值控制。在每个区间 $[t_i,t_{i+1})$ 上令}
\begin{equation}
    \tilde u_{\delta,h}(t)=u_i,
    \qquad u_i\in P.
\end{equation}
Вводится вспомогательная траектория, называемая поводырём:
\zh{引入一条辅助轨线，称为引导轨线：}
\begin{equation}
    z_0=\tilde y_0,
\end{equation}
\begin{equation}
    \frac{z_{i+1}-z_i}{t_{i+1}-t_i}
    =f(t_i,\tilde y_i)u_i+g(t_i,\tilde y_i),
    \qquad i=0,\ldots,N-1.
\end{equation}

Идея метода состоит в том, чтобы выбирать управление $u_i$ так, чтобы движение поводыря сближалось с наблюдаемой траекторией. Это реализует принцип экстремального сдвига Красовского. Для устойчивости к ошибкам наблюдений вводится стабилизатор Тихонова. На $i$-м шаге минимизируют функционал
\zh{该方法的思想是：选择控制 $u_i$，使引导轨线的运动逐步接近观测轨线。这实现了 Krasovskii 的极端位移原理。为了提高对观测误差的稳定性，引入 Tikhonov 稳定项。在第 $i$ 步，极小化泛函}
\begin{equation}
    T_i(v)=
    2\inner{z_i-\tilde y_i}{f(t_i,\tilde y_i)v+g(t_i,\tilde y_i)}
    +\alpha |v|^2,
    \qquad v\in P,
\end{equation}
где $\alpha>0$ --- параметр регуляризации. Управление $u_i\in P$ выбирается приближённо:
\zh{其中 $\alpha>0$ 是正则化参数。控制 $u_i\in P$ 按近似方式选取：}
\begin{equation}
    T_i(u_i)\le \inf_{v\in P}T_i(v)+\varepsilon,
\end{equation}
где $\varepsilon>0$ характеризует точность решения вспомогательной задачи минимизации.
\zh{其中 $\varepsilon>0$ 表示辅助极小化问题的求解精度。}

Смысл члена $\alpha |v|^2$ состоит в том, что он подавляет неустойчивость обратной задачи и одновременно направляет построение к нормальному управлению, то есть к управлению минимальной нормы.
\zh{项 $\alpha |v|^2$ 的意义在于：它抑制反问题的不稳定性，同时把构造过程引向正规控制，也就是引向范数最小的控制。}

Основная теорема сходимости имеет следующий смысл. Если
\zh{主要收敛定理的含义如下。如果}
\[
    \alpha\to0,
    \qquad \delta\to0,
    \qquad h\to0,
    \qquad \varepsilon\to0
\]
и параметры согласованы, например
\zh{并且这些参数满足协调条件，例如}
\begin{equation}
    \frac{\delta+h+\varepsilon}{\alpha}\to0,
\end{equation}
то построенная траектория поводыря сходится к точной наблюдаемой траектории:
\zh{那么构造出的引导轨线收敛到精确观测轨线：}
\begin{equation}
    \norm{z_h-y^r}_{C[0,T]}\to0,
\end{equation}
а построенные кусочно-постоянные управления сходятся к нормальному управлению:
\zh{同时，构造出的分段常值控制收敛到正规控制：}
\begin{equation}
    \norm{\tilde u_{\delta,h}-u_*}_{L_2^r(0,T)}\to0.
\end{equation}

Таким образом, метод динамической регуляризации заменяет неустойчивое прямое восстановление управления устойчивой пошаговой процедурой. На каждом шаге используется только уже поступившая информация, поэтому метод имеет динамический характер и может применяться при последовательной обработке наблюдений.
\zh{因此，动态正则化方法用稳定的逐步过程代替了不稳定的直接控制恢复。在每一步中只使用已经获得的信息，所以该方法具有动态性质，并且适用于对观测数据的逐步处理。}

\plainheading{Краткая версия для запоминания}{记忆版要点}

\begin{enumerate}
    \item Модель: $\dot y=f(t,y)u+g(t,y)$, $u(t)\in P$.
    \zh{模型为 $\dot y=f(t,y)u+g(t,y)$，其中 $u(t)\in P$。}
    \item Прямая задача: по $u(\cdot)$ найти $y(\cdot)$.
    \zh{正问题：由 $u(\cdot)$ 求 $y(\cdot)$。}
    \item Обратная задача: по наблюдаемой траектории восстановить $u(\cdot)$.
    \zh{反问题：由观测轨线恢复 $u(\cdot)$。}
    \item Обратная задача некорректна: возможны неединственность и неустойчивость.
    \zh{反问题是不适定的，可能出现非唯一性和不稳定性。}
    \item Вводится нормальное управление минимальной $L_2$-нормы.
    \zh{引入 $L_2$ 范数最小的正规控制。}
    \item По дискретным наблюдениям строится поводырь.
    \zh{根据离散观测构造引导轨线。}
    \item На каждом шаге решается регуляризованная задача минимизации с членом $\alpha |v|^2$.
    \zh{每一步求解带有 $\alpha |v|^2$ 项的正则化极小化问题。}
    \item При согласовании параметров построенное управление сходится к нормальному управлению.
    \zh{在参数协调趋于零时，构造出的控制收敛到正规控制。}
\end{enumerate}

\ticket{2}{Обратные задачи, связанные с динамическими системами, описываемыми уравнениями с частными производными. Метод динамической регуляризации для задач восстановления характеристик параболических уравнений.}{与偏微分方程描述的动力系统有关的反问题。用于恢复抛物型方程特性的动态正则化方法。}

\qpart{1}{Обратные задачи для систем с уравнениями в частных производных}{偏微分方程系统的反问题}

Рассмотрим типичную обратную задачу для параболического уравнения на примере уравнения теплопроводности:
\zh{以热传导方程为例，考虑抛物型方程中的一个典型反问题：}
\begin{equation}
\begin{cases}
    y_t(t,x)=a^2y_{xx}(t,x)+u(t,x), &(t,x)\in Q=(0,T)\times(0,l),\\
    y(t,0)=y(t,l)=0,\\
    y(0,x)=y_0(x).
\end{cases}
\end{equation}
Здесь $y(t,x)$ --- состояние системы, $u(t,x)$ --- неизвестная правая часть или распределённое управление.
\zh{这里 $y(t,x)$ 是系统状态，$u(t,x)$ 是未知右端项，或者说是分布式控制。}

Прямая задача состоит в том, чтобы по заданному $u(t,x)$ найти решение $y(t,x)$. Обратная задача состоит в восстановлении $u(t,x)$ по наблюдениям состояния $y(t,x)$.
\zh{正问题是：给定 $u(t,x)$，求解 $y(t,x)$。反问题是：根据对状态 $y(t,x)$ 的观测来恢复 $u(t,x)$。}

Такая обратная задача является некорректной: восстановление правой части по зашумлённым наблюдениям неустойчиво, а решение может быть неединственным. Поэтому, как и в случае ОДУ, вводится нормальное управление, то есть управление минимальной нормы среди всех управлений, порождающих наблюдаемую траекторию.
\zh{这样的反问题是不适定的：根据带噪观测恢复右端项是不稳定的，而且解可能不唯一。因此，与常微分方程的情形一样，引入正规控制，即在所有产生观测轨线的控制中范数最小的控制。}

Перейдём к абстрактной операторной форме. Пусть
\zh{转到抽象算子形式。令}
\[
    H=L_2(0,l),\qquad V=H_0^1(0,l).
\]
Имеем цепочку вложений
\zh{有如下嵌入链：}
\begin{equation}
    V\subset H\simeq H^*\subset V^*,
\end{equation}
где вложение $V\subset H$ непрерывно, плотно и компактно.
\zh{其中嵌入 $V\subset H$ 是连续、稠密且紧的。}

Определим оператор
\zh{定义算子}
\[
    A:V\to V^*
\]
по формуле
\zh{其定义式为}
\begin{equation}
    \langle Ay,v\rangle_{V^*,V}
    =a^2\int_0^l y_x(x)v_x(x)\dd x,
    \qquad y,v\in V.
\end{equation}
Оператор $A$ является симметричным и положительно определённым в энергетическом смысле.
\zh{算子 $A$ 在能量意义下是对称且正定的。}

Тогда уравнение теплопроводности записывается как
\zh{于是热传导方程可以写成}
\begin{equation}
    \dot y(t)+Ay(t)=u(t),\qquad y(0)=y_0.
\end{equation}
Равенство понимается в пространстве $V^*$. При
\zh{该等式在空间 $V^*$ 中理解。当}
\[
    u\in L_2(0,T;H),\qquad y_0\in H
\]
естественный класс слабых решений имеет вид
\zh{自然的弱解空间为}
\begin{equation}
    y\in L_2(0,T;V),\qquad \dot y\in L_2(0,T;V^*).
\end{equation}
Отсюда следует
\zh{由此可得}
\begin{equation}
    y\in C([0,T];H),
\end{equation}
поэтому значения $y(t_i)\in H$ корректно определены.
\zh{因此 $y(t_i)\in H$ 的值是良定义的。}

Пусть допустимые управления задаются множеством
\zh{设可行控制集合为}
\begin{equation}
    U=\{u(\cdot)\in L_2(0,T;H): u(t)\in P\ \text{для почти всех }t\in[0,T]\},
\end{equation}
где $P\subset H$ --- выпуклое, замкнутое и ограниченное множество.
\zh{其中 $P\subset H$ 是凸的、闭的且有界的集合。}

Пусть $y^r(t)$ --- точная наблюдаемая траектория. Множество управлений, порождающих эту траекторию, обозначим
\zh{设 $y^r(t)$ 是精确观测轨线。把产生这条轨线的控制集合记为}
\begin{equation}
    U_* = \{u\in U:y(t,u)=y^r(t),\ 0\le t\le T\}.
\end{equation}
Нормальное управление определяется как
\zh{正规控制定义为}
\begin{equation}
    u_* = \operatorname*{arg\,min}_{u\in U_*}\norm{u}_{L_2(0,T;H)}.
\end{equation}
\zh{也就是在集合 $U_*$ 中 $L_2(0,T;H)$ 范数最小的控制。}

На практике точная траектория неизвестна. В дискретные моменты
\zh{在实践中，精确轨线未知。在离散时刻}
\[
    0=t_0<t_1<\cdots<t_N=T
\]
имеются наблюдения
\zh{给定观测数据}
\begin{equation}
    \norm{y^r(t_i)-\tilde y_i}_H\le \delta,
    \qquad i=0,\ldots,N.
\end{equation}
Положим
\zh{记}
\begin{equation}
    h=\max_{0\le i\le N-1}(t_{i+1}-t_i).
\end{equation}

\qpart{2}{Метод динамической регуляризации для восстановления характеристик параболических уравнений}{恢复抛物型方程特性的动态正则化方法}

Метод динамической регуляризации строит кусочно-постоянное управление
\zh{动态正则化方法构造分段常值控制}
\begin{equation}
    \tilde u_{\delta,h}(t)=u_i,
    \qquad t\in[t_i,t_{i+1}),
    \qquad u_i\in P.
\end{equation}

Вводится траектория-поводырь $z_h(t)$, которая на каждом промежутке удовлетворяет уравнению
\zh{引入引导轨线 $z_h(t)$，它在每个区间上满足方程}
\begin{equation}
    \dot z_h(t)+Az_h(t)=u_i,
    \qquad t\in(t_i,t_{i+1}),
\end{equation}
с начальным условием, полученным на предыдущем шаге.
\zh{其初始条件由前一步得到。}

На $i$-м шаге управление выбирается из регуляризованной задачи минимизации:
\zh{在第 $i$ 步，控制由如下正则化极小化问题确定：}
\begin{equation}
    \varphi_i(v)
    =2\inner{z_h(t_i)-\tilde y_i}{v}_H
    +\alpha\norm{v}_H^2
    \to \min_{v\in P},
\end{equation}
где $\alpha>0$ --- параметр регуляризации. Управление $u_i\in P$ выбирается с точностью $\varepsilon>0$:
\zh{其中 $\alpha>0$ 是正则化参数。控制 $u_i\in P$ 以精度 $\varepsilon>0$ 选取：}
\begin{equation}
    \varphi_i(u_i)\le \inf_{v\in P}\varphi_i(v)+\varepsilon.
\end{equation}

Первое слагаемое реализует идею экстремального сдвига: поводырь должен сближаться с наблюдаемой траекторией. Второе слагаемое $\alpha\norm{v}_H^2$ является стабилизатором Тихонова и обеспечивает устойчивость восстановления.
\zh{第一项体现极端位移思想：引导轨线应当接近观测轨线。第二项 $\alpha\norm{v}_H^2$ 是 Tikhonov 稳定项，用来保证恢复过程的稳定性。}

Смысл теоремы сходимости следующий. Если
\zh{收敛定理的含义如下。如果}
\[
    \delta\to0,
    \qquad h\to0,
    \qquad \varepsilon\to0,
    \qquad \alpha\to0
\]
и параметры согласованы, то
\zh{并且各参数满足协调条件，那么}
\begin{equation}
    z_h\to y^r \quad\text{в } C([0,T];H),
\end{equation}
а построенные управления сходятся к нормальному управлению:
\zh{构造出的控制收敛到正规控制：}
\begin{equation}
    \tilde u_{\delta,h}\to u_* \quad\text{в } L_2(0,T;H).
\end{equation}
В студенческих материалах используется условие вида
\zh{学生材料中使用如下形式的条件：}
\begin{equation}
    \frac{\sqrt h+\delta+\varepsilon}{\alpha}\to0,
\end{equation}
но точная форма оценки требует сверки с оригинальным текстом Осипова--Васильева--Потапова и поэтому помечается как требующая проверки.
\zh{不过估计的精确形式需要与 Osipov--Vasiliev--Potapov 的原文核对，因此这里标记为需要进一步检查。}

Итак, метод динамической регуляризации для параболического уравнения заменяет неустойчивое восстановление правой части $u(t,x)$ устойчивой пошаговой процедурой, основанной на траектории-поводыре и регуляризованной минимизации.
\zh{因此，对于抛物型方程，动态正则化方法把不稳定的右端项 $u(t,x)$ 恢复问题，替换成基于引导轨线和正则化极小化的稳定逐步过程。}

\plainheading{Краткая версия для запоминания}{记忆版要点}

\begin{enumerate}
    \item Модель: $y_t=a^2y_{xx}+u$ с нулевыми граничными условиями.
    \zh{模型为 $y_t=a^2y_{xx}+u$，边界条件为零。}
    \item Обратная задача: восстановить $u(t,x)$ по наблюдениям $y(t,x)$.
    \zh{反问题：根据 $y(t,x)$ 的观测恢复 $u(t,x)$。}
    \item Используется тройка $V\subset H\simeq H^*\subset V^*$, где $V=H_0^1(0,l)$, $H=L_2(0,l)$.
    \zh{使用三元组 $V\subset H\simeq H^*\subset V^*$，其中 $V=H_0^1(0,l)$，$H=L_2(0,l)$。}
    \item Оператор $A:V\to V^*$ задаётся формой $\langle Ay,v\rangle=a^2\int_0^l y_xv_x\dd x$.
    \zh{算子 $A:V\to V^*$ 由形式 $\langle Ay,v\rangle=a^2\int_0^l y_xv_x\dd x$ 给出。}
    \item Уравнение: $\dot y+Ay=u$.
    \zh{方程为 $\dot y+Ay=u$。}
    \item Строится поводырь $\dot z_h+Az_h=u_i$.
    \zh{构造引导轨线 $\dot z_h+Az_h=u_i$。}
    \item Управление $u_i\in P$ выбирается из регуляризованной задачи $2\langle z_h(t_i)-\tilde y_i,v\rangle_H+\alpha\norm{v}_H^2\to\min$.
    \zh{控制 $u_i\in P$ 从正则化问题 $2\langle z_h(t_i)-\tilde y_i,v\rangle_H+\alpha\norm{v}_H^2\to\min$ 中选取。}
    \item При согласовании параметров $z_h\to y^r$, $\tilde u\to u_*$.
    \zh{在参数协调趋于零时，$z_h\to y^r$，$\tilde u\to u_*$。}
\end{enumerate}

\ticket{3}{Обратные задачи, связанные с динамическими системами, описываемыми уравнениями с частными производными. Метод динамической регуляризации для задач восстановления характеристик гиперболических уравнений.}{与偏微分方程描述的动力系统有关的反问题。用于恢复双曲型方程特性的动态正则化方法。}

\qpart{1}{Обратные задачи для систем с уравнениями в частных производных}{偏微分方程系统的反问题}

Рассмотрим типичную обратную задачу для гиперболического уравнения на примере волнового уравнения:
\zh{以波动方程为例，考虑双曲型方程中的一个典型反问题：}
\begin{equation}
\begin{cases}
    y_{tt}(t,x)=a^2y_{xx}(t,x)+u(t,x), &(t,x)\in Q=(0,T)\times(0,l),\\
    y(t,0)=y(t,l)=0,\\
    y(0,x)=y_0(x),\\
    y_t(0,x)=y_1(x).
\end{cases}
\end{equation}
Здесь $u(t,x)$ --- неизвестная правая часть или распределённое управление.
\zh{这里 $u(t,x)$ 是未知右端项，或者说是分布式控制。}

Прямая задача состоит в том, чтобы по заданному $u(t,x)$ найти решение $y(t,x)$. Обратная задача состоит в восстановлении $u(t,x)$ по наблюдениям состояния $y$ или скорости $y_t$.
\zh{正问题是：给定 $u(t,x)$，求解 $y(t,x)$。反问题是：根据对状态 $y$ 或速度 $y_t$ 的观测来恢复 $u(t,x)$。}

В отличие от параболического случая, уравнение является вторым порядком по времени, поэтому естественное состояние системы есть пара
\zh{与抛物型情形不同，该方程关于时间是二阶的，因此系统的自然状态是一对量：}
\begin{equation}
    X(t)=(y(t),\dot y(t)).
\end{equation}

Введём пространства
\zh{引入空间}
\[
    H=L_2(0,l),\qquad V=H_0^1(0,l),
\]
и цепочку вложений
\zh{以及嵌入链}
\begin{equation}
    V\subset H\simeq H^*\subset V^*.
\end{equation}
Оператор
\zh{算子}
\[
    A:V\to V^*
\]
определяется по формуле
\zh{定义为}
\begin{equation}
    \langle Ay,v\rangle_{V^*,V}
    =a^2\int_0^l y_x(x)v_x(x)\dd x,
    \qquad y,v\in V.
\end{equation}
Тогда волновое уравнение записывается в абстрактной форме
\zh{于是波动方程可以写成抽象形式}
\begin{equation}
    \ddot y(t)+Ay(t)=u(t),
\end{equation}
\begin{equation}
    y(0)=y_0,
    \qquad \dot y(0)=y_1.
\end{equation}
Равенство понимается в пространстве $V^*$.
\zh{该等式在空间 $V^*$ 中理解。}

При
\zh{当}
\[
    y_0\in V,
    \qquad y_1\in H,
    \qquad u\in L_2(0,T;H)
\]
естественный энергетический класс решений имеет вид
\zh{自然的能量解空间为}
\begin{equation}
    y\in C([0,T];V),
    \qquad \dot y\in C([0,T];H),
\end{equation}
а уравнение выполняется в $V^*$.
\zh{并且方程在 $V^*$ 中成立。}

Множество допустимых управлений зададим как
\zh{可行控制集合定义为}
\begin{equation}
    U=\{u(\cdot)\in L_2(0,T;H): u(t)\in P\ \text{для почти всех }t\in[0,T]\},
\end{equation}
где $P\subset H$ --- выпуклое, замкнутое и ограниченное множество.
\zh{其中 $P\subset H$ 是凸的、闭的且有界的集合。}

Если $y^r$ --- точное наблюдаемое движение, то множество управлений, порождающих его, обозначается
\zh{如果 $y^r$ 是精确观测到的运动，则产生该运动的控制集合记为}
\begin{equation}
    U_* = \{u\in U:y(t,u)=y^r(t),\ 0\le t\le T\}.
\end{equation}
Нормальное управление определяется как управление минимальной нормы:
\zh{正规控制定义为范数最小的控制：}
\begin{equation}
    u_* = \operatorname*{arg\,min}_{u\in U_*}\norm{u}_{L_2(0,T;H)}.
\end{equation}

На практике точное движение неизвестно. Например, если наблюдается скорость, то в моменты
\zh{在实践中，精确运动未知。例如，如果观测的是速度，则在时刻}
\[
    0=t_0<t_1<\cdots<t_N=T
\]
имеются данные
\zh{给定数据满足}
\begin{equation}
    \norm{\dot y^r(t_i)-\tilde w_i}_H\le \delta.
\end{equation}
Если наблюдается само состояние, то вместо этого пишут
\zh{如果观测的是状态本身，则改写为}
\begin{equation}
    \norm{y^r(t_i)-\tilde y_i}_H\le \delta.
\end{equation}
На экзамене важно указать, какой именно тип наблюдений используется.
\zh{在考试回答中，需要明确说明所使用的是哪一种观测方式。}

\qpart{2}{Метод динамической регуляризации для восстановления характеристик гиперболических уравнений}{恢复双曲型方程特性的动态正则化方法}

Метод динамической регуляризации строит кусочно-постоянное управление
\zh{动态正则化方法构造分段常值控制}
\begin{equation}
    \tilde u_{\delta,h}(t)=u_i,
    \qquad t\in[t_i,t_{i+1}),
    \qquad u_i\in P,
\end{equation}
где
\zh{其中}
\begin{equation}
    h=\max_i(t_{i+1}-t_i).
\end{equation}

Вводится траектория-поводырь $z_h(t)$, удовлетворяющая на каждом промежутке уравнению
\zh{引入引导轨线 $z_h(t)$，它在每个区间上满足方程}
\begin{equation}
    \ddot z_h(t)+Az_h(t)=u_i,
    \qquad t\in(t_i,t_{i+1}).
\end{equation}
Начальные данные $z_h(t_i)$, $\dot z_h(t_i)$ берутся из предыдущего шага.
\zh{初始数据 $z_h(t_i)$ 与 $\dot z_h(t_i)$ 由前一步给出。}

Если наблюдается скорость, то естественная регуляризованная задача выбора управления имеет вид
\zh{如果观测的是速度，则自然的正则化控制选择问题为}
\begin{equation}
    \varphi_i(v)
    =2\inner{\dot z_h(t_i)-\tilde w_i}{v}_H
    +\alpha\norm{v}_H^2
    \to \min_{v\in P},
\end{equation}
где $\alpha>0$ --- параметр регуляризации. Управление $u_i\in P$ выбирается с точностью $\varepsilon>0$:
\zh{其中 $\alpha>0$ 是正则化参数。控制 $u_i\in P$ 以精度 $\varepsilon>0$ 选取：}
\begin{equation}
    \varphi_i(u_i)\le \inf_{v\in P}\varphi_i(v)+\varepsilon.
\end{equation}

Если наблюдается состояние $y(t_i)$, то записывается эквивалентная форма алгоритма, согласованная с выбранным способом наблюдения; в любом случае управление выбирается пошагово из задачи экстремального сдвига с добавлением стабилизатора Тихонова $\alpha\norm{v}_H^2$.
\zh{如果观测的是状态 $y(t_i)$，则写出与该观测方式相一致的等价算法形式；无论哪种情形，控制都是逐步从极端位移问题中选取，并加入 Tikhonov 稳定项 $\alpha\norm{v}_H^2$。}

Смысл этого стабилизатора тот же, что и в случае ОДУ и параболического уравнения: он подавляет неустойчивость обратной задачи и направляет решение к нормальному управлению минимальной нормы.
\zh{该稳定项的意义与常微分方程和抛物型方程中的情形相同：它抑制反问题的不稳定性，并把解引向范数最小的正规控制。}

Сходимость формулируется в энергетических нормах. При согласованном стремлении
\zh{收敛性在能量范数中表述。当参数协调趋于零，即}
\[
    \delta\to0,
    \qquad h\to0,
    \qquad \varepsilon\to0,
    \qquad \alpha\to0
\]
получают
\zh{可以得到}
\begin{equation}
    (z_h,\dot z_h)\to (y^r,\dot y^r)
\end{equation}
в соответствующем пространстве траекторий, например в энергетическом смысле, и
\zh{该收敛在相应的轨线空间中成立，例如在能量意义下成立，并且}
\begin{equation}
    \tilde u_{\delta,h}\to u_* \quad\text{в } L_2(0,T;H).
\end{equation}
Точная форма оценки и условия согласования параметров требуют сверки с оригинальным текстом Осипова--Васильева--Потапова или с работами Осипова--Кряжимского--Максимова по системам с распределёнными параметрами.
\zh{估计的精确形式以及参数协调条件需要与 Osipov--Vasiliev--Potapov 的原文，或者与 Osipov--Kryazhimskii--Maksimov 关于分布参数系统的工作进行核对。}

Итак, для гиперболического уравнения метод динамической регуляризации также заменяет некорректное восстановление неизвестной правой части устойчивой динамической процедурой. Главное отличие от параболического случая состоит в том, что состояние включает как $y$, так и $\dot y$, а сходимость должна формулироваться в энергетическом пространстве.
\zh{因此，对于双曲型方程，动态正则化方法同样把不适定的未知右端项恢复问题替换为稳定的动态过程。与抛物型情形的主要区别在于，状态同时包含 $y$ 和 $\dot y$，并且收敛性必须在能量空间中表述。}

\plainheading{Краткая версия для запоминания}{记忆版要点}
\begin{enumerate}
    \item Модель: $y_{tt}=a^2y_{xx}+u$ с нулевыми граничными условиями.
    \zh{模型为 $y_{tt}=a^2y_{xx}+u$，边界条件为零。}
    \item Обратная задача: восстановить $u(t,x)$ по наблюдениям состояния или скорости.
    \zh{反问题：根据状态或速度的观测恢复 $u(t,x)$。}
    \item Поскольку уравнение второго порядка по времени, состояние есть пара $(y,\dot y)$.
    \zh{由于方程关于时间为二阶，状态是一对量 $(y,\dot y)$。}
    \item Энергетическое пространство: $y\in C([0,T];H_0^1)$, $\dot y\in C([0,T];L_2)$.
    \zh{能量空间为 $y\in C([0,T];H_0^1)$，$\dot y\in C([0,T];L_2)$。}
    \item Абстрактная форма: $\ddot y+Ay=u$, где $A:V\to V^*$.
    \zh{抽象形式为 $\ddot y+Ay=u$，其中 $A:V\to V^*$。}
    \item Строится поводырь $\ddot z_h+Az_h=u_i$.
    \zh{构造引导轨线 $\ddot z_h+Az_h=u_i$。}
    \item Управление выбирается из задачи экстремального сдвига с регуляризатором $\alpha\norm{v}_H^2$.
    \zh{控制从带有正则项 $\alpha\norm{v}_H^2$ 的极端位移问题中选取。}
    \item Сходимость формулируется в энергетических нормах, а управление сходится к нормальному управлению.
    \zh{收敛性在能量范数中表述，控制收敛到正规控制。}
\end{enumerate}

\blocktitle{B}{Экономические модели и задачи оптимального управления}{Билеты 4--8}
\zh{B板块：经济模型与最优控制问题。第4--8题签。}

\ticket{4}{Одномерная модель Рамсея на бесконечном горизонте планирования с производственной функцией в форме Кобба--Дугласа. Задача о наилучших пропорциях производства и потребления. Особый режим. Качественный характер решения.}{无限规划期上的一维 Ramsey 模型，生产函数为 Cobb--Douglas 形式。生产与消费最佳比例问题。奇异机制。解的定性结构。}

\qpart{1}{Одномерная модель Рамсея на бесконечном горизонте с функцией Кобба--Дугласа}{无限规划期上的一维 Ramsey 模型与 Cobb--Douglas 函数}

\plainheading{Смысл задачи}{问题含义}

Рассматривается замкнутая экономика, в которой произведенный продукт распределяется между потреблением и инвестициями. Управление задает долю выпуска, направляемую на накопление капитала. Нужно найти оптимальную пропорцию между текущим потреблением и инвестициями на бесконечном горизонте планирования.
\zh{考虑一个封闭经济体，其中产出在消费和投资之间分配。控制变量表示产出中用于资本积累的份额。问题是在无限规划期上寻找当前消费与投资之间的最优比例。}

\qpart{2}{Задача о наилучших пропорциях производства и потребления}{生产与消费最佳比例问题}

\plainheading{Постановка модели}{模型设定}

Пусть
\zh{设}
\[
    x(t)>0
\]
-- удельный капитал, а
\zh{为单位资本，并设}
\[
    u(t)\in[0,1]
\]
-- доля выпуска, направляемая в инвестиции. Тогда \(1-u(t)\) -- доля выпуска, направляемая на потребление.
\zh{为用于投资的产出份额。那么 \(1-u(t)\) 就是用于消费的产出份额。}

Производственная функция Кобба--Дугласа в интенсивной форме:
\zh{Cobb--Douglas 生产函数的集约形式为}
\[
    F(x)=Ax^{\eps},\qquad A>0,\quad 0<\eps<1.
\]
Динамика капитала:
\zh{资本动态为}
\[
    \dot x(t)=u(t)F(x(t))-\mu x(t),\qquad x(0)=x_0>0,
\]
где \(\mu>0\) -- коэффициент выбытия капитала. Если модель получена из переменных совокупного капитала \(K\) и труда \(L\), то \(\mu\) может включать как амортизацию, так и рост трудовых ресурсов.
\zh{其中 \(\mu>0\) 是资本退出率。如果模型是由总资本 \(K\) 和劳动 \(L\) 的变量化简得到的，那么 \(\mu\) 可以同时包含折旧率和劳动资源增长率。}

Функционал:
\zh{目标泛函为}
\[
    J(u)=\int_0^\infty e^{-\nu t}(1-u(t))F(x(t))\,dt\to\max,
\]
где \(\nu>0\) -- коэффициент дисконтирования.
\zh{其中 \(\nu>0\) 是贴现系数。}

\qpart{3}{Особый режим}{奇异机制}

\plainheading{Принцип максимума и функция переключения}{极大值原理与切换函数}

Гамильтониан в приведенных к начальному моменту переменных:
\zh{用折现到初始时刻的变量写出的 Hamilton 函数为}
\[
    H=e^{-\nu t}(1-u)F(x)+\psi\bigl(uF(x)-\mu x\bigr).
\]
Выделим часть, зависящую от \(u\):
\zh{取出其中依赖于 \(u\) 的部分：}
\[
    H=e^{-\nu t}F(x)-\psi\mu x+uF(x)\bigl(\psi-e^{-\nu t}\bigr).
\]
Так как \(F(x)>0\), функция переключения имеет вид
\zh{由于 \(F(x)>0\)，切换函数为}
\[
    \Phi(t)=\psi(t)-e^{-\nu t}.
\]
Следовательно,
\zh{因此}
\[
    u^*(t)=
    \begin{cases}
        1, & \Phi(t)>0,\\
        0, & \Phi(t)<0,\\
        \text{не определяется из условия максимума}, & \Phi(t)=0.
    \end{cases}
\]
\zh{也就是说，当 \(\Phi(t)>0\) 时选择全部投资，\(\Phi(t)<0\) 时选择全部消费，而 \(\Phi(t)=0\) 时极大值条件本身不能唯一确定控制。}

\qpart{4}{Качественный характер решения}{解的定性结构}

\plainheading{Текущая сопряженная переменная и особый режим}{当前值伴随变量与奇异机制}

Введем текущую сопряженную переменную
\zh{引入当前值伴随变量}
\[
    p(t)=e^{\nu t}\psi(t).
\]
Тогда текущий гамильтониан:
\zh{于是当前值 Hamilton 函数为}
\[
    M(x,u,p)=(1-u)F(x)+p\bigl(uF(x)-\mu x\bigr).
\]
Условие максимума зависит от знака \(p-1\):
\zh{极大值条件取决于 \(p-1\) 的符号：}
\[
    u^*(t)=
    \begin{cases}
        1, & p(t)>1,\\
        0, & p(t)<1,\\
        \text{особый режим}, & p(t)=1.
    \end{cases}
\]

На особом режиме
\zh{在奇异机制上}
\[
    p(t)\equiv1.
\]
Уравнение для текущей сопряженной переменной имеет вид
\zh{当前值伴随变量的方程为}
\[
    \dot p=\nu p-M_x(x,u,p).
\]
Так как
\zh{由于}
\[
    M_x=(1-u)F'(x)+p\bigl(uF'(x)-\mu\bigr),
\]
то при \(p=1\) получаем
\zh{所以当 \(p=1\) 时得到}
\[
    M_x=F'(x)-\mu.
\]
Из \(\dot p=0\) следует
\zh{由 \(\dot p=0\) 得}
\[
    0=\nu+\mu-F'(x).
\]
Поэтому особое состояние \(x_s\) определяется уравнением
\zh{因此奇异状态 \(x_s\) 由方程}
\[
    F'(x_s)=\mu+\nu
\]
задаётся. Для \(F(x)=Ax^\eps\):
\zh{确定。对于 \(F(x)=Ax^\eps\)，有}
\[
    A\eps x_s^{\eps-1}=\mu+\nu,
\]
откуда
\zh{因此}
\[
    x_s=\left(\frac{A\eps}{\mu+\nu}\right)^{\frac{1}{1-\eps}}.
\]
Особое управление определяется условием стационарности капитала:
\zh{奇异控制由资本保持平稳的条件确定：}
\[
    \dot x=0,
\]
то есть
\zh{也就是说}
\[
    u_sF(x_s)=\mu x_s.
\]
Следовательно,
\zh{因此}
\[
    u_s=\frac{\mu x_s}{F(x_s)}.
\]
Для функции Кобба--Дугласа это дает
\zh{对于 Cobb--Douglas 函数，这给出}
\[
    u_s=\frac{\mu\eps}{\mu+\nu}.
\]

\plainheading{Частный случай \texorpdfstring{\(F(x)=\sqrt{x}\)}{F(x)=sqrt(x)}}{特例 \texorpdfstring{\(F(x)=\sqrt{x}\)}{F(x)=sqrt(x)}}

В часто используемом частном случае
\zh{在常用的特例中}
\[
    A=1,
    \qquad
    \eps=\frac12,
\]
имеем
\zh{有}
\[
    F(x)=\sqrt{x}.
\]
Тогда
\zh{因此}
\[
    F'(x)=\frac{1}{2\sqrt{x}}.
\]
Из условия \(F'(x_s)=\mu+\nu\):
\zh{由条件 \(F'(x_s)=\mu+\nu\) 得}
\[
    \frac{1}{2\sqrt{x_s}}=\mu+\nu.
\]
Значит,
\zh{所以}
\[
    x_s=\frac{1}{4(\mu+\nu)^2}.
\]
Особое управление:
\zh{奇异控制为}
\[
    u_s=\frac{\mu x_s}{\sqrt{x_s}}=\mu\sqrt{x_s}=\frac{\mu}{2(\mu+\nu)}.
\]
На особом режиме
\zh{在奇异机制上}
\[
    x(t)\equiv x_s,
    \qquad
    u(t)\equiv u_s.
\]

\plainheading{Качественный характер решения}{解的定性结构}

Если начальный капитал меньше особого уровня,
\zh{如果初始资本低于奇异水平，即}
\[
    x_0<x_s,
\]
то сначала выбирается максимальное инвестирование
\zh{则首先选择最大投资}
\[
    u=1,
\]
чтобы увеличить капитал до \(x_s\), после чего траектория идет по особому режиму.
\zh{以便把资本提高到 \(x_s\)，之后轨线沿奇异机制运动。}

Если
\zh{如果}
\[
    x_0>x_s,
\]
то сначала выбирается
\zh{则首先选择}
\[
    u=0,
\]
то есть весь выпуск направляется на потребление, капитал уменьшается до \(x_s\), затем начинается особый режим.
\zh{即全部产出用于消费，使资本下降到 \(x_s\)，然后进入奇异机制。}

Если
\zh{如果}
\[
    x_0=x_s,
\]
то особый режим реализуется сразу:
\zh{则从一开始就实现奇异机制：}
\[
    x(t)\equiv x_s,
    \qquad
    u(t)\equiv u_s.
\]

Экономический смысл условия
\zh{条件}
\[
    F'(x_s)=\mu+\nu
\]
состоит в том, что на оптимальной магистрали предельная производительность капитала равна сумме дисконтирования и выбытия капитала. Это баланс между текущим потреблением и накоплением.
\zh{的经济含义是：在最优主干路径上，资本边际产出等于贴现率与资本退出率之和。这表示当前消费与资本积累之间的平衡。}

\plainheading{Краткая версия для запоминания}{记忆版要点}
\[
    \dot x=uF(x)-\mu x,
    \qquad
    J=\int_0^\infty e^{-\nu t}(1-u)F(x)\,dt\to\max.
\]
\[
    M=(1-u)F(x)+p(uF(x)-\mu x).
\]
Функция переключения в текущих переменных: \(p-1\).
\zh{用当前值变量表示时，切换函数为 \(p-1\)。}
\[
    p>1\Rightarrow u=1,
    \qquad
    p<1\Rightarrow u=0.
\]
Особый режим:
\zh{奇异机制为：}
\[
    p=1,
    \qquad
    F'(x_s)=\mu+\nu,
    \qquad
    u_sF(x_s)=\mu x_s.
\]
Для \(F(x)=\sqrt{x}\):
\zh{对于 \(F(x)=\sqrt{x}\)，有：}
\[
    x_s=\frac{1}{4(\mu+\nu)^2},
    \qquad
    u_s=\frac{\mu}{2(\mu+\nu)}.
\]

\ticket{5}{Двумерная задача распределения ресурсов в двухсекторной экономической модели с производственной функцией Кобба--Дугласа на конечном отрезке времени с максимизацией второй фазовой координаты в конечный момент времени. Особый режим. Качественный характер решения.}{有限时间区间上两部门 Cobb--Douglas 经济模型中的二维资源分配问题，目标是在终端时刻最大化第二个状态坐标。奇异机制。解的定性结构。}

\qpart{1}{Двумерная задача распределения ресурсов в двухсекторной модели Кобба--Дугласа}{两部门 Cobb--Douglas 模型中的二维资源分配问题}

\plainheading{Постановка задачи}{问题设定}

Рассмотрим задачу оптимального распределения ресурсов в двухсекторной экономической модели:
\zh{考虑两部门经济模型中的最优资源分配问题：}
\[
    \dot x_1(t)=u(t)F(x(t)),
    \qquad
    x_1(0)=x_{10}>0,
\]
\[
    \dot x_2(t)=(1-u(t))F(x(t)),
    \qquad
    x_2(0)=x_{20}>0,
\]
\[
    0\le u(t)\le 1,
    \qquad
    0\le t\le T.
\]
Здесь
\zh{这里}
\[
    x(t)=(x_1(t),x_2(t))\in\R^2_+
\]
характеризует уровни развития двух секторов экономики. Управление \(u(t)\) есть доля выпуска, направляемая в первый сектор, а \(1-u(t)\) -- доля выпуска, направляемая во второй сектор.
\zh{表示两个经济部门的发展水平。控制 \(u(t)\) 表示产出中投入第一部门的份额，而 \(1-u(t)\) 表示投入第二部门的份额。}

Целевой функционал:
\zh{目标泛函为}
\[
    J[u]=x_2(T)\to\max.
\]
Требуется максимизировать уровень развития второго сектора в конечный момент времени \(T\).
\zh{要求在终端时刻 \(T\) 最大化第二部门的发展水平。}

Производственная функция Кобба--Дугласа:
\zh{Cobb--Douglas 生产函数为}
\[
    F(x)=Ax_1^{\eps_1}x_2^{\eps_2},
\]
где
\zh{其中}
\[
    A>0,
    \qquad
    \eps_1>0,
    \qquad
    \eps_2>0,
    \qquad
    \eps_1+\eps_2=1.
\]
Постоянный множитель \(A>0\) можно устранить заменой масштаба времени, поэтому далее можно считать
\zh{常数因子 \(A>0\) 可以通过改变时间尺度消去，因此下面可以认为}
\[
    A=1,
    \qquad
    F(x)=x_1^{\eps_1}x_2^{\eps_2}.
\]

Так как \(\eps_1+\eps_2=1\), функция \(F\) положительно однородна степени 1:
\zh{因为 \(\eps_1+\eps_2=1\)，函数 \(F\) 是一次正齐次函数：}
\[
    F(\lambda x)=\lambda F(x),
    \qquad
    \lambda>0.
\]
Кроме того, функция Кобба--Дугласа вогнута на \(\R^2_+\). Это свойство используется при доказательстве оптимальности экстремального процесса.
\zh{此外，Cobb--Douglas 函数在 \(\R^2_+\) 上是凹函数。这个性质用于证明极值过程的最优性。}

\qpart{2}{Особый режим}{奇异机制}

\plainheading{Принцип максимума Понтрягина}{Pontryagin 极大值原理}

Так как функционал имеет тип Майера,
\zh{由于目标泛函属于 Mayer 型，}
\[
    \Phi(x)=x_2,
    \qquad
    \Phi'(x)=
    \begin{pmatrix}
        0\\1
    \end{pmatrix},
\]
терминальное условие для сопряженных переменных:
\zh{伴随变量的终端条件为}
\[
    \psi_1(T)=0,
    \qquad
    \psi_2(T)=1.
\]

Гамильтониан:
\zh{Hamilton 函数为}
\[
    H(x,\psi,u)=\psi_1uF(x)+\psi_2(1-u)F(x).
\]
Преобразуем:
\zh{将其改写为}
\[
    H(x,\psi,u)=F(x)\bigl(u(\psi_1-\psi_2)+\psi_2\bigr).
\]
Введем функцию переключения
\zh{引入切换函数}
\[
    \Pi(\psi)=\psi_1-\psi_2.
\]
Тогда
\zh{于是}
\[
    H(x,\psi,u)=F(x)\bigl(u\Pi(\psi)+\psi_2\bigr).
\]
Поскольку \(F(x)>0\) в положительном квадранте, условие максимума по \(u\in[0,1]\) дает
\zh{由于在正象限内 \(F(x)>0\)，对 \(u\in[0,1]\) 的极大值条件给出}
\[
    u^*(t)=
    \begin{cases}
        1, & \Pi(t)>0,\\
        0, & \Pi(t)<0,\\
        \text{не определяется однозначно}, & \Pi(t)=0.
    \end{cases}
\]
Если \(\Pi(t)\equiv0\) на некотором интервале, то возникает особый режим.
\zh{如果在某个区间上 \(\Pi(t)\equiv0\)，则出现奇异机制。}

Сопряженная система:
\zh{伴随系统为}
\[
    \dot\psi_1(t)=-F_{x_1}'(x(t))\bigl(u(t)\Pi(t)+\psi_2(t)\bigr),
\]
\[
    \dot\psi_2(t)=-F_{x_2}'(x(t))\bigl(u(t)\Pi(t)+\psi_2(t)\bigr).
\]

\qpart{3}{Качественный характер решения}{解的定性结构}

\plainheading{Особый режим}{奇异机制}

Пусть на некотором интервале выполнено
\zh{设在某个区间上满足}
\[
    \Pi(t)\equiv0.
\]
Тогда
\zh{则}
\[
    \psi_1(t)=\psi_2(t),
    \qquad
    \dot\psi_1(t)=\dot\psi_2(t).
\]
Из сопряженной системы получаем
\zh{由伴随系统得到}
\[
    F_{x_1}'(x(t))=F_{x_2}'(x(t)).
\]
Для функции Кобба--Дугласа
\zh{对于 Cobb--Douglas 函数，}
\[
    F_{x_1}'(x)=\frac{\eps_1}{x_1}F(x),
    \qquad
    F_{x_2}'(x)=\frac{\eps_2}{x_2}F(x).
\]
Следовательно,
\zh{因此}
\[
    \frac{\eps_1}{x_1}F(x)=\frac{\eps_2}{x_2}F(x).
\]
Так как \(F(x)>0\), получаем
\zh{由于 \(F(x)>0\)，得到}
\[
    \frac{x_1}{\eps_1}=\frac{x_2}{\eps_2}.
\]
Значит, особая магистраль в исходных координатах:
\zh{因此，原坐标下的奇异主干线为}
\[
    L_{\sng}=\left\{x\in\R^2_+:
    \frac{x_1}{\eps_1}=\frac{x_2}{\eps_2}\right\}.
\]

Чтобы траектория оставалась на этой магистрали, необходимо
\zh{为了使轨线保持在这条主干线上，需要}
\[
    \frac{\dot x_1}{\eps_1}=\frac{\dot x_2}{\eps_2}.
\]
Подставляя уравнения движения,
\zh{代入运动方程，得到}
\[
    \frac{u}{\eps_1}=\frac{1-u}{\eps_2}.
\]
Отсюда
\zh{由此可得}
\[
    \eps_2u=\eps_1(1-u),
\]
\[
    u(\eps_1+\eps_2)=\eps_1.
\]
Так как \(\eps_1+\eps_2=1\), получаем особое управление
\zh{由于 \(\eps_1+\eps_2=1\)，得到奇异控制}
\[
    u_{\sng}=\eps_1.
\]
Итак, особый режим задается условиями
\zh{所以，奇异机制由条件}
\[
    \frac{x_1}{\eps_1}=\frac{x_2}{\eps_2},
    \qquad
    u(t)=\eps_1
\]
задается.
\zh{给出。}

\plainheading{Нормированные координаты}{归一化坐标}

Для геометрического описания удобно ввести координаты
\zh{为了进行几何描述，方便引入坐标}
\[
    y_1=\frac{x_1}{\eps_1},
    \qquad
    y_2=\frac{x_2}{\eps_2}.
\]
Тогда
\zh{于是}
\[
    x_1=\eps_1y_1,
    \qquad
    x_2=\eps_2y_2.
\]
После устранения постоянного множителя в производственной функции система принимает вид
\zh{消去生产函数中的常数因子后，系统可写为}
\[
    \dot y_1=\frac{u}{\eps_1}F(y),
\]
\[
    \dot y_2=\frac{1-u}{\eps_2}F(y).
\]
Цель \(x_2(T)\to\max\) эквивалентна цели \(y_2(T)\to\max\), так как
\zh{目标 \(x_2(T)\to\max\) 等价于目标 \(y_2(T)\to\max\)，因为}
\[
    x_2(T)=\eps_2y_2(T).
\]
В координатах \(y\) особая магистраль имеет простой вид
\zh{在坐标 \(y\) 中，奇异主干线具有简单形式}
\[
    L_{\sng}^y=\{y\in\R^2_+: y_1=y_2\}.
\]

В этих координатах удобно положить
\zh{在这些坐标中，方便令}
\[
    q_1=\frac{\psi_1}{\eps_1},
    \qquad
    q_2=\frac{\psi_2}{\eps_2},
    \qquad
    \pi=q_1-q_2.
\]
Тогда
\zh{于是}
\[
    H(y,\psi,u)=F(y)(u\pi+q_2).
\]
Следовательно,
\zh{因此}
\[
    \pi>0\Rightarrow u=1,
    \qquad
    \pi<0\Rightarrow u=0,
\]
а на особой дуге
\zh{而在奇异弧上}
\[
    \pi\equiv0,
    \qquad
    y_1=y_2,
    \qquad
    u_{\sng}=\eps_1.
\]

\plainheading{Обоснование оптимальности}{最优性的证明思路}

Рассмотрим максимизированный гамильтониан
\zh{考虑最大化后的 Hamilton 函数}
\[
    M(x,\psi)=\max_{u\in[0,1]}H(x,\psi,u).
\]
Так как
\zh{由于}
\[
    H(x,\psi,u)=F(x)(u\Pi+\psi_2),
\]
получаем
\zh{得到}
\[
    M(x,\psi)=F(x)g(\psi),
\]
где
\zh{其中}
\[
    g(\psi)=\psi_2+\Pi^+,
    \qquad
    \Pi^+=\max\{\Pi,0\}.
\]
Эквивалентно,
\zh{等价地，}
\[
    g(\psi)=\max\{\psi_1,\psi_2\}
    =\frac{\psi_1+\psi_2+|\psi_1-\psi_2|}{2}.
\]
Для экстремальной тройки в рассматриваемой задаче получается
\zh{对于所考虑问题中的极值三元组，有}
\[
    g(\psi(t))>0,
    \qquad
    0\le t\le T.
\]

Пусть \((x(t),u(t),\psi(t))\) -- экстремальная тройка, а \((\widehat x(t),\widehat u(t))\) -- произвольный допустимый процесс с тем же начальным условием. Обозначим
\zh{设 \((x(t),u(t),\psi(t))\) 是极值三元组，\((\widehat x(t),\widehat u(t))\) 是具有相同初始条件的任意可行过程。记}
\[
    \Delta x(t)=\widehat x(t)-x(t),
    \qquad
    \Delta J=J[\widehat u]-J[u].
\]
Так как
\zh{因为}
\[
    J[u]=x_2(T),
    \qquad
    \psi(T)=
    \begin{pmatrix}
        0\\1
    \end{pmatrix},
\]
то
\zh{所以}
\[
    \Delta J=(\psi(T),\Delta x(T))-(\psi(0),\Delta x(0)).
\]
Поскольку \(\Delta x(0)=0\),
\zh{由于 \(\Delta x(0)=0\)，有}
\[
    \Delta J=\int_0^T \frac{d}{dt}(\psi(t),\Delta x(t))\,dt.
\]
Используя сопряженную систему и условие максимума, получаем
\zh{利用伴随系统和极大值条件，得到}
\[
    \Delta J\le
    \int_0^T
    \left[ F(\widehat x(t))-F(x(t))-\bigl(F'(x(t)),\widehat x(t)-x(t)\bigr)\right]
    g(\psi(t))\,dt.
\]
Так как функция Кобба--Дугласа \(F\) вогнута на \(\R^2_+\),
\zh{由于 Cobb--Douglas 函数 \(F\) 在 \(\R^2_+\) 上是凹的，}
\[
    F(\widehat x)-F(x)-\bigl(F'(x),\widehat x-x\bigr)\le0.
\]
Кроме того, \(g(\psi(t))>0\). Следовательно,
\zh{此外，\(g(\psi(t))>0\)。因此}
\[
    \Delta J\le0.
\]
То есть
\zh{也就是说}
\[
    J[\widehat u]\le J[u]
\]
для любого допустимого управления \(\widehat u\). Значит, экстремальный процесс является оптимальным.
\zh{对任意可行控制 \(\widehat u\) 都成立。因此，该极值过程就是最优过程。}

\plainheading{Качественный характер решения}{解的定性结构}

Для достаточно большого горизонта планирования оптимальная траектория имеет структуру
\zh{当规划期足够长时，最优轨线具有如下结构：}
\[
    \text{граничный участок}
    \quad\longrightarrow\quad
    \text{особая магистраль}
    \quad\longrightarrow\quad
    \text{терминальный участок}.
\]
В координатах \(y\) особая магистраль есть прямая
\zh{在坐标 \(y\) 中，奇异主干线是直线}
\[
    y_1=y_2.
\]

Если начальная точка уже лежит на особой магистрали,
\zh{如果初始点已经位于奇异主干线上，即}
\[
    y^0\in L_{\sng},
\]
то сначала движение идет по особой магистрали:
\zh{那么运动首先沿奇异主干线进行：}
\[
    u(t)=\eps_1,
\]
а затем на терминальном участке
\zh{随后在终端区段上取}
\[
    u(t)=0,
\]
поскольку конечная цель состоит в максимизации второй координаты \(y_2(T)\).
\zh{因为终端目标是最大化第二个坐标 \(y_2(T)\)。}

Если начальная точка лежит ниже особой магистрали, то сначала надо увеличить относительный уровень второго сектора, поэтому начальный участок имеет вид
\zh{如果初始点位于奇异主干线下方，那么首先需要提高第二部门的相对水平，因此初始区段为}
\[
    u(t)=0.
\]
Затем траектория выходит на особую магистраль и движется с управлением
\zh{之后轨线到达奇异主干线，并以控制}
\[
    u(t)=\eps_1
\]
движется. На последнем участке снова
\zh{继续运动。在最后区段再次取}
\[
    u(t)=0.
\]

Если начальная точка лежит выше особой магистрали, то сначала надо увеличить относительный уровень первого сектора, поэтому начальный участок имеет вид
\zh{如果初始点位于奇异主干线上方，那么首先需要提高第一部门的相对水平，因此初始区段为}
\[
    u(t)=1.
\]
Затем следует особая дуга
\zh{然后进入奇异弧}
\[
    u(t)=\eps_1,
\]
а в конце -- терминальный участок
\zh{最后进入终端区段}
\[
    u(t)=0.
\]

Если \(\tau\) -- момент входа на особую магистраль, а \(\theta\) -- момент выхода с нее на терминальный участок, то структура управления записывается так:
\zh{如果 \(\tau\) 是进入奇异主干线的时刻，\(\theta\) 是离开奇异主干线并进入终端区段的时刻，则控制结构可写为}
\[
    u(t)=
    \begin{cases}
        0\ \text{или}\ 1, & 0\le t<\tau,\\
        \eps_1, & \tau<t<\theta,\\
        0, & \theta<t\le T.
    \end{cases}
\]
Для случая, когда начальная точка уже лежит на особой магистрали, \(\tau=0\), и структура упрощается:
\zh{如果初始点已经位于奇异主干线上，则 \(\tau=0\)，结构简化为}
\[
    u(t)=
    \begin{cases}
        \eps_1, & 0\le t<\theta,\\
        0, & \theta<t\le T.
    \end{cases}
\]

Экономический смысл: сначала ресурсы перераспределяются так, чтобы вывести экономику на сбалансированную магистраль
\zh{其经济含义是：首先重新分配资源，使经济到达平衡主干线}
\[
    \frac{x_1}{\eps_1}=\frac{x_2}{\eps_2}.
\]
Затем экономика движется по особому режиму, где доля ресурсов, направляемая в первый сектор, равна его эластичности:
\zh{然后经济沿奇异机制运动，此时投向第一部门的资源份额等于该部门的弹性：}
\[
    u_{\sng}=\eps_1.
\]
На заключительном участке все ресурсы направляются во второй сектор:
\zh{在最后区段，全部资源投向第二部门：}
\[
    u=0,
\]
так как именно \(x_2(T)\) максимизируется.
\zh{因为目标正是最大化 \(x_2(T)\)。}

\plainheading{Краткая версия для запоминания}{记忆版要点}
\[
    \dot x_1=uF(x),
    \qquad
    \dot x_2=(1-u)F(x),
    \qquad
    J=x_2(T)\to\max.
\]
\[
    F(x)=Ax_1^{\eps_1}x_2^{\eps_2},
    \qquad
    \eps_1+\eps_2=1.
\]
\[
    H=F(x)(u(\psi_1-\psi_2)+\psi_2).
\]
\[
    \Pi=\psi_1-\psi_2.
\]
\[
    \Pi>0\Rightarrow u=1,
    \qquad
    \Pi<0\Rightarrow u=0.
\]
Особая дуга:
\zh{奇异弧：}
\[
    \Pi\equiv0\Rightarrow \psi_1=\psi_2,
    \quad
    \dot\psi_1=\dot\psi_2
    \Rightarrow
    F_{x_1}=F_{x_2}.
\]
\[
    \frac{\eps_1}{x_1}F=\frac{\eps_2}{x_2}F
    \Rightarrow
    \frac{x_1}{\eps_1}=\frac{x_2}{\eps_2}.
\]
\[
    \frac{\dot x_1}{\eps_1}=\frac{\dot x_2}{\eps_2}
    \Rightarrow
    \frac{u}{\eps_1}=\frac{1-u}{\eps_2}
    \Rightarrow
    u_{\sng}=\eps_1.
\]
Структура:
\zh{结构为：}
\[
    \text{boundary arc }(u=0\text{ or }1)
    \to
    \text{singular arc }(u=\eps_1)
    \to
    \text{terminal arc }(u=0).
\]
\zh{也就是：边界弧（$u=0$ 或 $u=1$）$\to$ 奇异弧（$u=\eps_1$）$\to$ 终端弧（$u=0$）。}



\newpage

\ticket{6}{Принцип максимума Понтрягина для задач на бесконечном интервале времени (общий случай). Основные соотношения принципа максимума. Условия трансверсальности на бесконечности.}{无限时间区间上最优控制问题的 Pontryagin 极大值原理（一般情形）。极大值原理的基本关系。无穷远处的横截条件。}

\qpart{1}{Принцип максимума Понтрягина на бесконечном интервале времени}{无限时间区间上的 Pontryagin 极大值原理}

\plainheading{Постановка задачи}{问题设定}

Рассмотрим задачу оптимального управления на бесконечном полуинтервале времени:
\zh{考虑定义在无限半轴上的最优控制问题：}
\[
    \dot x(t)=f(x(t),u(t)),
    \qquad
    u(t)\in U,
    \qquad
    x(0)=x_0,
    \qquad
    t\ge0,
\]
\[
    J(x,u)=\int_0^\infty e^{-\rho t}g(x(t),u(t))\,dt\to\max.
\]
Здесь \(x(t)\in G\subset\R^n\), \(u(t)\in U\subset\R^m\), \(U\) -- компакт, \(\rho\ge0\) -- параметр дисконтирования.
\zh{这里 \(x(t)\in G\subset\R^n\)，\(u(t)\in U\subset\R^m\)，集合 \(U\) 为紧集，\(\rho\ge0\) 为贴现参数。}

Допустимое управление -- измеримая функция \(u:[0,\infty)\to U\), для которой существует соответствующая траектория \(x(\cdot)\), удовлетворяющая системе. Допустимая пара \((x,u)\) называется оптимальной, если функционал \(J\) принимает на ней наибольшее значение.
\zh{可行控制是指一个可测函数 \(u:[0,\infty)\to U\)，并且存在与之对应、满足系统方程的轨线 \(x(\cdot)\)。如果泛函 \(J\) 在某个可行对 \((x,u)\) 上取得最大值，则称该可行对为最优的。}

Так как горизонт времени бесконечен, функционал является несобственным интегралом. Поэтому обычно предполагают условия, обеспечивающие его сходимость, например равномерную оценку хвоста:
\zh{由于时间区间是无限的，目标泛函是一个反常积分。因此通常需要假设保证其收敛的条件，例如如下的尾项一致估计：}
\[
    \int_T^\infty e^{-\rho t}|g(x(t),u(t))|\,dt\le\omega(T),
    \qquad
    \omega(T)\to0,
    \quad T\to\infty.
\]

\qpart{2}{Основные соотношения принципа максимума}{极大值原理的基本关系}

\plainheading{Основные соотношения принципа максимума}{极大值原理的基本关系}

Определим функцию Гамильтона--Понтрягина:
\zh{定义 Hamilton--Pontryagin 函数：}
\[
    H(x,t,u,\psi,\psi^0)
    =\langle f(x,u),\psi\rangle+\psi^0e^{-\rho t}g(x,u),
\]
и гамильтониан
\zh{并定义最大化后的 Hamilton 函数：}
\[
    H(x,t,\psi,\psi^0)=\sup_{u\in U}H(x,t,u,\psi,\psi^0).
\]
Здесь \(\psi(t)\in\R^n\) -- сопряженная переменная, \(\psi^0\ge0\) -- числовой множитель.
\zh{这里 \(\psi(t)\in\R^n\) 是伴随变量，\(\psi^0\ge0\) 是数值乘子。}

Пара \((\psi,\psi^0)\) должна быть нетривиальной:
\zh{伴随变量与乘子不能同时退化为零，即必须满足非平凡性条件：}
\[
    \|\psi(0)\|+\psi^0>0.
\]

Если \((x^*,u^*)\) -- оптимальная допустимая пара, то существует нетривиальная пара \((\psi,\psi^0)\), такая что выполняются основные соотношения ПМП:
\zh{如果 \((x^*,u^*)\) 是最优可行对，则存在非平凡对 \((\psi,\psi^0)\)，使得以下 Pontryagin 极大值原理的基本关系成立：}
\[
    \dot x^*(t)=f(x^*(t),u^*(t)),
\]
\[
    \dot\psi(t)=-\frac{\partial H}{\partial x}
    (x^*(t),t,u^*(t),\psi(t),\psi^0),
\]
то есть
\zh{也就是说，伴随方程可写为}
\[
    \dot\psi(t)
    =-\bigl(f_x(x^*(t),u^*(t))\bigr)^T\psi(t)
    -\psi^0e^{-\rho t}g_x(x^*(t),u^*(t)).
\]
Условие максимума:
\zh{极大值条件为：}
\[
    H(x^*(t),t,u^*(t),\psi(t),\psi^0)
    =\max_{u\in U}H(x^*(t),t,u,\psi(t),\psi^0)
\]
для почти всех \(t\ge0\).
\zh{该条件对几乎所有 \(t\ge0\) 成立。}

Если \(\psi^0>0\), можно перейти к нормальной форме и положить
\zh{如果 \(\psi^0>0\)，则可以转到正常形式，并令}
\[
    \psi^0=1.
\]
Если \(\psi^0=0\), имеет место анормальный случай, который нельзя заранее исключать без дополнительных предположений.
\zh{如果 \(\psi^0=0\)，则得到异常情形；在没有额外假设时，不能预先排除这种情况。}

\qpart{3}{Условия трансверсальности на бесконечности}{无穷远处的横截条件}

\plainheading{Условия трансверсальности на бесконечности}{无穷远处的横截条件}

Для задач на конечном интервале обычно есть конечные условия трансверсальности в правом конце. В задаче на бесконечном интервале таких условий автоматически нет.
\zh{对于有限区间上的问题，通常在右端点有相应的横截条件。而在无限时间区间问题中，这类端点条件不会自动出现。}

В приложениях часто используют условия вида
\zh{在应用中常使用如下形式的条件：}
\[
    \lim_{t\to\infty}\psi(t)=0
\]
или
\zh{或者}
\[
    \lim_{t\to\infty}\langle x^*(t),\psi(t)\rangle=0.
\]
Однако эти условия не являются следствием основных соотношений ПМП в общем случае. Они требуют дополнительных предположений.
\zh{但是在一般情形下，这些条件并不是 Pontryagin 极大值原理基本关系的直接推论。它们需要额外假设才能成立。}

Главное отличие бесконечного горизонта от конечного состоит в том, что нужно отдельно контролировать:
\zh{无限时间区间与有限时间区间的主要区别在于，需要另外控制以下三点：}
\begin{itemize}
    \item сходимость функционала;
    \zh{目标泛函的收敛性；}
    \item поведение сопряженной переменной на бесконечности;
    \zh{伴随变量在无穷远处的行为；}
    \item корректность выбранного условия трансверсальности.
    \zh{所选横截条件的正确性。}
\end{itemize}

\plainheading{Краткая версия для запоминания}{记忆版要点}

\[
    \dot x=f(x,u),
    \qquad
    J=\int_0^\infty e^{-\rho t}g(x,u)\,dt\to\max.
\]
\[
    H=\langle f(x,u),\psi\rangle+\psi^0e^{-\rho t}g(x,u).
\]
ПМП:
\zh{Pontryagin 极大值原理：}
\[
    \dot\psi=-H_x,
    \qquad
    H(x^*,u^*,\psi,\psi^0)=\max_{u\in U}H(x^*,u,\psi,\psi^0).
\]
Нетривиальность:
\zh{非平凡性条件：}
\[
    \|\psi(0)\|+\psi^0>0.
\]
Нормальный случай:
\zh{正常情形：}
\[
    \psi^0=1.
\]
Главная осторожность:
\zh{最需要注意的是：}
\[
    \lim_{t\to\infty}\psi(t)=0
\]
не является автоматическим условием в общем случае.
\zh{在一般情形下，上式并不是自动成立的横截条件。}

\newpage

\ticket{7}{Метод динамического программирования Беллмана и принцип максимума Понтрягина для задач на бесконечном интервале времени. Текущие сопряженные переменные. Экономический смысл соотношений принципа максимума.}{无限时间区间上 Bellman 动态规划方法与 Pontryagin 极大值原理。当前值伴随变量。极大值原理关系的经济含义。}

\qpart{1}{Метод динамического программирования Беллмана и принцип максимума Понтрягина}{Bellman 动态规划方法与 Pontryagin 极大值原理}

\plainheading{Метод Беллмана}{Bellman 方法}

Рассмотрим стационарную задачу на бесконечном горизонте:
\zh{考虑无限时间区间上的自治最优控制问题：}
\[
    \dot x(t)=f(x(t),u(t)),
    \qquad
    u(t)\in U,
    \qquad
    x(0)=x,
\]
\[
    J(x,u)=\int_0^\infty e^{-\rho t}g(x(t),u(t))\,dt\to\max.
\]

В методе динамического программирования вводится функция цены:
\zh{在动态规划方法中，引入价值函数：}
\[
    V(x)=\sup_{u(\cdot)}\int_0^\infty e^{-\rho t}g(x(t),u(t))\,dt.
\]
Она показывает максимальный дисконтированный выигрыш, который можно получить, если начальное состояние равно \(x\).
\zh{它表示当初始状态为 \(x\) 时可以获得的最大贴现收益。}

Если функция \(V\) достаточно гладкая, то для стационарной задачи она удовлетворяет уравнению Беллмана:
\zh{如果函数 \(V\) 足够光滑，则在自治情形下它满足 Bellman 方程：}
\[
    \rho V(x)=\max_{u\in U}\{g(x,u)+\langle V_x(x),f(x,u)\rangle\}.
\]
Если существует управление
\zh{如果存在控制}
\[
    u^*(x)\in\argmax_{u\in U}\{g(x,u)+\langle V_x(x),f(x,u)\rangle\},
\]
то при выполнении условий проверочной теоремы Беллмана обратная связь
\zh{那么在 Bellman 验证定理的条件成立时，反馈控制}
\[
    u(t)=u^*(x(t))
\]
является оптимальной.
\zh{就是最优控制。}

\qpart{2}{Текущие сопряженные переменные}{当前值伴随变量}

\plainheading{Связь с принципом максимума}{与极大值原理的联系}

Связь с ПМП задается через текущую сопряженную переменную. Пусть
\zh{它与 Pontryagin 极大值原理的联系通过当前值伴随变量给出。设}
\[
    p(t)=V_x(x^*(t)).
\]
Тогда \(p(t)\) называется текущей сопряженной переменной.
\zh{则 \(p(t)\) 称为当前值伴随变量。}

Она связана с обычной, то есть приведенной к начальному моменту, сопряженной переменной \(\psi(t)\) формулами
\zh{它与通常的、也就是折算到初始时刻的伴随变量 \(\psi(t)\) 之间有如下关系：}
\[
    \psi(t)=e^{-\rho t}p(t),
    \qquad
    p(t)=e^{\rho t}\psi(t).
\]

Введем гамильтониан в текущих ценах:
\zh{引入当前值 Hamilton 函数：}
\[
    M(x,u,p)=g(x,u)+\langle p,f(x,u)\rangle.
\]
Тогда условия ПМП в текущих переменных имеют вид
\zh{于是 Pontryagin 极大值原理在当前值变量中的形式为：}
\[
    \dot x^*(t)=f(x^*(t),u^*(t)),
\]
\[
    \dot p(t)=\rho p(t)-M_x(x^*(t),u^*(t),p(t)),
\]
то есть
\zh{即}
\[
    \dot p(t)=\rho p(t)-\bigl(f_x(x^*(t),u^*(t))\bigr)^Tp(t)-g_x(x^*(t),u^*(t)).
\]
Условие максимума:
\zh{极大值条件为：}
\[
    M(x^*(t),u^*(t),p(t))
    =\max_{u\in U}M(x^*(t),u,p(t))
\]
для почти всех \(t\ge0\).
\zh{该条件对几乎所有 \(t\ge0\) 成立。}

На оптимальной траектории уравнение Беллмана дает
\zh{在最优轨线上，Bellman 方程给出}
\[
    \rho V(x^*(t))=M(x^*(t),u^*(t),V_x(x^*(t))).
\]

\qpart{3}{Экономический смысл соотношений принципа максимума}{极大值原理关系的经济含义}

\plainheading{Экономический смысл}{经济含义}

Функция \(V(x)\) -- это внутренняя стоимость состояния \(x\), то есть максимальная будущая дисконтированная полезность. Градиент \(V_x(x)\) показывает предельную ценность дополнительной единицы состояния.
\zh{函数 \(V(x)\) 是状态 \(x\) 的内在价值，也就是从该状态出发能够得到的最大未来贴现效用。梯度 \(V_x(x)\) 表示状态增加一单位时的边际价值。}

В экономике \(p(t)\) интерпретируется как текущая теневая цена состояния, а \(\psi(t)\) -- как та же цена, приведенная к начальному моменту времени.
\zh{在经济学中，\(p(t)\) 可解释为状态的当前影子价格，而 \(\psi(t)\) 是同一个价格折算到初始时刻后的形式。}

Условие максимума означает, что оптимальное управление максимизирует сумму текущего выигрыша и прироста ценности состояния:
\zh{极大值条件的含义是：最优控制使“当前收益”和“状态价值增量”之和达到最大：}
\[
    g(x,u)+\langle p,f(x,u)\rangle.
\]

В терминах текущей сопряженной переменной условия трансверсальности часто записываются как
\zh{用当前值伴随变量表示时，横截条件常写为}
\[
    \lim_{t\to\infty}e^{-\rho t}p(t)=0
\]
или
\zh{或者}
\[
    \lim_{t\to\infty}e^{-\rho t}\langle x^*(t),p(t)\rangle=0.
\]
Но такие условия не являются автоматическими в общем случае и требуют дополнительных предположений.
\zh{但在一般情形下，这些条件并不是自动成立的，也需要额外假设。}

\plainheading{Важное замечание}{重要说明}

Метод Беллмана и ПМП имеют разный логический статус. ПМП обычно дает необходимые условия оптимальности. Проверочная теорема Беллмана при выполнении своих условий дает достаточное условие оптимальности.
\zh{Bellman 方法与 Pontryagin 极大值原理的逻辑地位不同。Pontryagin 极大值原理通常给出最优性的必要条件；而 Bellman 验证定理在其条件满足时给出最优性的充分条件。}

\plainheading{Краткая версия для запоминания}{记忆版要点}

\[
    V(x)=\sup_u\int_0^\infty e^{-\rho t}g(x(t),u(t))\,dt.
\]
Уравнение Беллмана:
\zh{Bellman 方程：}
\[
    \rho V(x)=\max_{u\in U}\{g(x,u)+\langle V_x(x),f(x,u)\rangle\}.
\]
Текущая сопряженная переменная:
\zh{当前值伴随变量：}
\[
    p(t)=V_x(x^*(t)).
\]
Связь переменных:
\zh{两个伴随变量之间的关系：}
\[
    \psi(t)=e^{-\rho t}p(t).
\]
Текущий гамильтониан:
\zh{当前值 Hamilton 函数：}
\[
    M(x,u,p)=g(x,u)+\langle p,f(x,u)\rangle.
\]
ПМП в current-value форме:
\zh{当前值形式下的 Pontryagin 极大值原理：}
\[
    \dot p=\rho p-M_x,
    \qquad
    M(x^*,u^*,p)=\max_{u\in U}M(x^*,u,p).
\]

\newpage

\ticket{8}{Достаточные условия оптимальности для задач на бесконечном интервале времени.}{无限时间区间上最优控制问题的最优性充分条件。}

\qpart{1}{Достаточные условия оптимальности для задач на бесконечном интервале времени}{无限时间区间问题的最优性充分条件}

\plainheading{Смысл вопроса}{问题含义}

Принцип максимума Понтрягина дает, вообще говоря, только необходимые условия оптимальности. Поэтому выполнение сопряженной системы и условия максимума само по себе не гарантирует оптимальность допустимого процесса. Для достаточности нужны дополнительные условия, обычно связанные с вогнутостью и корректным условием на бесконечности.
\zh{一般来说，Pontryagin 极大值原理只给出最优性的必要条件。因此，仅仅满足伴随系统和极大值条件，并不能保证一个可行过程就是最优的。为了得到充分性，还需要附加条件，这些条件通常与凹性以及无穷远处的正确条件有关。}

\plainheading{Постановка}{问题设定}

Рассмотрим задачу
\zh{考虑问题}
\[
    \dot x(t)=f(x(t),u(t)),
    \qquad
    u(t)\in U,
    \qquad
    x(0)=x_0,
\]
\[
    J(x,u)=\int_0^\infty e^{-\rho t}g(x(t),u(t))\,dt\to\max.
\]
Пусть \((x^*,u^*)\) -- допустимая пара. Предположим, что существует сопряженная переменная \(\psi(t)\), такая что в нормальной форме выполняются соотношения ПМП:
\zh{设 \((x^*,u^*)\) 是一个可行对。假设存在伴随变量 \(\psi(t)\)，使得正常形式下的 Pontryagin 极大值原理关系成立：}
\[
    \dot\psi(t)=-H_x(x^*(t),t,u^*(t),\psi(t)),
\]
\[
    H(x^*(t),t,u^*(t),\psi(t))
    =\max_{u\in U}H(x^*(t),t,u,\psi(t))
\]
для почти всех \(t\ge0\), где
\zh{其中上述条件对几乎所有 \(t\ge0\) 成立，并且}
\[
    H(x,t,u,\psi)=\langle f(x,u),\psi\rangle+e^{-\rho t}g(x,u).
\]

Введем максимизированный гамильтониан:
\zh{引入最大化后的 Hamilton 函数：}
\[
    \mathcal H(x,t,\psi(t))
    =\max_{u\in U}\{\langle f(x,u),\psi(t)\rangle+e^{-\rho t}g(x,u)\}.
\]

\plainheading{Достаточное условие оптимальности}{最优性的充分条件}

Одно из стандартных достаточных условий имеет следующий вид. Если:
\zh{一个标准的充分条件可以表述如下。如果满足：}
\begin{enumerate}
    \item множество состояний \(G\) выпукло;
    \zh{状态集合 \(G\) 是凸的；}
    \item для каждого \(t\ge0\) функция
    \zh{对每个 \(t\ge0\)，函数}
    \[
        x\mapsto \mathcal H(x,t,\psi(t))
    \]
    вогнута на \(G\);
    \zh{在 \(G\) 上关于 \(x\) 是凹函数；}
    \item для любой допустимой траектории \(x(\cdot)\) выполнено предельное условие
    \zh{对任意可行轨线 \(x(\cdot)\)，以下极限条件成立：}
    \[
        \liminf_{t\to\infty}\langle \psi(t),x(t)-x^*(t)\rangle\ge0;
    \]
\end{enumerate}
то допустимая пара \((x^*,u^*)\) является оптимальной.
\zh{那么可行对 \((x^*,u^*)\) 就是最优的。}

\plainheading{Идея доказательства}{证明思路}

Из условия максимума имеем
\zh{由极大值条件可得}
\[
    H(x^*(t),t,u^*(t),\psi(t))=\mathcal H(x^*(t),t,\psi(t)).
\]
Из вогнутости \(\mathcal H\) по \(x\) следует оценка первого порядка:
\zh{由 \(\mathcal H\) 关于 \(x\) 的凹性得到一阶估计：}
\[
    \mathcal H(x(t),t,\psi(t))-\mathcal H(x^*(t),t,\psi(t))
    \le
    \langle \mathcal H_x(x^*(t),t,\psi(t)),x(t)-x^*(t)\rangle.
\]
С учетом сопряженного уравнения эта оценка приводит к неравенству
\zh{结合伴随方程，上述估计推出不等式}
\[
    J_T(x,u)-J_T(x^*,u^*)
    \le
    -\langle \psi(T),x(T)-x^*(T)\rangle,
\]
где
\zh{其中}
\[
    J_T(x,u)=\int_0^T e^{-\rho t}g(x(t),u(t))\,dt.
\]
Переходя к пределу при \(T\to\infty\) и используя условие на бесконечности, получаем
\zh{令 \(T\to\infty\)，并使用无穷远处的条件，得到}
\[
    J(x,u)\le J(x^*,u^*).
\]
Следовательно, \((x^*,u^*)\) оптимальна.
\zh{因此，\((x^*,u^*)\) 是最优的。}

\plainheading{Достаточность через метод Беллмана}{通过 Bellman 方法得到充分性}

Другой способ получить достаточные условия -- метод Беллмана. Если существует гладкая функция \(V\), удовлетворяющая уравнению
\zh{得到充分条件的另一种方式是 Bellman 方法。如果存在光滑函数 \(V\)，满足方程}
\[
    \rho V(x)=\max_{u\in U}\{g(x,u)+\langle V_x(x),f(x,u)\rangle\},
\]
и управление \(u^*(x)\) достигает максимума в правой части, то при выполнении соответствующего условия на бесконечности обратная связь
\zh{并且控制 \(u^*(x)\) 使右端达到最大，那么在相应无穷远条件成立时，反馈控制}
\[
    u(t)=u^*(x(t))
\]
является оптимальной.
\zh{就是最优的。}

Таким образом, в задачах на бесконечном интервале достаточность обычно обеспечивается либо проверочной теоремой Беллмана, либо Arrow-type условием: нормальная форма ПМП, вогнутость максимизированного гамильтониана по фазовой переменной и корректное условие трансверсальности на бесконечности.
\zh{因此，在无限时间区间问题中，充分性通常由两类条件保证：一类是 Bellman 验证定理；另一类是 Arrow 型条件，即正常形式的 Pontryagin 极大值原理、最大化 Hamilton 函数关于状态变量的凹性，以及无穷远处正确的横截条件。}

\plainheading{Краткая версия для запоминания}{记忆版要点}

ПМП сам по себе -- необходимое условие, не достаточное.
\zh{Pontryagin 极大值原理本身只是必要条件，不是充分条件。}

Достаточность:
\zh{充分性可概括为：}
\[
\begin{aligned}
&\text{нормальная форма ПМП}
+\text{вогнутость }\mathcal H(x,t,\psi(t))\text{ по }x \\
&\quad +\text{условие на бесконечности}
\Rightarrow \text{оптимальность}.
\end{aligned}
\]
Главное предельное условие:
\zh{主要的极限条件是：}
\[
    \liminf_{t\to\infty}\langle \psi(t),x(t)-x^*(t)\rangle\ge0.
\]
Идея доказательства:
\zh{证明思路的核心不等式是：}
\[
    J_T(x,u)-J_T(x^*,u^*)
    \le
    -\langle\psi(T),x(T)-x^*(T)\rangle.
\]
После перехода к пределу:
\zh{取极限后得到：}
\[
    J(x,u)\le J(x^*,u^*).
\]

\blocktitle{C}{Геометрическая теория управления, сопряженные точки и особые траектории}{Билеты 9--10}
\zh{C板块：几何控制理论、共轭点与奇异轨线。第9--10题签。}

\ticket{9}{Введение в геометрическую теорию управления: гладкие многообразия, обыкновенные дифференциальные уравнения, векторные поля, коммутаторы, алгебра Ли векторных полей. Условие скобочной порождаемости для систем линейных по управлению.}{几何控制理论导论：光滑流形、常微分方程、向量场、交换子、向量场的 Lie 代数。控制线性系统的括号生成条件。}

\qpart{1}{Гладкие многообразия}{光滑流形}

\plainheading{Смысл вопроса}{问题含义}

В этом билете надо перейти от записи системы в \(\mathbb R^n\) к геометрическому языку. Состояние системы является точкой гладкого многообразия, правая часть дифференциального уравнения является векторным полем, а управляемая система представляет собой семейство векторных полей. Главная идея состоит в том, что коммутаторы векторных полей дают новые инфинитезимальные направления движения, которые могут отсутствовать среди мгновенно допустимых скоростей. Поэтому условие скобочной порождаемости связывает алгебру Ли векторных полей с управляемостью.
\zh{本题的关键是把系统在 \(\mathbb R^n\) 中的写法转化为几何语言。系统状态是光滑流形上的一个点，微分方程的右端是一个向量场，而受控系统则是一族向量场。核心思想是：向量场的交换子能够产生新的无穷小运动方向，这些方向可能并不属于瞬时可行速度集合。因此，括号生成条件把向量场的 Lie 代数与系统的可控性联系起来。}

\qpart{2}{Обыкновенные дифференциальные уравнения}{常微分方程}

\plainheading{Гладкие многообразия и касательные пространства}{光滑流形与切空间}

Гладкое \(n\)-мерное многообразие \(M\) --- это пространство, которое локально устроено как \(\mathbb R^n\) и снабжено согласованными гладкими координатными картами. Точка многообразия обозначается через \(q\in M\), а ее координаты в локальной карте --- через \(x=(x^1,\ldots,x^n)\).
\zh{光滑 \(n\) 维流形 \(M\) 是一个局部上与 \(\mathbb R^n\) 相同，并配有相容光滑坐标图的空间。流形上的点记为 \(q\in M\)，它在局部坐标图中的坐标记为 \(x=(x^1,\ldots,x^n)\)。}

Касательное пространство \(T_qM\) в точке \(q\in M\) --- это линейное пространство касательных векторов к гладким кривым, проходящим через \(q\). Если
\zh{点 \(q\in M\) 处的切空间 \(T_qM\) 是所有经过 \(q\) 的光滑曲线的切向量所组成的线性空间。如果}
\[
\gamma:(-\varepsilon,\varepsilon)\to M,
\qquad
\gamma(0)=q,
\]
то вектор скорости
\zh{那么速度向量}
\[
\dot\gamma(0)\in T_qM
\]
является касательным вектором в точке \(q\). В локальных координатах элементы \(T_qM\) имеют вид
\zh{就是点 \(q\) 处的切向量。在局部坐标中，\(T_qM\) 的元素可写为}
\[
\xi=\sum_{i=1}^n \xi^i\frac{\partial}{\partial x^i}\bigg|_q.
\]
Касательное расслоение обозначается
\zh{切丛记为}
\[
TM=\bigcup_{q\in M}T_qM.
\]

\qpart{3}{Векторные поля}{向量场}

\plainheading{Векторные поля, ОДУ и потоки}{向量场、常微分方程与流}

Гладкое векторное поле на \(M\) --- это гладкое отображение
\zh{流形 \(M\) 上的光滑向量场是一个光滑映射}
\[
X:M\to TM,
\qquad
X(q)\in T_qM.
\]
Множество всех гладких векторных полей на \(M\) обозначают \(\Vect(M)\).
\zh{\(M\) 上所有光滑向量场的集合记为 \(\Vect(M)\)。}

Векторное поле \(X\) задает обыкновенное дифференциальное уравнение на многообразии:
\zh{向量场 \(X\) 在流形上定义一个常微分方程：}
\[
\dot q=X(q),
\qquad
q\in M.
\]
Решением является кривая \(q(t)\), такая что
\zh{其解是一条曲线 \(q(t)\)，满足}
\[
\frac{dq(t)}{dt}=X(q(t)).
\]
Если решение зависит от начальной точки \(q_0\), то отображение
\zh{如果解依赖于初始点 \(q_0\)，则映射}
\[
P_X^t(q_0)=q(t;q_0)
\]
называется локальным потоком поля \(X\). Если поле полно, то поток определен для всех \(t\in\mathbb R\) и образует однопараметрическую группу диффеоморфизмов:
\zh{称为向量场 \(X\) 的局部流。如果向量场是完备的，则流对所有 \(t\in\mathbb R\) 都有定义，并形成一个单参数微分同胚群：}
\[
P_X^{t+s}=P_X^t\circ P_X^s,
\qquad
P_X^0=\operatorname{Id}.
\]

\qpart{4}{Коммутаторы}{交换子}

\plainheading{Управляемые системы на многообразиях}{流形上的受控系统}

В геометрической теории управления управляемая система рассматривается как семейство векторных полей
\zh{在几何控制理论中，受控系统被看作一族向量场：}
\[
\dot q=V_u(q),
\qquad
q\in M,
\quad u\in U,
\]
где \(M\) --- пространство состояний, \(U\) --- множество управляющих параметров.
\zh{其中 \(M\) 是状态空间，\(U\) 是控制参数集合。}

Если управление зависит от времени, \(u=u(t)\), то получаем неавтономное уравнение
\zh{如果控制依赖于时间，即 \(u=u(t)\)，则得到非自治方程}
\[
\dot q(t)=V_{u(t)}(q(t)).
\]
Траектории такой системы можно приближать траекториями, соответствующими кусочно-постоянным управлениям, то есть последовательным движениям вдоль потоков различных векторных полей.
\zh{这类系统的轨线可以用分段常值控制所对应的轨线来逼近，也就是沿不同向量场的流依次运动。}

Особенно важный класс --- системы, аффинные или линейные по управлению:
\zh{特别重要的一类系统是关于控制仿射或线性的系统：}
\[
\dot q=f_0(q)+\sum_{i=1}^m u_i f_i(q),
\qquad u=(u_1,\ldots,u_m).
\]
Если \(f_0=0\), система называется бездрейфовой:
\zh{如果 \(f_0=0\)，则称系统为无漂移系统：}
\[
\dot q=\sum_{i=1}^m u_i f_i(q).
\]
Для бездрейфовых систем условие скобочной порождаемости имеет прямую связь с управляемостью в смысле теоремы Чоу--Рашевского.
\zh{对于无漂移系统，括号生成条件与 Chow--Rashevskii 定理意义下的可控性有直接关系。}

\qpart{5}{Алгебра Ли векторных полей}{向量场的 Lie 代数}

\plainheading{Коммутатор векторных полей}{向量场的交换子}

Векторное поле можно понимать как дифференциальный оператор первого порядка на алгебре гладких функций:
\zh{向量场可以看作作用在光滑函数代数上的一阶微分算子：}
\[
X\varphi(q)=d\varphi(q)X(q),
\qquad
\varphi\in C^\infty(M).
\]
Коммутатором, или скобкой Ли, двух векторных полей \(X,Y\in\Vect(M)\) называется векторное поле
\zh{两个向量场 \(X,Y\in\Vect(M)\) 的交换子，也称 Lie 括号，是如下向量场：}
\[
[X,Y](\varphi)=X(Y\varphi)-Y(X\varphi),
\qquad
\varphi\in C^\infty(M).
\]
В локальных координатах, если
\zh{在局部坐标中，如果}
\[
X=\sum_{i=1}^n X^i(x)\frac{\partial}{\partial x^i},
\qquad
Y=\sum_{i=1}^n Y^i(x)\frac{\partial}{\partial x^i},
\]
то
\zh{则}
\[
[X,Y]=\sum_{i=1}^n
\left(
\sum_{j=1}^n
\left(
X^j\frac{\partial Y^i}{\partial x^j}
-
Y^j\frac{\partial X^i}{\partial x^j}
\right)
\right)
\frac{\partial}{\partial x^i}.
\]
Основные свойства скобки Ли:
\zh{Lie 括号的基本性质为：}
\[
[X,Y]=-[Y,X],
\qquad
[X,X]=0,
\]
\[
[X,[Y,Z]]+[Y,[Z,X]]+[Z,[X,Y]]=0.
\]
Последнее равенство называется тождеством Якоби.
\zh{最后一个等式称为 Jacobi 恒等式。}

\qpart{6}{Условие скобочной порождаемости для систем, линейных по управлению}{控制线性系统的括号生成条件}

\plainheading{Геометрический смысл коммутатора}{交换子的几何意义}

Коммутатор показывает некоммутативность потоков. Если за малое время двигаться последовательно вдоль \(X\), затем вдоль \(Y\), затем назад вдоль \(X\) и назад вдоль \(Y\), то главный нетривиальный член смещения имеет направление \([X,Y]\):
\zh{交换子反映了流的不可交换性。如果在很短时间内先沿 \(X\) 运动，再沿 \(Y\) 运动，然后沿 \(X\) 反向运动，最后沿 \(Y\) 反向运动，那么位移的主要非平凡项具有 \([X,Y]\) 的方向：}
\[
P_Y^{-\varepsilon}P_X^{-\varepsilon}P_Y^{\varepsilon}P_X^{\varepsilon}(q)
= q+\varepsilon^2[X,Y](q)+o(\varepsilon^2)
\]
в локальных координатах, с возможным изменением знака в зависимости от порядка композиции потоков.
\zh{这是局部坐标下的表达式；根据流的复合顺序不同，符号可能发生变化。}

Поэтому даже если мгновенные скорости системы лежат только в подпространстве
\zh{因此，即使系统的瞬时速度只位于子空间}
\[
\Delta(q)=\operatorname{span}\{f_1(q),\ldots,f_m(q)\}\subset T_qM,
\]
переключения управлений позволяют получать направления, связанные со скобками
\zh{通过控制的切换，仍然可以得到与 Lie 括号有关的方向：}
\[
[f_i,f_j],
\qquad
[f_i,[f_j,f_k]],
\qquad
\ldots
\]
Это ключевая идея неголономной управляемости: ограничения на мгновенные скорости не обязательно означают такие же ограничения на достижимые состояния.
\zh{这就是非完整约束可控性的核心思想：对瞬时速度的限制并不一定意味着对可达状态也有同样的限制。}

\plainheading{Алгебра Ли векторных полей}{向量场的 Lie 代数}

Пусть \(\mathcal F=\{f_1,\ldots,f_m\}\subset\Vect(M)\). Алгеброй Ли, порожденной этими векторными полями, называется наименьшее линейное пространство векторных полей, содержащее \(f_1,\ldots,f_m\) и замкнутое относительно скобки Ли:
\zh{设 \(\mathcal F=\{f_1,\ldots,f_m\}\subset\Vect(M)\)。由这些向量场生成的 Lie 代数，是包含 \(f_1,\ldots,f_m\) 并且对 Lie 括号封闭的最小向量场线性空间：}
\[
\Lie\{f_1,\ldots,f_m\}.
\]
В каждой точке \(q\in M\) можно рассмотреть подпространство
\zh{在每一点 \(q\in M\) 处，可以考虑如下子空间：}
\[
\Lie_q\{f_1,\ldots,f_m\}
=
\operatorname{span}\{X(q):X\in\Lie\{f_1,\ldots,f_m\}\}\subset T_qM.
\]
Оно порождается значениями исходных полей и всех повторных коммутаторов:
\zh{它由原始向量场的取值以及所有重复交换子的取值生成：}
\[
f_i(q),\quad [f_i,f_j](q),\quad [f_i,[f_j,f_k]](q),\quad \ldots
\]

\plainheading{Условие скобочной порождаемости}{括号生成条件}

Для бездрейфовой системы
\zh{对于无漂移系统}
\[
\dot q=\sum_{i=1}^m u_i f_i(q)
\]
условие скобочной порождаемости имеет вид
\zh{括号生成条件写为}
\[
\Lie_q\{f_1,\ldots,f_m\}=T_qM
\qquad \text{для всех }q\in M.
\]
Иначе говоря, исходные поля и их повторные коммутаторы должны порождать все касательное пространство в каждой точке.
\zh{换言之，在每一点上，原始向量场及其所有重复 Lie 括号必须生成整个切空间。}

Если \(\dim M=n\), то это условие эквивалентно ранговому условию
\zh{如果 \(\dim M=n\)，则该条件等价于秩条件：}
\[
\rank\{f_i(q),[f_i,f_j](q),[f_i,[f_j,f_k]](q),\ldots\}=n.
\]

В случае системы с дрейфом
\zh{对于带漂移项的系统}
\[
\dot q=f_0(q)+\sum_{i=1}^m u_i f_i(q)
\]
аналогичное условие
\zh{类似条件}
\[
\Lie_q\{f_0,f_1,\ldots,f_m\}=T_qM
\]
обычно является условием локальной достижимости или доступности, но не всегда полной управляемости. Наличие дрейфа \(f_0\) делает задачу тоньше: из ранга алгебры Ли не всегда следует возможность двигаться назад по времени вдоль дрейфа.
\zh{通常给出局部可达性或可及性条件，但并不总是意味着完全可控。漂移项 \(f_0\) 的存在使问题更加细致：Lie 代数满秩并不总能推出系统可以沿漂移方向反向运动。}

\plainheading{Теорема Чоу--Рашевского и управляемость}{Chow--Rashevskii 定理与可控性}

Стандартная формулировка для бездрейфовой системы такова. Пусть \(M\) --- связное гладкое многообразие, а система имеет вид
\zh{对于无漂移系统，标准表述如下。设 \(M\) 是连通光滑流形，系统具有形式}
\[
\dot q=\sum_{i=1}^m u_i f_i(q),
\]
где управления позволяют двигаться в обе стороны вдоль полей \(f_i\). Если для всех \(q\in M\)
\zh{并且控制允许沿每个向量场 \(f_i\) 的正反两个方向运动。如果对所有 \(q\in M\) 都有}
\[
\Lie_q\{f_1,\ldots,f_m\}=T_qM,
\]
то любые две точки одной связной компоненты можно соединить допустимой траекторией системы. В локальной форме это означает, что множество достижимости имеет полный размер.
\zh{那么同一连通分支中的任意两点都可以用系统的可行轨线连接。在局部形式下，这意味着可达集具有满维。}

Экзаменационный смысл: геометрическая теория управления переводит вопрос об управляемости в вопрос о ранге системы векторных полей и их скобок Ли. Если мгновенных направлений мало, но их скобки порождают всё касательное пространство, то система может быть управляема.
\zh{考试中应强调：几何控制理论把可控性问题转化为向量场系统及其 Lie 括号的秩问题。即使瞬时方向数量较少，只要它们的括号能够生成整个切空间，系统仍然可能是可控的。}

\plainheading{Пример: модель автомобиля}{例子：汽车模型}

Простейшая модель автомобиля на \(M=\mathbb R^2\times S^1\) задается векторными полями
\zh{在 \(M=\mathbb R^2\times S^1\) 上，最简单的汽车模型由以下向量场给出：}
\[
f_1(q)=
\begin{pmatrix}
\cos\theta\\
\sin\theta\\
0
\end{pmatrix},
\qquad
f_2(q)=
\begin{pmatrix}
0\\
0\\
1
\end{pmatrix},
\qquad q=(x_1,x_2,\theta).
\]
Здесь \(f_1\) соответствует движению вперед-назад, а \(f_2\) --- вращению автомобиля. Их коммутатор равен
\zh{这里 \(f_1\) 对应汽车的前后运动，\(f_2\) 对应汽车的转向或旋转。二者的交换子为}
\[
[f_1,f_2](q)=
\begin{pmatrix}
\sin\theta\\
-\cos\theta\\
0
\end{pmatrix}
\]
с точностью до знака в зависимости от принятого порядка скобки. Векторы
\zh{根据 Lie 括号顺序约定的不同，符号可能相差一个负号。向量}
\[
f_1(q),\quad f_2(q),\quad [f_1,f_2](q)
\]
порождают всё пространство \(T_qM\) размерности 3. Следовательно, хотя автомобиль не может двигаться мгновенно вбок, боковое направление возникает через маневр, соответствующий коммутатору.
\zh{生成整个三维切空间 \(T_qM\)。因此，虽然汽车不能瞬时横向运动，但横向方向可以通过与交换子相对应的机动动作产生。}

\plainheading{Краткая схема}{简要框架}
\begin{enumerate}
    \item Геометрическая теория управления рассматривает системы на гладких многообразиях.
    \zh{几何控制理论研究光滑流形上的系统。}
    \item Векторное поле \(X\) задает ОДУ \(\dot q=X(q)\) и поток \(P_X^t\).
    \zh{向量场 \(X\) 给出常微分方程 \(\dot q=X(q)\) 以及流 \(P_X^t\)。}
    \item Управляемая система --- семейство векторных полей: \(\dot q=V_u(q)\).
    \zh{受控系统是一族向量场：\(\dot q=V_u(q)\)。}
    \item Для системы, линейной по управлению: \(\dot q=f_0(q)+\sum u_i f_i(q)\).
    \zh{对于控制线性系统：\(\dot q=f_0(q)+\sum u_i f_i(q)\)。}
    \item Скобка Ли: \([X,Y](\varphi)=X(Y\varphi)-Y(X\varphi)\).
    \zh{Lie 括号为：\([X,Y](\varphi)=X(Y\varphi)-Y(X\varphi)\)。}
    \item Коммутаторы дают направления, достижимые за счет быстрых переключений управлений.
    \zh{交换子给出通过快速切换控制而可达的方向。}
    \item Условие скобочной порождаемости: \(\Lie_q\{f_1,\ldots,f_m\}=T_qM\).
    \zh{括号生成条件为：\(\Lie_q\{f_1,\ldots,f_m\}=T_qM\)。}
    \item Для бездрейфовых систем это ведет к управляемости по теореме Чоу--Рашевского; для систем с дрейфом обычно дает доступность, а не автоматически полную управляемость.
    \zh{对无漂移系统，它通过 Chow--Rashevskii 定理推出可控性；对带漂移系统，它通常给出可及性，而不自动推出完全可控性。}
\end{enumerate}

\newpage

\ticket{10}{Понятие сопряженной точки. Достаточные условия слабого минимума в классической задаче вариационного исчисления.}{共轭点的概念。经典变分法问题中弱极小的充分条件。}

\qpart{1}{Понятие сопряженной точки}{共轭点的概念}

\plainheading{Смысл вопроса}{问题含义}

Вопрос требует раскрыть два понятия: что такое сопряженная точка и какие условия обеспечивают слабый минимум в классической задаче вариационного исчисления. Основная логика:
\zh{本题需要说明两个概念：什么是共轭点，以及在经典变分法问题中哪些条件保证弱极小。基本逻辑是：}
\[
\begin{aligned}
&\text{уравнение Эйлера}
\Rightarrow \text{вторая вариация}
\Rightarrow \text{условие Лежандра} \\
&\Rightarrow \text{уравнение Якоби}
\Rightarrow \text{сопряженные точки}
\Rightarrow \text{слабый минимум}.
\end{aligned}
\]
\zh{Euler 方程 \(\Rightarrow\) 二阶变分 \(\Rightarrow\) Legendre 条件 \(\Rightarrow\) Jacobi 方程 \(\Rightarrow\) 共轭点 \(\Rightarrow\) 弱极小。}

\qpart{2}{Достаточные условия слабого минимума в классической задаче вариационного исчисления}{经典变分法问题中弱极小的充分条件}

\plainheading{Постановка задачи}{问题设定}

Рассмотрим простейшую классическую задачу вариационного исчисления:
\zh{考虑最简单的经典变分法问题：}
\[
J[x]=\int_{t_0}^{t_1} f(t,x(t),\dot x(t))\,dt\to\min,
\]
\[
x(t_0)=a,
\qquad
x(t_1)=b,
\]
где
\zh{其中}
\[
x(t)\in\mathbb R^n,
\qquad
f\in C^2.
\]

Кривая \(\widehat x(\cdot)\) называется слабым локальным минимумом, если существует \(\varepsilon>0\) такое, что для любой допустимой кривой \(x(\cdot)\), удовлетворяющей тем же граничным условиям и условию
\zh{如果存在 \(\varepsilon>0\)，使得对任何满足相同边界条件并满足下列条件的可行曲线 \(x(\cdot)\)，都有不等式成立，则曲线 \(\widehat x(\cdot)\) 称为弱局部极小：}
\[
\|x-\widehat x\|_{C^1}<\varepsilon,
\]
выполнено
\zh{即}
\[
J[x]\ge J[\widehat x].
\]
Иными словами, слабый минимум --- это минимум относительно малых \(C^1\)-возмущений.
\zh{换言之，弱极小是相对于小的 \(C^1\) 扰动而言的极小。}

\plainheading{Уравнение Эйлера--Лагранжа}{Euler--Lagrange 方程}

Первое необходимое условие экстремума --- уравнение Эйлера--Лагранжа:
\zh{极值的第一个必要条件是 Euler--Lagrange 方程：}
\[
\frac{d}{dt}f_{\dot x}(t,\widehat x,\dot{\widehat x})
-
f_x(t,\widehat x,\dot{\widehat x})=0.
\]
Кривая, удовлетворяющая этому уравнению и граничным условиям, называется экстремалью.
\zh{满足该方程和边界条件的曲线称为极值曲线。}

\plainheading{Вторая вариация}{二阶变分}

Для исследования минимума рассматриваются вариации
\zh{为了研究极小性，考虑如下变分：}
\[
x_\lambda(t)=\widehat x(t)+\lambda h(t),
\]
где
\zh{其中}
\[
h(t_0)=h(t_1)=0.
\]
На экстремали первая вариация равна нулю. Главную роль играет вторая вариация:
\zh{在极值曲线上，一阶变分等于零。因此关键作用由二阶变分承担：}
\[
\delta^2J[h]
=
\int_{t_0}^{t_1}
\left(
\langle A(t)\dot h,\dot h\rangle
+
2\langle C(t)\dot h,h\rangle
+
\langle B(t)h,h\rangle
\right)dt,
\]
где все матрицы вычисляются на экстремали \((t,\widehat x(t),\dot{\widehat x}(t))\):
\zh{其中所有矩阵都在极值曲线 \((t,\widehat x(t),\dot{\widehat x}(t))\) 上计算：}
\[
A(t)=f_{\dot x\dot x}(t,\widehat x,\dot{\widehat x}),
\]
\[
B(t)=f_{xx}(t,\widehat x,\dot{\widehat x}),
\]
\[
C(t)=f_{x\dot x}(t,\widehat x,\dot{\widehat x}).
\]
Матрицы \(A(t)\) и \(B(t)\) симметричны, а \(C(t)\) вообще говоря не обязана быть симметричной.
\zh{矩阵 \(A(t)\) 和 \(B(t)\) 是对称的，而 \(C(t)\) 一般不必对称。}

\plainheading{Условие Лежандра}{Legendre 条件}

Необходимое условие Лежандра для минимума имеет вид
\zh{极小的 Legendre 必要条件为}
\[
A(t)=f_{\dot x\dot x}(t,\widehat x,\dot{\widehat x})\ge0.
\]
Усиленное условие Лежандра:
\zh{加强 Legendre 条件为：}
\[
A(t)>0,
\]
то есть существует \(\alpha>0\) такое, что
\zh{也就是说，存在 \(\alpha>0\)，使得}
\[
\langle A(t)\xi,\xi\rangle\ge \alpha|\xi|^2
\]
для всех \(\xi\in\mathbb R^n\) и всех \(t\in[t_0,t_1]\).
\zh{对所有 \(\xi\in\mathbb R^n\) 以及所有 \(t\in[t_0,t_1]\) 成立。}

\plainheading{Уравнение Якоби и сопряженные точки}{Jacobi 方程与共轭点}

Рассмотрим присоединенную задачу для второй вариации. Ее уравнение Эйлера имеет вид
\zh{考虑与二阶变分相联系的变分问题。它的 Euler 方程为}
\[
-\frac{d}{dt}\left(A(t)\dot h(t)+C^T(t)h(t)\right)
+
C(t)\dot h(t)+B(t)h(t)=0.
\]
Это уравнение называется уравнением Якоби для исходной вариационной задачи.
\zh{该方程称为原变分问题的 Jacobi 方程。}

Точка \(\tau\in(t_0,t_1]\) называется сопряженной с точкой \(t_0\) вдоль экстремали \(\widehat x(\cdot)\), если существует нетривиальное решение \(h(t)\not\equiv0\) уравнения Якоби, такое что
\zh{如果存在 Jacobi 方程的非平凡解 \(h(t)\not\equiv0\)，并且满足}
\[
h(t_0)=0,
\qquad
h(\tau)=0,
\]
то точка \(\tau\in(t_0,t_1]\) называется сопряженной с точкой \(t_0\) вдоль экстремали \(\widehat x(\cdot)\).
\zh{则点 \(\tau\in(t_0,t_1]\) 称为沿极值曲线 \(\widehat x(\cdot)\) 与点 \(t_0\) 共轭的点。}

Иначе говоря, сопряженная точка --- это момент, в котором возникает нетривиальная вариация, обращающаяся в нуль в двух концах и удовлетворяющая уравнению Якоби.
\zh{换言之，共轭点是这样一个时刻：存在一个满足 Jacobi 方程、在两个端点为零的非平凡变分。}

Геометрически наличие сопряженной точки означает, что семейство близких экстремалей, выходящих из одной начальной точки, начинает иметь пересечение или фокусировку. Аналитически появление сопряженной точки связано с потерей положительной определенности второй вариации.
\zh{从几何上看，共轭点的存在意味着从同一起点出发的邻近极值曲线族开始相交或聚焦。从分析上看，共轭点的出现与二阶变分正定性的丧失有关。}

Условие Якоби для слабого минимума состоит в отсутствии точек, сопряженных с \(t_0\), на интервале
\zh{弱极小的 Jacobi 条件是：在下列区间上不存在与 \(t_0\) 共轭的点：}
\[
(t_0,t_1].
\]

\plainheading{Достаточное условие слабого минимума}{弱极小的充分条件}

Пусть \(\widehat x(\cdot)\) --- экстремаль, удовлетворяющая уравнению Эйлера--Лагранжа и заданным граничным условиям. Если:
\zh{设 \(\widehat x(\cdot)\) 是满足 Euler--Lagrange 方程和给定边界条件的极值曲线。如果：}
\begin{enumerate}
    \item выполнено усиленное условие Лежандра:
    \zh{加强 Legendre 条件成立：}
    \[
    A(t)=f_{\dot x\dot x}(t,\widehat x,\dot{\widehat x})>0;
    \]
    \item на интервале \((t_0,t_1]\) нет точек, сопряженных с \(t_0\),
    \zh{在区间 \((t_0,t_1]\) 上不存在与 \(t_0\) 共轭的点，}
\end{enumerate}
то вторая вариация положительно определена:
\zh{则二阶变分正定：}
\[
\delta^2J[h]>0
\]
для всех ненулевых допустимых вариаций \(h\), \(h(t_0)=h(t_1)=0\). Следовательно, экстремаль \(\widehat x(\cdot)\) доставляет строгий слабый локальный минимум функционала \(J\).
\zh{该不等式对所有非零可行变分 \(h\) 成立，其中 \(h(t_0)=h(t_1)=0\)。因此，极值曲线 \(\widehat x(\cdot)\) 给出泛函 \(J\) 的严格弱局部极小。}

Таким образом, для слабого минимума недостаточно только уравнения Эйлера. Нужно проверить вторые условия: усиленное условие Лежандра и отсутствие сопряженных точек. Эти условия обеспечивают положительную определенность второй вариации, а значит, малые \(C^1\)-отклонения от экстремали увеличивают значение функционала.
\zh{因此，仅有 Euler 方程不足以保证弱极小。还必须检查二阶条件：加强 Legendre 条件以及没有共轭点。这些条件保证二阶变分正定，从而说明对极值曲线的小 \(C^1\) 偏离会增大泛函值。}

\plainheading{Краткая схема}{简要框架}
\[
\text{Euler}+
\text{усиленное условие Лежандра}+
\text{нет сопряженных точек}
\Rightarrow
\text{строгий слабый минимум}.
\]
\zh{Euler 方程 + 加强 Legendre 条件 + 无共轭点 \(\Rightarrow\) 严格弱极小。}

\ticket{11}{Особые траектории. Метод построения особых управлений.}{奇异轨线。奇异控制的构造方法。}

\qpart{1}{Особые траектории}{奇异轨线}

\subsection*{Смысл вопроса}
\zh{问题含义}
Вопрос требует объяснить, что такое особое управление и особая траектория, почему в особом режиме условие максимума Понтрягина не определяет управление, и как управление строится из условий
\[
H_u=0,\qquad \frac{d}{dt}H_u=0,\qquad \frac{d^2}{dt^2}H_u=0,\ldots
\]
до появления управления в явном виде.
\zh{本题要求说明什么是奇异控制和奇异轨线，为什么在奇异状态下 Pontryagin 极大值原理中的最大化条件不能直接确定控制，以及如何从 \(H_u=0\) 及其逐次时间导数为零这些条件出发，一直推到控制变量显式出现为止。}

\qpart{2}{Метод построения особых управлений}{奇异控制的构造方法}

\subsection*{Система, линейная по управлению}
\zh{关于控制线性的系统}
Рассмотрим задачу оптимального управления, в которой управление входит в систему линейно. Для простоты сначала возьмем скалярное управление:
\[
\dot x(t)=f_0(t,x(t))+u(t)f_1(t,x(t)),
\qquad u(t)\in [a,b].
\]
В автономном случае:
\[
\dot x=f_0(x)+u f_1(x).
\]
Пусть Hamiltonian имеет вид
\[
H(x,\psi,u)=H_0(x,\psi)+uH_1(x,\psi),
\]
где
\[
H_0(x,\psi)=\langle \psi,f_0(x)\rangle,
\qquad
H_1(x,\psi)=\langle \psi,f_1(x)\rangle.
\]
Если имеется интегральный функционал, линейный по управлению, то соответствующий член также включается в $H_0$ и $H_1$. Главное свойство сохраняется: Hamiltonian является линейной функцией управления $u$.
\zh{考虑控制以线性方式进入系统的最优控制问题。为简明起见，先取标量控制 \(u(t)\in[a,b]\)。在自治情形下，系统写为 \(\dot x=f_0(x)+u f_1(x)\)。Hamiltonian 可写成 \(H=H_0+uH_1\)，其中 \(H_0=\langle\psi,f_0\rangle\)，\(H_1=\langle\psi,f_1\rangle\)。若积分型目标函数中也含有关于控制的线性项，则相应部分并入 \(H_0,H_1\)。关键点是：Hamiltonian 仍然是控制 \(u\) 的线性函数。}

\subsection*{Функция переключения}
\zh{切换函数}
Функция
\[
\Phi(t)=H_u(x(t),\psi(t),u(t))
\]
называется функцией переключения. В случае
\[
H=H_0+uH_1
\]
имеем
\[
\Phi(t)=H_1(x(t),\psi(t)).
\]
\zh{函数 \(\Phi(t)=H_u(x(t),\psi(t),u(t))\) 称为切换函数。在 \(H=H_0+uH_1\) 的线性情形下，切换函数就是 \(\Phi(t)=H_1(x(t),\psi(t))\)。}

По принципу максимума Понтрягина оптимальное управление должно максимизировать Hamiltonian по $u\in[a,b]$:
\[
H(x(t),\psi(t),u(t))=
\max_{v\in[a,b]}H(x(t),\psi(t),v).
\]
Если
\[
\Phi(t)>0,
\]
то максимум достигается при
\[
u(t)=b.
\]
Если
\[
\Phi(t)<0,
\]
то максимум достигается при
\[
u(t)=a.
\]
Это обычные bang-bang участки.
\zh{根据 Pontryagin 极大值原理，最优控制必须在区间 \([a,b]\) 上使 Hamiltonian 取最大值。因此，当 \(\Phi(t)>0\) 时最大值在 \(u(t)=b\) 处取得；当 \(\Phi(t)<0\) 时最大值在 \(u(t)=a\) 处取得。这些区间就是通常的 bang-bang 控制区间。}

\subsection*{Особый режим}
\zh{奇异状态}
Особый режим возникает, если на некотором интервале $I=[t_1,t_2]$ выполнено тождество
\[
\Phi(t)\equiv0,
\qquad t\in I.
\]
В этом случае условие максимума первого порядка не определяет управление, потому что Hamiltonian не зависит от $u$ в первом порядке. Соответствующее управление называется особым управлением, а соответствующая траектория $x(t)$ называется особой траекторией.

Итак, особая траектория --- это проекция на фазовое пространство такой экстремали принципа максимума, на которой условие максимума не задает управление однозначно на некотором интервале.
\zh{若在某个区间 \(I=[t_1,t_2]\) 上恒有 \(\Phi(t)\equiv0\)，则出现奇异状态。这时一阶最大化条件不能确定控制，因为 Hamiltonian 在一阶意义下不依赖于 \(u\)。对应的控制称为奇异控制，对应的轨线 \(x(t)\) 称为奇异轨线。换言之，奇异轨线就是极大值原理的某条极值曲线在相空间中的投影，并且在某个区间上最大化条件不能唯一给出控制。}

\subsection*{Метод построения особого управления}
\zh{构造奇异控制的方法}
Метод состоит в последовательном дифференцировании функции переключения. На особом интервале должно выполняться
\[
\Phi(t)=0.
\]
Так как это тождество по времени, то также
\[
\dot\Phi(t)=0,\qquad \ddot\Phi(t)=0,
\]
и так далее. Дифференцирование производится вдоль экстремальной системы:
\[
\dot x=H_\psi,
\qquad
\dot\psi=-H_x.
\]
\zh{方法的核心是对切换函数逐次求导。在奇异区间上必须有 \(\Phi(t)=0\)。既然这是关于时间的恒等式，也必须有 \(\dot\Phi(t)=0,\ddot\Phi(t)=0\)，依此类推。求导要沿着极值系统进行，即沿着 \(\dot x=H_\psi,\dot\psi=-H_x\) 计算。}

Если управление $u$ впервые явно появляется в производной порядка $k$, то из условия
\[
\left(\frac{d}{dt}\right)^k\Phi(t)=0
\]
можно выразить особое управление:
\[
u_{\sng}(t)=\cdots.
\]
\zh{如果控制 \(u\) 第一次在第 \(k\) 阶导数中显式出现，就可以由 \(\left(\frac{d}{dt}\right)^k\Phi(t)=0\) 解出奇异控制 \(u_{\sng}(t)\)。}

В типичной теории особых экстремалей для задач, линейных по управлению, первый порядок, в котором управление появляется существенно, является четным. Если управление впервые появляется в производной порядка $2q$, то число $q$ называется порядком особого управления или особой дуги.
\zh{在控制线性问题的典型奇异极值理论中，控制真正出现的最低阶通常是偶数阶。若控制第一次在 \(2q\) 阶导数中出现，则 \(q\) 称为该奇异控制或奇异弧的阶。}

Формально это записывают так:
\[
\frac{\partial}{\partial u}
\left(\frac{d}{dt}\right)^k H_u=0,
\qquad k=0,1,\ldots,2q-1,
\]
но
\[
\frac{\partial}{\partial u}
\left(\frac{d}{dt}\right)^{2q}H_u\ne0.
\]
Тогда особое управление находится из уравнения
\[
\left(\frac{d}{dt}\right)^{2q}H_u=0.
\]
После нахождения $u_{\sng}$ необходимо проверить:
\begin{enumerate}
    \item допустимость: $u_{\sng}(t)\in[a,b]$;
    \item выполнение исходных условий PMP;
    \item выполнение дополнительных необходимых условий второго порядка, например условия Kelley или generalized Legendre--Clebsch.
\end{enumerate}
\zh{形式上，在 \(k=0,1,\ldots,2q-1\) 阶时，对 \(u\) 的偏导仍为零，而在 \(2q\) 阶时不再为零。于是奇异控制由 \(\left(\frac{d}{dt}\right)^{2q}H_u=0\) 得到。求出 \(u_{\sng}\) 后，还必须检查它是否可行，即 \(u_{\sng}(t)\in[a,b]\)，是否满足原来的 PMP 条件，以及是否满足 Kelley 条件或广义 Legendre--Clebsch 条件等二阶必要条件。}

\subsection*{Условие Kelley}
\zh{Kelley 条件}
При максимизационной форме принципа максимума часто используется условие Kelley:
\[
(-1)^q
\frac{\partial}{\partial u}
\left(\frac{d}{dt}\right)^{2q}H_u
<0
\]
на особой дуге. Это условие является необходимым условием локальной оптимальности особой дуги. Однако знак этого условия зависит от соглашения о Hamiltonian и от того, рассматривается задача максимизации или минимизации. Поэтому при использовании другой конвенции знак может измениться.
\zh{在极大化形式的极大值原理中，常用 Kelley 条件作为奇异弧局部最优性的必要条件。需要注意，条件中的符号取决于 Hamiltonian 的约定，以及问题写成最大化还是最小化；若采用不同约定，不等号方向可能改变。}

\subsection*{Запись через скобки Пуассона}
\zh{用 Poisson 括号表示}
В автономном случае удобно использовать скобки Пуассона. Пусть
\[
h_0(x,\psi)=\langle \psi,f_0(x)\rangle,
\qquad
h_1(x,\psi)=\langle \psi,f_1(x)\rangle.
\]
Тогда
\[
H=h_0+u h_1,
\]
а switching function равна
\[
\Phi=h_1.
\]
Производная функции $h$ вдоль гамильтоновой системы выражается через скобку Пуассона:
\[
\dot h=\{H,h\}.
\]
Поэтому
\[
\dot\Phi=\{H,h_1\}.
\]
Так как
\[
\{h_1,h_1\}=0,
\]
получаем
\[
\dot\Phi=\{h_0,h_1\}.
\]
\zh{在自治情形下，用 Poisson 括号写法很方便。令 \(h_0=\langle\psi,f_0\rangle\)，\(h_1=\langle\psi,f_1\rangle\)，则 \(H=h_0+uh_1\)，切换函数为 \(\Phi=h_1\)。任意函数 \(h\) 沿 Hamilton 系统的导数可写为 \(\dot h=\{H,h\}\)。因此 \(\dot\Phi=\{H,h_1\}\)。由于 \(\{h_1,h_1\}=0\)，得到 \(\dot\Phi=\{h_0,h_1\}\)。}
Следующая производная:
\[
\ddot\Phi
=
\{H,\{h_0,h_1\}\}
=
\{h_0,\{h_0,h_1\}\}
+
u\{h_1,\{h_0,h_1\}\}.
\]
Если
\[
\{h_1,\{h_0,h_1\}\}\ne0,
\]
то особое управление первого порядка находится из условия $\ddot\Phi=0$:
\[
u_{\sng}
=
-
\frac{\{h_0,\{h_0,h_1\}\}}
{\{h_1,\{h_0,h_1\}\}}.
\]
После этого нужно проверить
\[
u_{\sng}\in[a,b]
\]
и вторые необходимые условия оптимальности.
\zh{再求一次导数，得到 \(\ddot\Phi=\{h_0,\{h_0,h_1\}\}+u\{h_1,\{h_0,h_1\}\}\)。若分母项 \(\{h_1,\{h_0,h_1\}\}\ne0\)，则一阶奇异控制可由 \(\ddot\Phi=0\) 显式求出。之后仍要检查 \(u_{\sng}\in[a,b]\) 以及二阶必要最优性条件。}

\subsection*{Векторное управление}
\zh{向量控制}
Если система имеет вид
\[
\dot x=f_0(x)+\sum_{i=1}^m u_i f_i(x),
\]
то
\[
H=h_0+\sum_{i=1}^m u_i h_i,
\qquad
h_i(x,\psi)=\langle \psi,f_i(x)\rangle.
\]
Функции $h_i(t)$ являются функциями переключения. Если на некотором интервале
\[
h_i(t)\equiv0
\]
для одной или нескольких компонент управления, то соответствующие компоненты являются особыми. Для их построения последовательно дифференцируют тождества $h_i=0$ до тех пор, пока соответствующие компоненты управления не появятся явно.
\zh{若系统含有向量控制 \(\dot x=f_0(x)+\sum_{i=1}^m u_i f_i(x)\)，则 \(H=h_0+\sum_{i=1}^m u_i h_i\)，其中 \(h_i=\langle\psi,f_i\rangle\)。各 \(h_i(t)\) 是相应控制分量的切换函数。若某些分量在某个区间上满足 \(h_i(t)\equiv0\)，这些分量就是奇异的；构造时要对恒等式 \(h_i=0\) 逐次求导，直到相应控制分量显式出现。}

\subsection*{Chattering / Fuller phenomenon}
\zh{Chattering / Fuller 现象}
В окрестности особых дуг могут возникать участки с очень частыми переключениями bang-bang управления. В некоторых задачах появляется бесконечное число переключений на конечном интервале времени; это явление называют chattering или явлением Фуллера. Для основного экзаменационного изложения достаточно знать, что оно связано с сопряжением bang-bang участков с особыми дугами и с условиями высокого порядка.
\zh{在奇异弧附近，bang-bang 控制可能出现非常频繁的切换。在某些问题中，有限时间区间内会出现无限多次切换，这称为 chattering 或 Fuller 现象。考试层面需要知道的是：它与 bang-bang 区间和奇异弧的拼接、高阶必要条件有关。}

\subsection*{Итог}
\zh{总结}
Особое управление не угадывается и не задается произвольно. Оно выводится из тождества функции переключения:
\[
H_u=0
\]
и его последовательных производных:
\[
H_u=0,
\quad \frac{d}{dt}H_u=0,
\quad
\frac{d^2}{dt^2}H_u=0,
\quad \ldots
\]
до появления управления. Затем найденное управление проверяется на допустимость и на выполнение дополнительных необходимых условий оптимальности.
\zh{奇异控制不是猜出来的，也不能任意指定。它由切换函数恒等式 \(H_u=0\) 及其逐次导数为零推出，直到控制变量显式出现。随后必须检查求得的控制是否可行，以及是否满足附加的必要最优性条件。}

\subsection*{Краткая схема}
\zh{简要框架}
\[
\Phi\equiv0
\Rightarrow
\Phi=\dot\Phi=\ddot\Phi=\cdots=0
\Rightarrow
u_{\sng}
\Rightarrow
\text{проверка допустимости и Kelley}.
\]
\zh{\(\Phi\equiv0 \Rightarrow \Phi=\dot\Phi=\ddot\Phi=\cdots=0 \Rightarrow u_{\sng} \Rightarrow\) 检查可行性和 Kelley 条件。}

\blocktitle{D}{Граничное управление, наблюдение и вариационный метод}{Билеты 12--14}
\section*{Замечание о надежности}
\zh{可靠性说明}
Настоящий текст является экзаменационной версией, составленной по официальным формулировкам билетов, студенческим конспектам 2017 и 2026 гг., материалам по функциональному анализу и методам оптимизации, а также по общим фактам теории управляемости волновых уравнений. Полный текст книги \textit{Васильев Ф.П., Куржанский М.А., Потапов М.М., Разгулин А.В. Приближенное решение двойственных задач управления и наблюдения} пока не использован, поэтому численные константы в неравенствах наблюдаемости и отдельные детали конечномерной аппроксимации желательно сверить по этому источнику.
\zh{本文是面向考试的版本，依据官方题目表、2017 与 2026 年学生笔记、泛函分析和优化方法材料，以及波动方程可控性理论的一般事实整理而成。Васильев、Куржанский、Потапов、Разгулин 关于控制与观测对偶问题近似解的完整著作尚未逐页使用，因此观测不等式中的具体常数和有限维近似的个别细节，最终最好再同该文献核对。}

\ticket{12}{Постановка задач граничного управления и наблюдения для волнового уравнения. Выбор функциональных пространств. Двойственность. Свойства управляемости и наблюдаемости и их связь с критическим моментом времени.}{波动方程边界控制与观测问题的设定。函数空间的选择。对偶性。可控性、可观测性及其与临界时间的关系。}

\qpart{1}{Постановка задач граничного управления и наблюдения для волнового уравнения}{波动方程边界控制与观测问题的设定}

Рассмотрим волновое уравнение на отрезке $(0,l)$:
\[
 y_{tt}(t,x)=y_{xx}(t,x)-q(x)y(t,x),
 \qquad 0<t<T,\quad 0<x<l.
\]
Задача граничного управления состоит в выборе граничных функций
\[
 y(t,0)=u_0(t),\qquad y(t,l)=u_1(t),
\]
так, чтобы решение, выходя из нулевого начального состояния
\[
 y(0,x)=0,\qquad y_t(0,x)=0,
\]
достигло в момент $T$ заданного конечного состояния
\[
 y(T,x)=f_0(x),\qquad y_t(T,x)=f_1(x).
\]
Обозначим
\[
 u=(u_0,u_1).
\]
Естественное пространство граничных управлений:
\[
 H=L_2(0,T)\times L_2(0,T).
\]
\zh{考虑区间 \((0,l)\) 上的波动方程。边界控制问题是选择边界函数 \(y(t,0)=u_0(t)\)、\(y(t,l)=u_1(t)\)，使从零初始状态出发的解在时刻 \(T\) 到达给定终端状态 \((f_0,f_1)\)。记 \(u=(u_0,u_1)\)，边界控制的自然空间为 \(H=L_2(0,T)\times L_2(0,T)\)。}

\qpart{2}{Выбор функциональных пространств}{函数空间的选择}

Так как управление задано на границе, решение имеет более слабую регулярность, чем при распределенном управлении. В слабой постановке удобно брать пространство конечных состояний
\[
 F=L_2(0,l)\times H^{-1}(0,l),
\]
а сопряженное пространство
\[
 F^*=H_0^1(0,l)\times L_2(0,l).
\]
Тогда оператор управления задается формулой
\[
 A:H\to F,\qquad Au=(y(T,\cdot),y_t(T,\cdot)).
\]
Задача точного управления принимает операторный вид
\[
 Au=f,\qquad f=(f_0,f_1)\in F.
\]
Система называется вполне управляемой в момент $T$, если
\[
 R(A)=F,
\]
то есть для любого конечного состояния $f\in F$ существует управление $u\in H$, переводящее систему в это состояние.
\zh{由于控制作用在边界上，解的正则性通常弱于分布控制情形。弱形式中常取终端状态空间 \(F=L_2(0,l)\times H^{-1}(0,l)\)，其对偶空间为 \(F^*=H_0^1(0,l)\times L_2(0,l)\)。控制算子 \(A:H\to F\) 将边界控制映到终端状态 \(Au=(y(T,\cdot),y_t(T,\cdot))\)。精确控制问题写成算子方程 \(Au=f\)。若 \(R(A)=F\)，即任意终端状态都可由某个控制达到，则系统在时刻 \(T\) 完全可控。}

Рассмотрим теперь сопряженную задачу наблюдения. Пусть
\[
 v=(v_0,v_1)\in F^*=H_0^1(0,l)\times L_2(0,l).
\]
Сопряженная волновая задача имеет вид
\[
 p_{tt}(t,x)=p_{xx}(t,x)-q(x)p(t,x),
\]
\[
 p(t,0)=p(t,l)=0,
\]
\[
 p(T,x)=v_0(x),\qquad p_t(T,x)=-v_1(x).
\]
Наблюдаемыми сигналами являются граничные производные:
\[
 g_0(t)=p_x(t,0),\qquad g_1(t)=-p_x(t,l).
\]
Иными словами,
\[
 A^*v=(p_x(\cdot,0),-p_x(\cdot,l)).
\]
Здесь
\[
 A^*:F^*\to H^*
\]
- оператор, сопряженный к оператору управления $A$.
\zh{再看对偶的观测问题。给定 \(v=(v_0,v_1)\in F^*\)，令 \(p\) 满足相应的伴随波动方程和齐次边界条件，终端数据为 \(p(T,x)=v_0(x)\)、\(p_t(T,x)=-v_1(x)\)。可观测信号是边界导数 \(p_x(t,0)\) 与 \(-p_x(t,l)\)，也就是 \(A^*v=(p_x(\cdot,0),-p_x(\cdot,l))\)。这里 \(A^*:F^*\to H^*\) 是控制算子 \(A\) 的伴随算子。}

\qpart{3}{Двойственность}{对偶性}

Связь управления и наблюдения выражается тождеством двойственности:
\[
 \langle Au,v\rangle_{F,F^*}=\langle u,A^*v\rangle_H.
\]
В развернутом виде это означает
\[
 \left\langle (y(T),y_t(T)),(v_0,v_1)\right\rangle_{F,F^*}
 =
 \int_0^T\left[u_0(t)p_x(t,0)-u_1(t)p_x(t,l)\right]dt.
\]
Это тождество показывает, что задача управления и задача наблюдения являются двойственными: управление описывает достижение конечного состояния, а наблюдение описывает восстановление сопряженного состояния по граничному сигналу.
\zh{控制与观测之间的联系由对偶恒等式 \(\langle Au,v\rangle_{F,F^*}=\langle u,A^*v\rangle_H\) 表示。展开后，终端状态与对偶终端数据的配对等于边界控制和边界观测信号的积分配对。这个恒等式说明：控制问题描述如何到达终端状态，观测问题描述如何由边界信号恢复伴随状态，两者互为对偶。}

Полная наблюдаемость сопряженной системы означает, что по сигналу
\[
 A^*v=(p_x(\cdot,0),-p_x(\cdot,l))
\]
можно устойчиво восстановить данные $v$. В сильной форме это выражается неравенством наблюдаемости:
\[
 \|A^*v\|_{H^*}^2\ge \mu \|v\|_{F^*}^2,
 \qquad \mu>0.
\]
Это неравенство означает, что ненулевое сопряженное состояние не может иметь нулевой или слишком малый наблюдаемый сигнал.
\zh{伴随系统完全可观测，意味着可以由观测信号 \(A^*v\) 稳定恢复数据 \(v\)。强形式上，这写成观测不等式 \(\|A^*v\|_{H^*}^2\ge \mu\|v\|_{F^*}^2\)，其中 \(\mu>0\)。它表示非零的伴随状态不可能产生零的或过小的观测信号。}

\qpart{4}{Управляемость, наблюдаемость и критический момент времени}{可控性、可观测性与临界时间}

Для волнового уравнения принципиальную роль играет конечная скорость распространения волн. Если время наблюдения или управления слишком мало, волна не успевает распространиться по всему отрезку, и точная управляемость невозможна. Для отрезка длины $l$ и двустороннего граничного управления критическое время обычно равно
\[
 T^*=l.
\]
Если управление или наблюдение проводится только с одного конца, критическое время обычно равно
\[
 T^*=2l.
\]
При
\[
 T>T^*
\]
выполняется неравенство наблюдаемости, и исходная система точно управляема. При
\[
 T<T^*
\]
точная управляемость и полная наблюдаемость нарушаются, поскольку часть энергии волны не успевает достичь наблюдаемой границы.
\zh{对波动方程来说，波的有限传播速度起决定作用。如果控制或观测时间太短，波尚未传遍整个区间，精确可控性就不可能成立。长度为 \(l\) 的区间上，双端边界控制的临界时间通常为 \(T^*=l\)；若只从一端控制或观测，临界时间通常为 \(T^*=2l\)。当 \(T>T^*\) 时观测不等式成立，原系统精确可控；当 \(T<T^*\) 时，部分波能量来不及到达可观测边界，因此精确可控性和完全可观测性失效。}

Основная схема билета:
\[
 \text{точная управляемость } A
 \quad\Longleftrightarrow\quad
 \text{наблюдаемость } A^*
 \quad\Longleftrightarrow\quad
 \text{неравенство наблюдаемости}.
\]
\zh{本题的核心链条是：控制算子 \(A\) 的精确可控性 \(\Longleftrightarrow\) 伴随算子 \(A^*\) 的可观测性 \(\Longleftrightarrow\) 观测不等式成立。}

\subsection*{Что важно помнить}
\zh{需要记住的要点}
\begin{enumerate}
\item Оператор управления: $A:H\to F$, $Au=(y(T),y_t(T))$.
\item Сопряженный оператор наблюдения: $A^*:F^*\to H^*$.
\item Главное тождество: $\langle Au,v\rangle=\langle u,A^*v\rangle$.
\item Управляемость и наблюдаемость связаны через неравенство наблюдаемости.
\item Критическое время возникает из-за конечной скорости распространения волн.
\end{enumerate}
\zh{控制算子为 \(A:H\to F\)，伴随观测算子为 \(A^*:F^*\to H^*\)。基本恒等式是 \(\langle Au,v\rangle=\langle u,A^*v\rangle\)。可控性和可观测性通过观测不等式联系起来；临界时间来自波的有限传播速度。}

\newpage
\ticket{13}{Неравенство наблюдаемости как критерий управляемости и наблюдаемости. Истокопредставимость нормальных решений задач управления и наблюдения при наличии неравенства наблюдаемости. Описание вариационного метода решения линейных операторных уравнений: условия применимости, алгоритм, сходимость.}{作为可控性与可观测性判据的观测不等式。在观测不等式成立时，控制与观测问题正规解的源表示性。线性算子方程变分解法：适用条件、算法与收敛性。}

\qpart{1}{Неравенство наблюдаемости как критерий управляемости и наблюдаемости}{作为可控性与可观测性判据的观测不等式}

Пусть $H$ и $F$ - гильбертовы пространства, а
\[
 A:H\to F
\]
- линейный ограниченный оператор. Задача управления записывается как операторное уравнение
\[
 Au=f,\qquad f\in F,
\]
где $u\in H$ - искомое управление.
\zh{设 \(H\) 和 \(F\) 为 Hilbert 空间，\(A:H\to F\) 为线性有界算子。控制问题可写成算子方程 \(Au=f\)，其中 \(f\in F\)，未知量 \(u\in H\) 是要寻找的控制。}

Сопряженный оператор имеет вид
\[
 A^*:F^*\to H^*.
\]
Задача наблюдения состоит в восстановлении элемента $v\in F^*$ по наблюдаемому сигналу
\[
 A^*v\in H^*.
\]
\zh{伴随算子写为 \(A^*:F^*\to H^*\)。观测问题是在给定可观测信号 \(A^*v\in H^*\) 的情况下恢复元素 \(v\in F^*\)。}

Система называется вполне управляемой, если
\[
 R(A)=F,
\]
то есть уравнение $Au=f$ разрешимо для любого $f\in F$.
\zh{若 \(R(A)=F\)，也就是方程 \(Au=f\) 对任意 \(f\in F\) 都可解，则称系统完全可控。}

Критерием точной управляемости является неравенство наблюдаемости:
\[
 \|A^*v\|_{H^*}^2\ge \mu\|v\|_{F^*}^2
 \qquad \forall v\in F^*,\quad \mu>0.
\]
Это неравенство означает, что ненулевое состояние сопряженной системы не может быть невидимым при наблюдении.
\zh{精确可控性的判据是观测不等式 \(\|A^*v\|_{H^*}^2\ge\mu\|v\|_{F^*}^2\)。它表示伴随系统的非零状态不可能在观测中完全不可见。}

В гильбертовых пространствах это связано с теоремой о замкнутом образе. Если
\[
 \|A^*v\|_{H^*}\ge c\|v\|_{F^*},
\]
то $A^*$ инъективен и его образ замкнут. Это эквивалентно сюръективности оператора $A$, то есть полной управляемости:
\[
 R(A)=F.
\]
Если же известно только
\[
 N(A^*)=\{0\},
\]
то это обычно дает лишь плотность образа $R(A)$, то есть приближенную управляемость, но не обязательно точную управляемость. Для точной управляемости нужна именно оценка снизу, то есть неравенство наблюдаемости.
\zh{在 Hilbert 空间中，这与闭值域定理有关。如果 \(\|A^*v\|_{H^*}\ge c\|v\|_{F^*}\)，则 \(A^*\) 单射且值域闭，这等价于 \(A\) 满射，即系统完全可控。若只知道 \(N(A^*)=\{0\}\)，通常只能得到 \(R(A)\) 稠密，也就是近似可控，而不一定精确可控。精确可控性需要真正的下界估计，即观测不等式。}

\qpart{2}{Истокопредставимость нормальных решений при наличии неравенства наблюдаемости}{观测不等式成立时正规解的源表示性}

Если уравнение
\[
 Au=f
\]
разрешимо, то среди всех его решений выделяют нормальное решение:
\[
 u_*=
 \operatorname*{arg\,min}\{\|u\|_H: Au=f\}.
\]
Нормальное решение имеет минимальную норму в пространстве управлений.
\zh{若方程 \(Au=f\) 可解，则在所有解中选取控制空间范数最小的解 \(u_*\)，称为正规解。正规解就是满足约束 \(Au=f\) 的最小范数控制。}

При выполнении неравенства наблюдаемости нормальное решение является истокопредставимым. Это означает, что существует элемент $v_*\in F^*$ такой, что
\[
 u_*=J_HA^*v_*,
\]
где
\[
 J_H:H^*\to H
\]
- изоморфизм Рисса.
\zh{当观测不等式成立时，正规解具有源表示性：存在 \(v_*\in F^*\)，使 \(u_*=J_HA^*v_*\)。这里 \(J_H:H^*\to H\) 是 Riesz 同构。}

Элемент $v_*$ можно найти из вариационной задачи. Рассмотрим функционал
\[
 I(v)=\|A^*v\|_{H^*}^2-2\langle f,v\rangle_{F,F^*},
 \qquad v\in F^*.
\]
Если известно априорное ограничение
\[
 \|v_*\|_{F^*}\le r,
\]
то минимизацию проводят на шаре
\[
 B_r=\{v\in F^*: \|v\|_{F^*}\le r\}.
\]
Искомый элемент $v_*$ является решением задачи
\[
 I(v)\to\min,\qquad v\in B_r.
\]
После его нахождения нормальное управление восстанавливается по формуле
\[
 u_*=J_HA^*v_*.
\]
\zh{元素 \(v_*\) 可由变分问题求出。考虑泛函 \(I(v)=\|A^*v\|_{H^*}^2-2\langle f,v\rangle_{F,F^*}\)。如果已知先验界 \(\|v_*\|_{F^*}\le r\)，就在球 \(B_r\) 上最小化该泛函。求得 \(v_*\) 后，再用 \(u_*=J_HA^*v_*\) 恢复正规控制。}

Смысл метода: вместо непосредственного решения уравнения $Au=f$ ищется сопряженная переменная $v$, для которой нормальное решение имеет источник
\[
 u_*=J_HA^*v.
\]
\zh{方法的意义在于：不直接求解 \(Au=f\)，而是先寻找伴随变量 \(v\)，使正规解可写成 \(u_*=J_HA^*v\) 的源形式。}

\qpart{3}{Вариационный метод решения линейных операторных уравнений: условия применимости, алгоритм, сходимость}{线性算子方程的变分解法：适用条件、算法与收敛性}

Если данные и оператор заданы приближенно, используют вариационный метод с конечномерными аппроксимациями. Пусть построены пространства
\[
 H_N\subset H,
 \qquad
 F_N^*\subset F^*,
\]
и операторы
\[
 A_N:H_N\to F_N,
 \qquad
 A_N^*:F_N^*\to H_N^*,
\]
так, чтобы сохранялась взаимная сопряженность:
\[
 \langle A_Nu_N,v_N\rangle=\langle u_N,A_N^*v_N\rangle.
\]
Это ключевое требование: нельзя независимо дискретизировать управление и наблюдение, иначе нарушается двойственность.
\zh{若数据和算子只近似给出，就使用带有限维近似的变分方法。构造有限维空间 \(H_N\subset H\)、\(F_N^*\subset F^*\)，以及近似算子 \(A_N\)、\(A_N^*\)。关键要求是保持互为伴随的离散对偶关系 \(\langle A_Nu_N,v_N\rangle=\langle u_N,A_N^*v_N\rangle\)。控制和观测不能彼此独立离散化，否则会破坏对偶性。}

Конечномерная вариационная задача имеет вид
\[
 I_N(v_N)=\|A_N^*v_N\|_{H_N^*}^2-2\langle f_N,v_N\rangle
 \to\min,
 \qquad v_N\in B_r^N.
\]
После решения этой квадратичной задачи управление восстанавливается:
\[
 u_N=J_{H_N}A_N^*v_N.
\]
\zh{有限维变分问题是最小化二次泛函 \(I_N(v_N)=\|A_N^*v_N\|_{H_N^*}^2-2\langle f_N,v_N\rangle\)。解出该二次问题后，用 \(u_N=J_{H_N}A_N^*v_N\) 恢复近似控制。}

Условия применимости метода:
\begin{enumerate}
\item операторное уравнение $Au=f$ разрешимо;
\item нормальное решение истокопредставимо: $u_*=J_HA^*v_*$;
\item выполнено неравенство наблюдаемости: $\|A^*v\|_{H^*}^2\ge \mu\|v\|_{F^*}^2$;
\item построены согласованные взаимно сопряженные аппроксимации $A_N,A_N^*$;
\item приближенные правые части сходятся: $f_N\to f$;
\item точность минимизации стремится к нулю.
\end{enumerate}
При выполнении этих условий приближенные управления сходятся к нормальному решению:
\[
 \|u_N-u_*\|_H\to0.
\]
\zh{适用条件包括：原算子方程 \(Au=f\) 可解；正规解具有源表示 \(u_*=J_HA^*v_*\)；观测不等式成立；近似算子 \(A_N,A_N^*\) 互为伴随；右端 \(f_N\to f\)；最小化误差趋于零。在这些条件下，近似控制收敛到正规解，即 \(\|u_N-u_*\|_H\to0\)。}

\subsection*{Что важно помнить}
\zh{需要记住的要点}
\begin{enumerate}
\item $A:H\to F$, $A^*:F^*\to H^*$.
\item Точная управляемость требует оценки снизу $\|A^*v\|\ge c\|v\|$.
\item $N(A^*)=\{0\}$ само по себе дает только приближенную управляемость, а не точную.
\item Нормальное решение имеет вид $u_*=J_HA^*v_*$.
\item $v_*$ ищется из квадратичного функционала $I(v)=\|A^*v\|^2-2\langle f,v\rangle$.
\end{enumerate}
\zh{要点是：\(A:H\to F\)，\(A^*:F^*\to H^*\)；精确可控性需要 \(\|A^*v\|\ge c\|v\|\) 这样的下界估计；仅有 \(N(A^*)=\{0\}\) 只保证近似可控；正规解写成 \(u_*=J_HA^*v_*\)，其中 \(v_*\) 由二次泛函最小化求得。}

\newpage
\ticket{14}{Вывод неравенства наблюдаемости для простейшего уравнения колебаний однородной струны. Реализация вариационного метода: конструкции конечномерных взаимно сопряжённых приближений к операторам управления и наблюдения; схема решения конечномерной задачи квадратичной минимизации.}{最简单齐次弦振动方程观测不等式的推导。变分方法的实现：控制算子与观测算子的有限维互伴随近似构造；有限维二次极小化问题的求解方案。}

\qpart{1}{Вывод неравенства наблюдаемости для уравнения колебаний однородной струны}{齐次弦振动方程观测不等式的推导}

Рассмотрим простейшее уравнение колебаний однородной струны:
\[
 p_{tt}(t,x)=p_{xx}(t,x),
 \qquad 0<t<T,\quad 0<x<l,
\]
с однородными граничными условиями
\[
 p(t,0)=p(t,l)=0.
\]
Это сопряженная задача наблюдения для гранично управляемого волнового уравнения. Терминальные данные задаются в момент $T$:
\[
 p(T,x)=v_0(x),
 \qquad
 p_t(T,x)=-v_1(x).
\]
Энергия решения:
\[
 E(t)=\int_0^l\left(p_t^2(t,x)+p_x^2(t,x)\right)dx.
\]
Для однородного волнового уравнения энергия сохраняется:
\[
 E(t)=E(T)=E(0).
\]
Наблюдаемыми величинами являются граничные производные:
\[
 p_x(t,0),\qquad -p_x(t,l).
\]
Поэтому оператор наблюдения имеет вид
\[
 A^*v=(p_x(\cdot,0),-p_x(\cdot,l)).
\]
\zh{考虑最简单的齐次弦振动方程 \(p_{tt}=p_{xx}\)，并取齐次边界条件 \(p(t,0)=p(t,l)=0\)。这是边界可控波动方程对应的伴随观测问题，终端数据在时刻 \(T\) 给定。解的能量为 \(E(t)=\int_0^l(p_t^2+p_x^2)\,dx\)，对齐次波动方程能量守恒。可观测量是边界导数 \(p_x(t,0)\) 和 \(-p_x(t,l)\)，因此 \(A^*v=(p_x(\cdot,0),-p_x(\cdot,l))\)。}

Для однородной струны используется энергетическое тождество или метод множителей. Идея состоит в том, что энергия внутри отрезка выражается через поток энергии через границу. Поскольку волны распространяются с конечной скоростью, за время
\[
 T>l
\]
при двустороннем наблюдении вся энергия успевает проявиться на концах струны.
\zh{对齐次弦，可用能量恒等式或乘子方法。基本思想是：区间内部的能量可以通过边界能流表示。由于波以有限速度传播，当双端观测时间满足 \(T>l\) 时，全部能量都来得及在弦的两端表现出来。}

В результате получается неравенство наблюдаемости:
\[
 \int_0^T\left(|p_x(t,0)|^2+|p_x(t,l)|^2\right)dt
 \ge
 \mu\left(\|v_0\|_{H_0^1(0,l)}^2+\|v_1\|_{L_2(0,l)}^2\right),
\]
где $\mu>0$.
\zh{于是得到观测不等式：边界导数的时间积分下界控制终端数据的能量范数，其中 \(\mu>0\)。}

В нормировке, используемой в студенческих материалах, для случая $q(x)=0$ указана оценка
\[
 \|A^*v\|^2\ge \mu\|v\|^2,
 \qquad
 \mu=\frac{2(T-l)}{l}>0,
 \qquad T>l.
\]
Точная числовая константа зависит от выбранной нормы энергии и нормировки оператора наблюдения, поэтому в окончательной версии ее желательно сверить с основной литературой 10. Существенно, что положительная константа существует именно при $T>l$.
\zh{在学生材料采用的规范化下，当 \(q(x)=0\) 时给出的估计为 \(\|A^*v\|^2\ge\mu\|v\|^2\)，其中 \(\mu=2(T-l)/l>0\)，要求 \(T>l\)。具体数值常数依赖于能量范数和观测算子的规范化，因此最终版本最好同主要文献核对；实质结论是：正的观测常数正是在 \(T>l\) 时存在。}

\qpart{2}{Конструкции конечномерных взаимно сопряжённых приближений к операторам управления и наблюдения}{控制算子与观测算子的有限维互伴随近似构造}

Это неравенство означает, что энергия любой ненулевой сопряженной волны обязательно наблюдается на границе. Следовательно, исходная задача граничного управления является точно управляемой при $T>l$.

Если же
\[
 T<l,
\]
то существует часть волны, которая не успевает достичь наблюдаемых концов за время $T$. Поэтому неравенство наблюдаемости нарушается и точная управляемость невозможна.
\zh{该不等式表示任何非零伴随波的能量都必然能在边界上被观测到。因此当 \(T>l\) 时，原边界控制问题精确可控。若 \(T<l\)，存在一部分波在时间 \(T\) 内来不及到达被观测的端点，所以观测不等式失效，精确可控性不可能成立。}

\subsection*{Конечномерные взаимно сопряженные приближения}
\zh{有限维互伴随近似}
Для численного решения операторного уравнения
\[
 Au=f
\]
строят конечномерные аппроксимации оператора управления и оператора наблюдения.

Пусть
\[
 H_N\subset H
\]
- конечномерное пространство управлений, а
\[
 F_N^*\subset F^*
\]
- конечномерное пространство сопряженных состояний. Нужно построить дискретные операторы
\[
 A_N:H_N\to F_N,
 \qquad
 A_N^*:F_N^*\to H_N^*,
\]
так, чтобы выполнялось точное дискретное тождество двойственности:
\[
 \langle A_Nu_N,v_N\rangle_{F_N,F_N^*}
 =
 \langle u_N,A_N^*v_N\rangle_{H_N,H_N^*}.
\]
Это и называется построением взаимно сопряженных приближений. Главное требование: оператор наблюдения $A_N^*$ должен быть действительно сопряженным к $A_N$ в выбранных конечномерных скалярных произведениях. Его нельзя выбирать независимо.
\zh{为数值求解 \(Au=f\)，构造控制算子和观测算子的有限维近似。设 \(H_N\subset H\) 是有限维控制空间，\(F_N^*\subset F^*\) 是有限维伴随状态空间，需要构造 \(A_N:H_N\to F_N\) 和 \(A_N^*:F_N^*\to H_N^*\)，使离散对偶恒等式严格成立。这就是互伴随近似的构造。核心要求是：在选定的有限维内积下，\(A_N^*\) 必须确实是 \(A_N\) 的伴随，不能独立选取。}

\qpart{3}{Схема решения конечномерной задачи квадратичной минимизации}{有限维二次极小化问题的求解方案}

В матричной форме, если скалярные произведения заданы матрицами Грама $G_H$ и $G_F$, сопряженный оператор записывается как
\[
 A_N^*=G_H^{-1}A_N^T G_F.
\]
Если базисы ортонормированы, эта формула упрощается до обычного транспонирования.
\zh{在矩阵形式中，若内积由 Gram 矩阵 \(G_H,G_F\) 给出，则伴随算子写为 \(A_N^*=G_H^{-1}A_N^T G_F\)。如果基是标准正交的，该公式退化为普通转置。}

\subsection*{Конечномерная квадратичная задача}
\zh{有限维二次问题}
После построения взаимно сопряженных операторов вариационный метод сводится к минимизации квадратичного функционала:
\[
 I_N(v_N)=\|A_N^*v_N\|_{H_N^*}^2
 -2\langle f_N,v_N\rangle_{F_N,F_N^*}
 \to\min.
\]
Если известна априорная оценка
\[
 \|v_*\|\le r,
\]
то минимизация ведется на шаре:
\[
 \|v_N\|\le r.
\]
В координатах
\[
 v_N=\sum_{j=1}^m c_j\varphi_j
\]
функционал принимает вид
\[
 I_N(c)=c^TG_Nc-2b_N^Tc,
\]
где $G_N$ - симметричная неотрицательно определенная матрица, соответствующая оператору $A_NA_N^*$, а $b_N$ - вектор правой части, соответствующий $f_N$.
\zh{构造出互伴随算子后，变分方法化为二次泛函最小化。若有先验估计 \(\|v_*\|\le r\)，则在球 \(\|v_N\|\le r\) 上最小化。把 \(v_N\) 展开为有限维基下的坐标 \(c\)，泛函成为 \(I_N(c)=c^TG_Nc-2b_N^Tc\)，其中 \(G_N\) 是与 \(A_NA_N^*\) 对应的对称半正定矩阵，\(b_N\) 是对应 \(f_N\) 的右端向量。}

Если $G_N$ положительно определена, минимум находится из системы линейных уравнений
\[
 G_Nc=b_N.
\]
Если есть ограничение $\|v_N\|\le r$, то решается конечномерная задача квадратичного программирования на шаре.

После нахождения $v_N$ управление восстанавливается по формуле
\[
 u_N=J_{H_N}A_N^*v_N.
\]
При согласованности аппроксимаций, выполнении неравенства наблюдаемости, сходимости данных $f_N\to f$ и стремлении ошибки минимизации к нулю получаем
\[
 \|u_N-u_*\|_H\to0.
\]
\zh{若 \(G_N\) 正定，最小值由线性方程组 \(G_Nc=b_N\) 给出。若有约束 \(\|v_N\|\le r\)，则求解球上的有限维二次规划问题。得到 \(v_N\) 后，通过 \(u_N=J_{H_N}A_N^*v_N\) 恢复控制。若近似相容、观测不等式成立、数据 \(f_N\to f\)，且最小化误差趋于零，则有 \(\|u_N-u_*\|_H\to0\)。}

\subsection*{Что важно помнить}
\zh{需要记住的要点}
\begin{enumerate}
\item Для однородной струны наблюдение идет через $p_x(t,0)$ и $-p_x(t,l)$.
\item При двустороннем наблюдении ключевое условие: $T>l$.
\item Студенческая нормировка дает $\mu=2(T-l)/l$, но константу надо сверять с первоисточником.
\item В конечномерной схеме главное - сохранить сопряженность: $\langle A_Nu_N,v_N\rangle=\langle u_N,A_N^*v_N\rangle$.
\item Конечная задача сводится к квадратичной минимизации $I_N(c)=c^TG_Nc-2b_N^Tc$.
\end{enumerate}
\zh{齐次弦的观测量是 \(p_x(t,0)\) 与 \(-p_x(t,l)\)。双端观测的关键条件是 \(T>l\)。学生材料中的规范化给出 \(\mu=2(T-l)/l\)，但常数应同原始文献核对。有限维方案中最重要的是保持伴随性：\(\langle A_Nu_N,v_N\rangle=\langle u_N,A_N^*v_N\rangle\)。最后的问题化为二次极小化 \(I_N(c)=c^TG_Nc-2b_N^Tc\)。}

\blocktitle{E}{Пространства Соболева и задача оптимального нагрева стержня}{Билеты 15--17}
\section*{Замечание об источниках и надежности}
\zh{可靠性说明}
Основная литература по конкретной задаче оптимального нагрева стержня и лекционный источник по билетам 16--17 в текущем наборе материалов не были найдены. Поэтому билеты 16--17 ниже составлены на основе официальной формулировки вопроса, студенческих конспектов, стандартной теории слабых решений параболических уравнений и общих методов оптимизации. Если в лекциях преподавателя уравнение или функционал записаны с иными коэффициентами, например \(x_t=a^2x_{ss}+u\) или \(J=\frac12\|x(T)-b\|^2\), соответствующие коэффициенты в сопряженной задаче и градиенте нужно изменить согласованно.
\zh{当前材料中没有找到关于最优加热杆具体问题的主要文献，也没有找到第16--17题对应的完整讲义来源。因此下面第16--17题依据官方题目、学生笔记、抛物型方程弱解的标准理论以及一般优化方法整理。如果教师讲义中方程或泛函采用不同系数，例如 \(x_t=a^2x_{ss}+u\) 或 \(J=\frac12\|x(T)-b\|^2\)，则伴随问题和梯度中的相应系数也需要同步调整。}

\newpage

\ticket{15}{Обобщенная производная по Соболеву. Пространства Соболева. Теоремы вложения.}{Sobolev 广义导数。Sobolev 空间。嵌入定理。}
\qpart{1}{Обобщенная производная по Соболеву}{Sobolev 广义导数}

Пусть \(\Omega\subset\R^n\) -- открытая область. Обозначим через
\[
C_c^\infty(\Omega)
\]
пространство бесконечно дифференцируемых функций с компактным носителем в \(\Omega\). Эти функции называются пробными функциями.
\zh{设 \(\Omega\subset\R^n\) 为开区域。记 \(C_c^\infty(\Omega)\) 为在 \(\Omega\) 内具有紧支集的无穷可微函数空间，这些函数称为试验函数。}

Пусть
\[
u\in L_{1,\mathrm{loc}}(\Omega).
\]
Для мультииндекса
\[
\alpha=(\alpha_1,\ldots,\alpha_n),\qquad
|\alpha|=\alpha_1+\cdots+\alpha_n
\]
обозначим
\[
D^\alpha=
\frac{\partial^{|\alpha|}}{\partial x_1^{\alpha_1}\cdots \partial x_n^{\alpha_n}}.
\]
\zh{设 \(u\in L_{1,\mathrm{loc}}(\Omega)\)。对多重指标 \(\alpha=(\alpha_1,\ldots,\alpha_n)\)，记 \(|\alpha|=\alpha_1+\cdots+\alpha_n\)，并用 \(D^\alpha\) 表示相应的 \(|\alpha|\) 阶偏导算子。}
Функция
\[
v_\alpha\in L_{1,\mathrm{loc}}(\Omega)
\]
называется обобщенной, или слабой, производной функции \(u\) порядка \(\alpha\), если для любой пробной функции
\[
\varphi\in C_c^\infty(\Omega)
\]
выполнено интегральное тождество
\[
\int_\Omega u(x)D^\alpha\varphi(x)\,dx
=
(-1)^{|\alpha|}
\int_\Omega v_\alpha(x)\varphi(x)\,dx.
\]
В этом случае пишут
\[
D^\alpha u=v_\alpha
\]
в обобщенном смысле.
\zh{若存在 \(v_\alpha\in L_{1,\mathrm{loc}}(\Omega)\)，使得对任意试验函数 \(\varphi\in C_c^\infty(\Omega)\) 都有上述积分恒等式，则称 \(v_\alpha\) 为 \(u\) 的 \(\alpha\) 阶广义导数或弱导数，并记为 \(D^\alpha u=v_\alpha\)。}

Это определение получается из формулы интегрирования по частям. Главное отличие состоит в том, что производная переносится с функции \(u\) на гладкую пробную функцию \(\varphi\). Поэтому обобщенная производная может существовать даже тогда, когда классической производной у функции нет.
\zh{这个定义来自分部积分公式。关键区别在于，把导数从可能不光滑的函数 \(u\) 转移到光滑的试验函数 \(\varphi\) 上。因此，即使函数没有经典导数，也可能存在广义导数。}

В одномерном случае, если
\[
f\in L_{1,\mathrm{loc}}(a,b),
\]
то функция \(\omega\in L_{1,\mathrm{loc}}(a,b)\) называется обобщенной производной \(f\), если
\[
\int_a^b f(x)\varphi'(x)\,dx
=
-
\int_a^b \omega(x)\varphi(x)\,dx
\]
для всех
\[
\varphi\in C_c^\infty(a,b).
\]
\zh{一维情形中，若 \(f\in L_{1,\mathrm{loc}}(a,b)\)，并且 \(\omega\in L_{1,\mathrm{loc}}(a,b)\) 对所有 \(\varphi\in C_c^\infty(a,b)\) 满足 \(\int_a^b f\varphi' dx=-\int_a^b \omega\varphi dx\)，则 \(\omega\) 称为 \(f\) 的广义导数。}

Основные свойства обобщенной производной:
\begin{enumerate}[leftmargin=2em]
\item обобщенная производная единственна с точностью до равенства почти всюду;
\item операция взятия обобщенной производной линейна:
\[
D^\alpha(\lambda u+
\mu v)=\lambda D^\alpha u+\mu D^\alpha v;
\]
\item если классическая производная существует и локально интегрируема, то она совпадает с обобщенной производной;
\item если на связном интервале \(u'=0\) в обобщенном смысле, то \(u\) равна константе почти всюду;
\item обобщенная производная позволяет рассматривать негладкие решения дифференциальных уравнений.
\end{enumerate}
\zh{广义导数的基本性质包括：在几乎处处相等意义下唯一；求广义导数是线性运算；若经典导数存在且局部可积，则它与广义导数一致；若连通区间上 \(u'=0\) 的广义意义成立，则 \(u\) 几乎处处为常数；广义导数使我们能够研究微分方程的非光滑解。}

\subsection*{Пространства Соболева}
\zh{Sobolev 空间}
Пусть
\[
1\le p\le\infty,
\qquad
m\in\mathbb N.
\]
Пространством Соболева \(W_p^m(\Omega)\) называется множество функций
\[
u\in L_p(\Omega),
\]
у которых все обобщенные производные
\[
D^\alpha u,
\qquad |\alpha|\le m,
\]
существуют и принадлежат \(L_p(\Omega)\). То есть
\[
W_p^m(\Omega)=
\{u\in L_p(\Omega):D^\alpha u\in L_p(\Omega),\ |\alpha|\le m\}.
\]
\zh{Sobolev 空间 \(W_p^m(\Omega)\) 是由这样的函数组成：函数本身属于 \(L_p(\Omega)\)，并且所有阶数不超过 \(m\) 的广义导数 \(D^\alpha u\) 都存在且属于 \(L_p(\Omega)\)。}

\qpart{2}{Пространства Соболева}{Sobolev 空间}

При \(1\le p<\infty\) норма задается формулой
\[
\|u\|_{W_p^m(\Omega)}=
\left(
\sum_{|\alpha|\le m}
\|D^\alpha u\|_{L_p(\Omega)}^p
\right)^{1/p}.
\]
При \(p=\infty\):
\[
\|u\|_{W_\infty^m(\Omega)}=
\max_{|\alpha|\le m}
\|D^\alpha u\|_{L_\infty(\Omega)}.
\]
\zh{当 \(1\le p<\infty\) 时，\(W_p^m(\Omega)\) 的范数由所有 \(|\alpha|\le m\) 的广义导数的 \(L_p\) 范数组合而成；当 \(p=\infty\) 时，则取这些导数的 \(L_\infty\) 范数最大值。}

При \(p=2\) пространство Соболева обозначают
\[
H^m(\Omega)=W_2^m(\Omega).
\]
Это гильбертово пространство со скалярным произведением
\[
(u,v)_{H^m(\Omega)}=
\sum_{|\alpha|\le m}
\int_\Omega D^\alpha u(x)D^\alpha v(x)\,dx.
\]
\zh{当 \(p=2\) 时，Sobolev 空间记作 \(H^m(\Omega)=W_2^m(\Omega)\)。这是 Hilbert 空间，其内积为所有 \(|\alpha|\le m\) 的广义导数内积之和。}

Например,
\[
H^1(a,b)=\{u\in L_2(a,b):u'\in L_2(a,b)\},
\]
и
\[
\|u\|_{H^1(a,b)}^2=
\int_a^b\left(u^2(x)+(u'(x))^2\right)dx.
\]
\zh{例如，\(H^1(a,b)\) 由平方可积且弱导数也平方可积的函数组成，范数平方为函数本身和导数的 \(L_2\) 范数平方之和。}

Пространство
\[
H_0^1(\Omega)
\]
определяется как замыкание \(C_c^\infty(\Omega)\) в норме \(H^1(\Omega)\):
\[
H_0^1(\Omega)=\overline{C_c^\infty(\Omega)}^{\,H^1}.
\]
Для достаточно регулярной области это означает, что функции из \(H_0^1(\Omega)\) имеют нулевой след на границе:
\[
u|_{\partial\Omega}=0
\]
в смысле следа.
\zh{空间 \(H_0^1(\Omega)\) 定义为 \(C_c^\infty(\Omega)\) 在 \(H^1(\Omega)\) 范数下的闭包。对足够规则的区域，这意味着其中的函数在边界上的迹为零。}

\subsection*{Одномерный случай}
\zh{一维情形}
Если
\[
u\in H^1(a,b),
\]
то у \(u\) есть абсолютно непрерывный представитель, и для любых
\[
\alpha,\beta\in[a,b]
\]
выполнено
\[
u(\beta)-u(\alpha)=\int_\alpha^\beta u'(x)\,dx.
\]
Отсюда следует непрерывное вложение
\[
H^1(a,b)\hookrightarrow C[a,b].
\]
То есть каждая функция из \(H^1(a,b)\) имеет непрерывного представителя, и существует константа \(C>0\), такая что
\[
\|u\|_{C[a,b]}\le C\|u\|_{H^1(a,b)}.
\]
Это важно для краевых задач: в одномерном случае значения функции Соболева в точках имеют корректный смысл.
\zh{若 \(u\in H^1(a,b)\)，则 \(u\) 有绝对连续代表，并满足 Newton--Leibniz 公式。因此有连续嵌入 \(H^1(a,b)\hookrightarrow C[a,b]\)，即每个 \(H^1\) 函数都有连续代表，且 \(\|u\|_{C[a,b]}\le C\|u\|_{H^1(a,b)}\)。这对边值问题很重要，因为一维中 Sobolev 函数在点上的取值有严格意义。}

\subsection*{Теоремы вложения Соболева}
\zh{Sobolev 嵌入定理}
Теоремы вложения Соболева показывают, что наличие обобщенных производных повышает регулярность функции. Пусть \(\Omega\subset\R^n\) -- ограниченная область с достаточно регулярной границей.
\zh{Sobolev 嵌入定理说明：拥有一定阶数的广义导数会提高函数的正则性。下面设 \(\Omega\subset\R^n\) 是边界足够规则的有界区域。}

\qpart{3}{Теоремы вложения}{嵌入定理}

Если
\[
mp<n,
\]
то имеет место непрерывное вложение
\[
W_p^m(\Omega)\hookrightarrow L_q(\Omega)
\]
для подходящих \(q\), в частности для
\[
q\le p^*=\frac{np}{n-mp}.
\]
Частный случай при \(m=1\):
\[
W_p^1(\Omega)\hookrightarrow L_{p^*}(\Omega),
\qquad
p^*=\frac{np}{n-p},
\qquad p<n.
\]
\zh{若 \(mp<n\)，则有连续嵌入 \(W_p^m(\Omega)\hookrightarrow L_q(\Omega)\)，其中 \(q\) 可取到临界指数 \(p^*=\frac{np}{n-mp}\)。特别地，当 \(m=1\)、\(p<n\) 时，有 \(W_p^1(\Omega)\hookrightarrow L_{p^*}(\Omega)\)，\(p^*=\frac{np}{n-p}\)。}

Если
\[
mp>n,
\]
то функции из \(W_p^m(\Omega)\) имеют более регулярные представители. В частности, при достаточном превышении гладкости над размерностью возникают вложения в пространства непрерывных или гёльдеровых функций:
\[
W_p^m(\Omega)\hookrightarrow C^{0,\gamma}(\overline\Omega)
\]
для некоторого \(\gamma>0\). В частности, при \(m=1\), \(p>n\) имеем вложение
\[
W_p^1(\Omega)\hookrightarrow C(\overline\Omega).
\]
\zh{若 \(mp>n\)，则 \(W_p^m(\Omega)\) 中的函数有更规则的代表。光滑性相对维数足够高时，会嵌入连续函数空间或 Hölder 空间，例如 \(W_p^m(\Omega)\hookrightarrow C^{0,\gamma}(\overline\Omega)\)。特别地，当 \(m=1,p>n\) 时，有 \(W_p^1(\Omega)\hookrightarrow C(\overline\Omega)\)。}

Для \(p=2\) часто используют следующую форму. Если
\[
m-\frac n2>0,
\]
то \(H^m(\Omega)\) вкладывается в пространство непрерывных функций при достаточной регулярности \(\Omega\).
\zh{对 \(p=2\) 常用的表述是：若 \(m-\frac n2>0\)，且区域足够规则，则 \(H^m(\Omega)\) 嵌入连续函数空间。}

Кроме непрерывных вложений существуют компактные вложения Реллиха--Кондрашова. В типичной форме
\[
W_p^m(\Omega)\Subset W_q^l(\Omega)
\]
при более строгом условии регулярности, например когда
\[
m>l,
\qquad
m-\frac np>l-\frac nq.
\]
Здесь знак \(\Subset\) означает компактное вложение: из любой ограниченной последовательности в более сильном пространстве можно выбрать сходящуюся подпоследовательность в более слабом пространстве.
\zh{除连续嵌入外，还有 Rellich--Kondrachov 紧嵌入。例如在更强的正则性条件下 \(W_p^m(\Omega)\Subset W_q^l(\Omega)\)。符号 \(\Subset\) 表示紧嵌入：强空间中的任意有界序列，都能在较弱空间中选出收敛子列。}

Для задач математической физики эти теоремы имеют принципиальное значение. Они позволяют корректно задавать граничные условия через следы, доказывать существование слабых решений, получать компактность минимизирующих последовательностей и переходить к пределу в вариационных задачах и задачах оптимального управления.
\zh{这些定理对数学物理问题非常关键。它们使我们能够通过迹严格给出边界条件，证明弱解存在性，得到极小化序列的紧性，并在变分问题和最优控制问题中合法地取极限。}

\subsection*{Связь с параболическими задачами}
\zh{与抛物型问题的联系}
В задачах типа уравнения теплопроводности часто используется гильбертов тройник
\[
H_0^1(0,l)\subset L_2(0,l)\simeq L_2^*(0,l)\subset H^{-1}(0,l).
\]
Если
\[
y\in L_2(0,T;H_0^1(0,l)),
\qquad
y_t\in L_2(0,T;H^{-1}(0,l)),
\]
то
\[
y\in C([0,T];L_2(0,l)).
\]
Именно это позволяет корректно говорить о начальном условии \(y(0)\) и терминальном значении \(y(T)\) в задачах управления нагревом стержня.
\zh{热传导方程这类抛物型问题中常使用 Hilbert 三元组 \(H_0^1(0,l)\subset L_2(0,l)\simeq L_2^*(0,l)\subset H^{-1}(0,l)\)。若 \(y\in L_2(0,T;H_0^1)\) 且 \(y_t\in L_2(0,T;H^{-1})\)，则 \(y\in C([0,T];L_2)\)。这正是加热杆控制问题中能够严格谈论初始值 \(y(0)\) 和终端值 \(y(T)\) 的原因。}

\subsection*{Краткий итог}
\zh{简要总结}
Обобщенная производная по Соболеву -- это производная, определенная через интегральное тождество с пробными функциями. Пространства Соболева \(W_p^m(\Omega)\) состоят из функций, у которых обобщенные производные до порядка \(m\) лежат в \(L_p\). Теоремы вложения показывают, в какие более регулярные пространства попадают функции Соболева при заданных \(m,p,n\). Эти понятия являются основой слабой постановки краевых задач для уравнений в частных производных.
\zh{Sobolev 广义导数是通过试验函数积分恒等式定义的导数。Sobolev 空间 \(W_p^m(\Omega)\) 由广义导数直到 \(m\) 阶都属于 \(L_p\) 的函数组成。嵌入定理说明，在给定 \(m,p,n\) 时 Sobolev 函数可进入哪些更规则的空间。这些概念是偏微分方程边值问题弱形式的基础。}

\newpage

\ticket{16}{Постановка задачи об оптимальном нагреве стержня. Существование и единственность обобщенного решения краевой задачи. Теорема Вейерштрасса.}{最优加热杆问题的设定。边值问题广义解的存在唯一性。Weierstrass 定理。}
\qpart{1}{Постановка задачи об оптимальном нагреве стержня}{最优加热杆问题的设定}

Рассмотрим стержень длины \(l\). Пусть
\[
Q=(0,T)\times(0,l),
\]
где \(t\in(0,T)\) -- время, \(s\in(0,l)\) -- пространственная координата. Через
\[
x=x(s,t)
\]
обозначим температуру стержня, а через
\[
u=u(s,t)
\]
-- распределенное управление, то есть интенсивность нагрева.
\zh{考虑长度为 \(l\) 的杆，记 \(Q=(0,T)\times(0,l)\)，其中 \(t\) 是时间，\(s\) 是空间坐标。用 \(x=x(s,t)\) 表示杆的温度，用 \(u=u(s,t)\) 表示分布控制，也就是加热强度。}

Задача оптимального нагрева стержня имеет вид
\[
J(u)=\int_0^l\bigl(x(s,T;u)-b(s)\bigr)^2\,ds\to\inf_{u\in U},
\]
где \(b(s)\) -- заданное желаемое распределение температуры в конечный момент времени \(T\).
\zh{最优加热问题要求最小化终端温度 \(x(s,T;u)\) 与给定目标温度分布 \(b(s)\) 的 \(L_2\) 距离平方。}

Состояние \(x\) определяется уравнением теплопроводности
\[
x_t(s,t)=x_{ss}(s,t)+u(s,t),
\qquad (s,t)\in Q,
\]
с краевыми условиями
\[
x(0,t)=x(l,t)=0,
\qquad 0<t<T,
\]
и начальным условием
\[
x(s,0)=\varphi(s),
\qquad 0<s<l.
\]
Множество допустимых управлений задается как замкнутый шар в \(L_2(Q)\):
\[
U=\{u\in L_2(Q):\|u\|_{L_2(Q)}\le R\}.
\]
Физический смысл задачи: нужно подобрать распределенный нагрев \(u(s,t)\), ограниченный по энергии, так чтобы конечная температура \(x(s,T)\) была как можно ближе к заданному профилю \(b(s)\) в норме \(L_2(0,l)\).
\zh{状态由热传导方程 \(x_t=x_{ss}+u\) 决定，并满足零边界条件和初始条件 \(x(s,0)=\varphi(s)\)。允许控制集取为 \(L_2(Q)\) 中的闭球 \(\|u\|_{L_2(Q)}\le R\)。物理意义是：在加热能量受限的情况下，选择分布加热 \(u(s,t)\)，使终端温度尽可能接近期望剖面 \(b(s)\)。}

\subsection*{Функциональные пространства}
\zh{函数空间}
Положим
\[
H=L_2(0,l),
\qquad
V=H_0^1(0,l),
\qquad
V^*=H^{-1}(0,l).
\]
Тогда имеем гильбертов тройник
\[
V\subset H\simeq H^*\subset V^*.
\]
Оператор теплопроводности задается слабой формой
\[
a(y,v)=\int_0^l y_s(s)v_s(s)\,ds,
\qquad y,v\in H_0^1(0,l).
\]
Соответствующий оператор
\[
A:V\to V^*
\]
определяется равенством
\[
\langle Ay,v\rangle_{V^*,V}=\int_0^l y_s(s)v_s(s)\,ds.
\]
Тогда уравнение теплопроводности можно записать абстрактно:
\[
\dot x(t)+Ax(t)=u(t),
\qquad x(0)=\varphi.
\]
Здесь \(u(t)\in H=L_2(0,l)\), а также \(u(t)\) можно рассматривать как элемент \(V^*\).
\zh{取 \(H=L_2(0,l)\)、\(V=H_0^1(0,l)\)、\(V^*=H^{-1}(0,l)\)，得到 Hilbert 三元组 \(V\subset H\simeq H^*\subset V^*\)。热传导算子通过弱形式 \(a(y,v)=\int_0^l y_s v_s\,ds\) 定义，即 \(\langle Ay,v\rangle=\int_0^l y_s v_s\,ds\)。于是方程可抽象写为 \(\dot x(t)+Ax(t)=u(t)\)，\(x(0)=\varphi\)。这里 \(u(t)\in H\)，也可自然看作 \(V^*\) 中的元素。}

\subsection*{Обобщенное решение}
\zh{广义解}
Пусть
\[
u\in L_2(Q)=L_2(0,T;L_2(0,l)),
\qquad
\varphi\in L_2(0,l).
\]
Функция \(x=x(s,t)\) называется обобщенным решением краевой задачи, если
\[
x\in L_2(0,T;H_0^1(0,l)),
\]
\[
x_t\in L_2(0,T;H^{-1}(0,l)),
\]
\[
x(0)=\varphi\quad \text{в }L_2(0,l),
\]
и для почти всех \(t\in(0,T)\) выполнено тождество
\[
\langle x_t(t),v\rangle_{H^{-1},H_0^1}
+
\int_0^l x_s(s,t)v_s(s)\,ds
=
\int_0^l u(s,t)v(s)\,ds
\]
для всех
\[
v\in H_0^1(0,l).
\]
Это слабая форма уравнения \(x_t=x_{ss}+u\).
\zh{若 \(u\in L_2(Q)\)、\(\varphi\in L_2(0,l)\)，函数 \(x\) 称为该边值问题的广义解，若 \(x\in L_2(0,T;H_0^1(0,l))\)，\(x_t\in L_2(0,T;H^{-1}(0,l))\)，初始条件在 \(L_2\) 中成立，并且对几乎所有 \(t\) 和任意 \(v\in H_0^1(0,l)\) 满足弱恒等式。这正是方程 \(x_t=x_{ss}+u\) 的弱形式。}

\qpart{2}{Существование и единственность обобщенного решения краевой задачи}{边值问题广义解的存在唯一性}

В интегральной форме, близкой к записи студенческого варианта, это можно записать так: для всех достаточно гладких пробных функций \(\Phi(s,t)\), обращающихся в нуль при \(s=0,l\),
\[
\int_0^l x(s,T)\Phi(s,T)\,ds
+
\iint_Q\left(-x(s,t)\Phi_t(s,t)+x_s(s,t)\Phi_s(s,t)\right)dsdt
\]
\[
=
\int_0^l \varphi(s)\Phi(s,0)\,ds
+
\iint_Q u(s,t)\Phi(s,t)\,dsdt.
\]
Главное преимущество такой формулировки состоит в том, что она не требует существования классических производных \(x_t\) и \(x_{ss}\) как обычных функций.
\zh{也可用接近学生版本的积分形式来写：对所有足够光滑且在 \(s=0,l\) 处为零的试验函数 \(\Phi(s,t)\)，满足相应的积分恒等式。这种表述的主要优点是，它不要求 \(x_t\) 和 \(x_{ss}\) 作为普通函数存在。}

Из стандартной теоремы о функциях из пространства
\[
L_2(0,T;H_0^1(0,l)),
\qquad
x_t\in L_2(0,T;H^{-1}(0,l))
\]
следует вложение
\[
x\in C([0,T];L_2(0,l)).
\]
Поэтому значения \(x(0)\) и \(x(T)\) имеют корректный смысл в \(L_2(0,l)\). Именно это делает функционал
\[
J(u)=\|x(T;u)-b\|_{L_2(0,l)}^2
\]
строго определенным.
\zh{标准定理给出：若 \(x\in L_2(0,T;H_0^1)\) 且 \(x_t\in L_2(0,T;H^{-1})\)，则 \(x\in C([0,T];L_2)\)。因此 \(x(0)\) 和 \(x(T)\) 在 \(L_2(0,l)\) 中有严格意义，目标泛函 \(J(u)=\|x(T;u)-b\|_{L_2}^2\) 也就定义良好。}

\subsection*{Существование обобщенного решения}
\zh{广义解的存在性}
Для каждого
\[
u\in L_2(Q),
\qquad
\varphi\in L_2(0,l)
\]
существует обобщенное решение
\[
x\in L_2(0,T;H_0^1(0,l)),
\qquad
x_t\in L_2(0,T;H^{-1}(0,l)).
\]
Это следует из стандартной теории линейных параболических уравнений. Билинейная форма
\[
a(y,v)=\int_0^l y_s v_s\,ds
\]
непрерывна и коэрцитивна на \(H_0^1(0,l)\). Правая часть \(u\in L_2(0,T;L_2(0,l))\) принадлежит также \(L_2(0,T;H^{-1}(0,l))\). Поэтому задача
\[
\dot x+Ax=u,
\qquad x(0)=\varphi
\]
имеет решение в указанном энергетическом классе.
\zh{对每个 \(u\in L_2(Q)\) 和 \(\varphi\in L_2(0,l)\)，存在满足上述能量类条件的广义解。这来自线性抛物型方程的标准理论。双线性型 \(a(y,v)=\int_0^l y_s v_s\,ds\) 在 \(H_0^1(0,l)\) 上连续且强制，右端 \(u\in L_2(0,T;L_2)\) 也属于 \(L_2(0,T;H^{-1})\)，因此抽象问题 \(\dot x+Ax=u\)、\(x(0)=\varphi\) 有解。}

\subsection*{Единственность обобщенного решения}
\zh{广义解的唯一性}
Пусть \(x_1\) и \(x_2\) -- два обобщенных решения при одном и том же управлении \(u\) и одном и том же начальном условии \(\varphi\). Рассмотрим их разность
\[
z=x_1-x_2.
\]
Тогда \(z\) удовлетворяет однородной задаче
\[
z_t=z_{ss},
\qquad z(0,t)=z(l,t)=0,
\qquad z(s,0)=0.
\]
В слабой форме
\[
\langle z_t(t),v\rangle+
\int_0^l z_s(s,t)v_s(s)\,ds=0.
\]
Берем \(v=z(t)\). Тогда получаем энергетическое равенство
\[
\frac12\frac{d}{dt}\|z(t)\|_{L_2(0,l)}^2+
\|z_s(t)\|_{L_2(0,l)}^2=0.
\]
Интегрируя по времени от \(0\) до \(t\), получаем
\[
\frac12\|z(t)\|_{L_2(0,l)}^2+
\int_0^t\|z_s(\tau)\|_{L_2(0,l)}^2\,d\tau=0.
\]
Следовательно,
\[
z(t)=0
\]
для всех \(t\in[0,T]\) в \(L_2(0,l)\). Значит,
\[
x_1=x_2.
\]
Итак, обобщенное решение краевой задачи существует и единственно.
\zh{设 \(x_1,x_2\) 是同一控制和同一初始条件下的两个广义解，令 \(z=x_1-x_2\)。则 \(z\) 满足齐次问题，初值和边界值均为零。在弱形式中取 \(v=z(t)\)，得到能量等式。对时间积分后，\(\|z(t)\|_{L_2}^2\) 和 \(\int_0^t\|z_s(\tau)\|_{L_2}^2d\tau\) 的和为零，所以 \(z(t)=0\)，从而 \(x_1=x_2\)。因此广义解存在且唯一。}

\subsection*{Теорема Вейерштрасса и существование оптимального управления}
\zh{Weierstrass 定理与最优控制存在性}
Обозначим через
\[
S:u\mapsto x(T;u)
\]
оператор, который каждому управлению \(u\in L_2(Q)\) ставит в соответствие конечное состояние \(x(T;u)\in L_2(0,l)\). Поскольку уравнение теплопроводности линейно,
\[
x(T;u)=x^0(T)+S_Tu,
\]
где \(x^0\) -- решение при \(u=0\), а \(S_T\) -- линейный непрерывный оператор. Функционал можно записать так:
\[
J(u)=\|x^0(T)+S_Tu-b\|_{L_2(0,l)}^2.
\]
Следовательно, \(J\) является выпуклым и слабо полунепрерывным снизу на \(L_2(Q)\).
\zh{记 \(S:u\mapsto x(T;u)\) 为把控制映到终端状态的算子。由于热方程线性，有 \(x(T;u)=x^0(T)+S_Tu\)，其中 \(x^0\) 是 \(u=0\) 时的解，\(S_T\) 是线性连续算子。因此目标泛函可写成 \(J(u)=\|x^0(T)+S_Tu-b\|_{L_2}^2\)，它在 \(L_2(Q)\) 上是凸的并且弱下半连续。}

\qpart{3}{Теорема Вейерштрасса}{Weierstrass 定理}

Множество допустимых управлений
\[
U=\{u\in L_2(Q):\|u\|_{L_2(Q)}\le R\}
\]
является непустым, выпуклым, замкнутым и ограниченным в гильбертовом пространстве \(L_2(Q)\). Так как \(L_2(Q)\) рефлексивно, множество \(U\) слабо компактно.
\zh{允许控制集 \(U=\{u\in L_2(Q):\|u\|_{L_2(Q)}\le R\}\) 在 Hilbert 空间 \(L_2(Q)\) 中非空、凸、闭且有界。由于 \(L_2(Q)\) 自反，\(U\) 弱紧。}

Пусть
\[
\{u_k\}\subset U
\]
-- минимизирующая последовательность:
\[
J(u_k)\to \inf_{u\in U}J(u).
\]
Из слабой компактности \(U\) можно выбрать подпоследовательность, которую снова обозначим \(\{u_k\}\), такую что
\[
u_k\rightharpoonup u_*
\quad \text{в }L_2(Q),
\]
причем
\[
u_*\in U.
\]
Так как \(J\) слабо полунепрерывен снизу,
\[
J(u_*)\le \liminf_{k\to\infty}J(u_k).
\]
Но \(\{u_k\}\) -- минимизирующая последовательность, поэтому
\[
\lim_{k\to\infty}J(u_k)=\inf_{u\in U}J(u).
\]
Следовательно,
\[
J(u_*)=\inf_{u\in U}J(u).
\]
Значит, оптимальное управление существует.
\zh{取极小化序列 \(\{u_k\}\subset U\)。由弱紧性可选出弱收敛子列 \(u_k\rightharpoonup u_*\)，且 \(u_*\in U\)。由于 \(J\) 弱下半连续，有 \(J(u_*)\le\liminf J(u_k)\)。而 \(\{u_k\}\) 是极小化序列，因此 \(J(u_*)=\inf_{u\in U}J(u)\)，最优控制存在。}

Это и есть применение теоремы Вейерштрасса в слабой форме: слабо полунепрерывный снизу функционал достигает минимума на слабо компактном множестве.

Важно: из этого рассуждения следует существование оптимального управления, но не обязательно его единственность. Единственность оптимального управления требует дополнительных условий, например строгой выпуклости функционала по \(u\).
\zh{这正是弱形式下 Weierstrass 定理的应用：弱下半连续泛函在弱紧集合上达到最小值。需要注意，这个论证给出最优控制的存在性，但不必然给出唯一性；唯一性需要额外条件，例如泛函关于 \(u\) 严格凸。}

\subsection*{Краткий итог}
\zh{简要总结}
Задача оптимального нагрева стержня сводится к минимизации функционала
\[
J(u)=\|x(T;u)-b\|_{L_2(0,l)}^2
\]
по управлению \(u\in U\subset L_2(Q)\). Для каждого \(u\in L_2(Q)\) краевая задача для уравнения теплопроводности имеет единственное обобщенное решение
\[
x\in L_2(0,T;H_0^1(0,l)),
\qquad
x_t\in L_2(0,T;H^{-1}(0,l)).
\]
Благодаря вложению
\[
x\in C([0,T];L_2(0,l))
\]
конечное состояние \(x(T)\) корректно определено. Существование оптимального управления следует из слабой компактности множества \(U\) и слабой полунепрерывности снизу функционала \(J\).
\zh{最优加热杆问题归结为在 \(U\subset L_2(Q)\) 上最小化 \(J(u)=\|x(T;u)-b\|_{L_2(0,l)}^2\)。对每个 \(u\in L_2(Q)\)，热方程边值问题有唯一广义解，且 \(x\in C([0,T];L_2)\)，所以终端状态 \(x(T)\) 有意义。最优控制存在性来自允许集 \(U\) 的弱紧性和泛函 \(J\) 的弱下半连续性。}

\newpage

\ticket{17}{Градиент целевого функционала в задаче о нагреве стержня. Критерий оптимальности. Регуляризованный метод проекции градиента.}{加热杆问题中目标泛函的梯度。最优性判据。正则化投影梯度法。}
\qpart{1}{Градиент целевого функционала в задаче о нагреве стержня}{加热杆问题中目标泛函的梯度}

Рассмотрим задачу оптимального нагрева стержня:
\[
J(u)=\int_0^l (x(s,T;u)-b(s))^2\,ds\to\inf_{u\in U},
\]
где состояние \(x=x(s,t;u)\) удовлетворяет уравнению теплопроводности
\[
x_t=x_{ss}+u(s,t),
\qquad (s,t)\in Q=(0,T)\times(0,l),
\]
с краевыми условиями
\[
x(0,t)=x(l,t)=0,
\]
и начальным условием
\[
x(s,0)=\varphi(s).
\]
Множество допустимых управлений:
\[
U=\{u\in L_2(Q):\|u\|_{L_2(Q)}\le R\}.
\]
\zh{考虑最优加热杆问题：在允许控制集 \(U\) 上最小化终端偏差 \(J(u)=\int_0^l(x(s,T;u)-b(s))^2ds\)。状态 \(x\) 满足热传导方程 \(x_t=x_{ss}+u\)、零边界条件和初始条件 \(x(s,0)=\varphi(s)\)，允许控制集是 \(L_2(Q)\) 中半径为 \(R\) 的闭球。}

\subsection*{Градиент целевого функционала}
\zh{目标泛函的梯度}
Пусть \(u\in U\), а \(x=x(s,t;u)\) -- соответствующее решение прямой задачи. Возьмем приращение управления
\[
u+\varepsilon v,
\qquad v\in L_2(Q).
\]
\zh{令 \(u\in U\)，\(x=x(s,t;u)\) 为对应的正问题解。对控制作扰动 \(u+\varepsilon v\)，则状态的一阶变分 \(z=\frac{d}{d\varepsilon}x(u+\varepsilon v)|_{\varepsilon=0}\) 满足线性化方程 \(z_t=z_{ss}+v\)，并满足零边界条件和零初始条件。}
Тогда вариация состояния
\[
z=\frac{d}{d\varepsilon}x(u+\varepsilon v)\bigg|_{\varepsilon=0}
\]
\zh{目标泛函的一阶变分为 \(\delta J(u;v)=2\int_0^l(x(s,T;u)-b(s))z(s,T)\,ds\)。为了把它写成关于 \(v\) 的积分，引入伴随函数 \(\psi\)，它满足反向热方程、零边界条件和终端条件 \(\psi(s,T)=2(x(s,T;u)-b(s))\)。}
удовлетворяет линеаризованной задаче
\[
z_t=z_{ss}+v,
\]
\[
z(0,t)=z(l,t)=0,
\]
\[
z(s,0)=0.
\]
Первая вариация функционала равна
\[
\delta J(u;v)=
2\int_0^l (x(s,T;u)-b(s))z(s,T)\,ds.
\]

Чтобы выразить ее через \(v\), вводят сопряженную функцию \(\psi=\psi(s,t)\), решающую задачу
\[
-\psi_t=\psi_{ss},
\qquad (s,t)\in Q,
\]
\[
\psi(0,t)=\psi(l,t)=0,
\]
\[
\psi(s,T)=2(x(s,T;u)-b(s)).
\]
Эквивалентно первое уравнение можно записать как
\[
-\psi_t-\psi_{ss}=0.
\]
Умножая уравнение для \(z\) на \(\psi\), интегрируя по \(Q\) и используя интегрирование по частям, получаем
\[
\delta J(u;v)=
\iint_Q \psi(s,t)v(s,t)\,dsdt.
\]
Следовательно, производная Фреше функционала имеет вид
\[
\langle J'(u),v\rangle_{L_2(Q)}=
\iint_Q \psi(s,t)v(s,t)\,dsdt.
\]
Так как пространство управлений есть \(L_2(Q)\), то градиент функционала равен
\[
J'(u)=\psi.
\]
\zh{将 \(z\) 的方程乘以 \(\psi\)，在 \(Q\) 上积分并分部积分，就得到 \(\delta J(u;v)=\iint_Q\psi(s,t)v(s,t)\,dsdt\)。因此 Frechet 导数满足 \(\langle J'(u),v\rangle_{L_2(Q)}=\iint_Q\psi v\)，而控制空间是 \(L_2(Q)\)，所以梯度就是 \(J'(u)=\psi\)。}

Итак, для вычисления градиента нужно решить прямую задачу для \(x(s,t;u)\), задать терминальное условие
\[
\psi(s,T)=2(x(s,T;u)-b(s)),
\]
решить сопряженную задачу назад по времени и получить \(J'(u)=\psi\).
\zh{因此，计算梯度的步骤是：先对给定控制 \(u\) 解正问题得到 \(x\)，再用终端条件 \(\psi(s,T)=2(x(s,T;u)-b(s))\) 解伴随问题，最后得到 \(J'(u)=\psi\)。}

\subsection*{Критерий оптимальности}
\zh{最优性判据}
Пусть
\[
U=\{u\in L_2(Q):\|u\|\le R\}
\]
-- замкнутый выпуклый шар. Если \(u_*\in U\) -- оптимальное управление, то для выпуклого дифференцируемого функционала необходимое и достаточное условие оптимальности имеет вид вариационного неравенства:
\[
\langle J'(u_*),u-u_*\rangle_{L_2(Q)}\ge0
\qquad \forall u\in U.
\]
Так как \(J'(u_*)=\psi_*\), это можно записать как
\[
\iint_Q \psi_*(s,t)(u(s,t)-u_*(s,t))\,dsdt\ge0
\qquad \forall u\in U.
\]
Эквивалентная проекционная форма:
\[
u_*=P_U\bigl(u_* -\beta J'(u_*)\bigr)
\qquad \forall \beta>0.
\]
\zh{设 \(U=\{u\in L_2(Q):\|u\|\le R\}\) 是闭凸球。若 \(u_*\in U\) 为最优控制，且 \(J\) 凸且可微，则最优性的充要条件是变分不等式 \(\langle J'(u_*),u-u_*\rangle_{L_2(Q)}\ge0\) 对所有 \(u\in U\) 成立。因为 \(J'(u_*)=\psi_*\)，这也可写成 \(\iint_Q\psi_*(u-u_*)\,dsdt\ge0\)。等价的投影形式为 \(u_*=P_U(u_*-\beta J'(u_*))\)。}

\qpart{2}{Критерий оптимальности}{最优性判据}

Для шара \(U\) условие можно записать в форме Куна--Таккера. Рассмотрим ограничение
\[
\|u\|^2-R^2\le0
\]
и функцию Лагранжа
\[
L(u,\lambda)=J(u)+\lambda(\|u\|^2-R^2),
\qquad \lambda\ge0.
\]
Тогда оптимальность \(u_*\) эквивалентна существованию \(\lambda_*\ge0\), такого что
\[
J'(u_*)+2\lambda_*u_*=0,
\]
\[
\lambda_*(\|u_*\|^2-R^2)=0,
\]
\[
\|u_*\|\le R.
\]
Если ограничение не активно, то
\[
\|u_*\|<R,
\qquad
\lambda_*=0,
\qquad
J'(u_*)=0.
\]
Если ограничение активно, то \(\|u_*\|=R\). Из условия \(J'(u_*)+2\lambda_*u_*=0\) следует
\[
\lambda_*=\frac{\|J'(u_*)\|}{2R}.
\]
Если используется запись
\[
J'(u_*)+\mu_*u_*=0,
\]
то \(\mu_*=2\lambda_*\), и в активном случае
\[
\mu_*=\frac{\|J'(u_*)\|}{R}.
\]
Именно это различие объясняет разные коэффициенты в студенческих версиях.
\zh{对球形约束，也可写成 Kuhn--Tucker 条件。对约束 \(\|u\|^2-R^2\le0\)，构造 Lagrange 函数 \(L(u,\lambda)=J(u)+\lambda(\|u\|^2-R^2)\)。最优性等价于存在 \(\lambda_*\ge0\)，使 \(J'(u_*)+2\lambda_*u_*=0\)、互补松弛条件 \(\lambda_*(\|u_*\|^2-R^2)=0\) 和可行性 \(\|u_*\|\le R\) 成立。若约束不活跃，则 \(J'(u_*)=0\)；若约束活跃，则 \(\|u_*\|=R\)，并由记号差异产生学生版本中不同的系数写法。}

Критерий оптимальности можно также записать в седловой форме. Точка \((u_*,\lambda_*)\) является седловой точкой функции
\[
L(u,\lambda)=J(u)+\lambda(\|u\|^2-R^2),
\qquad \lambda\ge0,
\]
если
\[
L(u_*,\lambda)\le L(u_*,\lambda_*)\le L(u,\lambda_*)
\]
для всех \(u\in L_2(Q)\), \(\lambda\ge0\). В этом случае \(u_*\) является оптимальным управлением.
\zh{最优性还可写成鞍点形式：若 \((u_*,\lambda_*)\) 是 Lagrange 函数的鞍点，即对所有 \(u\in L_2(Q)\)、\(\lambda\ge0\) 有 \(L(u_*,\lambda)\le L(u_*,\lambda_*)\le L(u,\lambda_*)\)，则 \(u_*\) 是最优控制。}

\subsection*{Регуляризованный метод проекции градиента}
\zh{正则化投影梯度法}
Обычный метод проекции градиента для задачи
\[
J(u)\to\inf,
\qquad u\in U
\]
имеет вид
\[
u_{k+1}=P_U\left(u_k-\beta_kJ'(u_k)\right),
\qquad k=0,1,\ldots .
\]
Здесь \(P_U\) -- метрическая проекция на замкнутое выпуклое множество \(U\), а \(\beta_k>0\) -- шаг.
\zh{普通投影梯度法用于在闭凸集 \(U\) 上最小化 \(J(u)\)，迭代为 \(u_{k+1}=P_U(u_k-\beta_kJ'(u_k))\)。这里 \(P_U\) 是到闭凸集 \(U\) 的度量投影，\(\beta_k>0\) 是步长。}

\qpart{3}{Регуляризованный метод проекции градиента}{正则化投影梯度法}

В регуляризованном методе вводится функция Тихонова
\[
T_k(u)=J(u)+\alpha_k\|u\|^2,
\qquad \alpha_k>0.
\]
Ее градиент:
\[
T_k'(u)=J'(u)+2\alpha_k u.
\]
Если градиент вычисляется приближенно, используют \(J_k'(u)\) и предполагают оценку ошибки вида
\[
\|J_k'(u)-J'(u)\|\le \delta_k(1+\|u\|).
\]
Тогда регуляризованный метод проекции градиента задается формулой
\[
u_{k+1}=P_U\left(u_k-\beta_k\bigl(J_k'(u_k)+2\alpha_k u_k\bigr)\right).
\]
Если градиент точный, то \(J_k'(u_k)=J'(u_k)\), и формула принимает вид
\[
u_{k+1}=P_U\left(u_k-\beta_k\bigl(J'(u_k)+2\alpha_k u_k\bigr)\right).
\]
В задаче нагрева стержня это означает
\[
u_{k+1}=P_U\left(u_k-\beta_k\bigl(\psi_k+2\alpha_k u_k\bigr)\right),
\]
где \(\psi_k\) -- решение сопряженной задачи, соответствующее управлению \(u_k\).
\zh{正则化方法引入 Tikhonov 函数 \(T_k(u)=J(u)+\alpha_k\|u\|^2\)，其梯度为 \(T_k'(u)=J'(u)+2\alpha_k u\)。若梯度近似计算，则使用 \(J_k'(u)\)，并要求误差与正则化参数相协调。迭代公式为 \(u_{k+1}=P_U(u_k-\beta_k(J_k'(u_k)+2\alpha_k u_k))\)。在加热杆问题中，\(J_k'(u_k)\) 对应伴随问题解 \(\psi_k\)。}

\subsection*{Проекция на шар}
\zh{到球的投影}
Если
\[
U=\{u\in L_2(Q):\|u\|\le R\},
\]
то проекция имеет явный вид:
\[
P_U(w)=
\begin{cases}
w, & \|w\|\le R,\\[2mm]
\dfrac{R}{\|w\|}w, & \|w\|>R.
\end{cases}
\]
Поэтому один шаг метода можно записать так:
\[
w_k=u_k-\beta_k(\psi_k+2\alpha_k u_k),
\]
\[
u_{k+1}=P_U(w_k).
\]
\zh{若 \(U=\{u\in L_2(Q):\|u\|\le R\}\)，投影有显式形式：当 \(\|w\|\le R\) 时 \(P_U(w)=w\)；当 \(\|w\|>R\) 时 \(P_U(w)=\frac{R}{\|w\|}w\)。因此一次迭代可先算 \(w_k=u_k-\beta_k(\psi_k+2\alpha_k u_k)\)，再令 \(u_{k+1}=P_U(w_k)\)。}

\subsection*{Условия сходимости}
\zh{收敛条件}
Типичные условия сходимости регуляризованного метода:
\begin{enumerate}[leftmargin=2em]
\item \(U\) -- выпуклое замкнутое множество в гильбертовом пространстве;
\item \(J\) выпукл и непрерывно дифференцируем по Фреше;
\item множество оптимальных решений непусто;
\item градиент \(J'\) удовлетворяет условию типа Липшица или коэрцитивности;
\item параметры регуляризации стремятся к нулю:
\[
\alpha_k\to0;
\]
\item ошибка приближенного градиента согласована с регуляризацией:
\[
\frac{\delta_k}{\alpha_k}\to0;
\]
\item шаги \(\beta_k\) выбраны так, чтобы обеспечивать устойчивость и убывание регуляризованного функционала.
\end{enumerate}
При этих условиях последовательность \(u_k\) сходится к нормальному оптимальному управлению, то есть к оптимальному управлению минимальной нормы:
\[
u_k\to u_*,
\]
где
\[
u_*=
\operatorname*{arg\,min}\{\|u\|:u\in U^*\}.
\]
Здесь \(U^*\) -- множество оптимальных управлений.
\zh{正则化投影梯度法的典型收敛条件包括：\(U\) 是 Hilbert 空间中的闭凸集；\(J\) 凸且 Frechet 连续可微；最优解集合非空；梯度满足 Lipschitz 型条件或强制性条件；正则化参数 \(\alpha_k\to0\)；近似梯度误差与正则化协调，例如 \(\delta_k/\alpha_k\to0\)；步长 \(\beta_k\) 选取得足以保证稳定性和正则化泛函下降。在这些条件下，\(u_k\) 收敛到正规最优控制，即最优解集合中范数最小的元素。}

\subsection*{Краткий итог}
\zh{简要总结}
В задаче нагрева стержня градиент целевого функционала находится через сопряженную задачу:
\[
J'(u)=\psi.
\]
Оптимальность на выпуклом множестве \(U\) задается вариационным неравенством
\[
\langle J'(u_*),u-u_*\rangle\ge0
\qquad \forall u\in U.
\]
Для шара это эквивалентно условиям Куна--Таккера
\[
J'(u_*)+2\lambda_*u_*=0,
\]
\[
\lambda_*(\|u_*\|^2-R^2)=0,
\]
\[
\lambda_*\ge0,
\qquad
\|u_*\|\le R.
\]
Регуляризованный метод проекции градиента имеет вид
\[
u_{k+1}=P_U\left(u_k-\beta_k(J_k'(u_k)+2\alpha_k u_k)\right).
\]
\zh{在加热杆问题中，目标泛函梯度通过伴随问题求得：\(J'(u)=\psi\)。凸集 \(U\) 上的最优性由变分不等式 \(\langle J'(u_*),u-u_*\rangle\ge0\) 描述。若 \(U\) 为球，则等价于 Kuhn--Tucker 条件 \(J'(u_*)+2\lambda_*u_*=0\)、\(\lambda_*(\|u_*\|^2-R^2)=0\)、\(\lambda_*\ge0\)、\(\|u_*\|\le R\)。正则化投影梯度法的迭代为 \(u_{k+1}=P_U(u_k-\beta_k(J_k'(u_k)+2\alpha_k u_k))\)。}

\blocktitle{F}{Линейно-квадратичные дифференциальные игры и максимум по Парето}{Билеты 18--20}
\section*{Замечание об источниках и надежности}
\zh{可靠性说明}
Программа государственного экзамена указывает в качестве основных источников для этого блока книги В.И.~Жуковского и соавторов по линейно-квадратичным дифференциальным играм, гарантированным решениям конфликтов, уравновешиванию конфликтов, а также дополнительные источники по векторной оптимизации и многокритериальным динамическим системам. Полные тексты некоторых книг в текущем наборе материалов отсутствуют. Поэтому ниже используются: формулировки билетов, студенческие версии 2017 и 2026 гг., а также открытые статьи В.И.~Жуковского и соавторов, непосредственно относящиеся к билетам 18--20.
\zh{国家考试大纲把 Жуковский 及其合作者关于线性二次微分博弈、冲突保证解、冲突均衡化的著作，以及向量优化和多准则动力系统方面的补充资料列为本板块主要来源。当前材料中缺少部分著作全文，因此下面主要依据题目表、2017 与 2026 年学生版本，以及与第18--20题直接相关的 Жуковский 等人的公开文章整理。}

Для билета 18 главным источником является статья \emph{«Об одной нерешенной задаче в матричных обыкновенных дифференциальных уравнениях»}. Для билета 19 главным источником является статья \emph{«Berge--Vaisman equilibrium for one linear-quadratic differential game»}. Для билета 20 используются открытые материалы о максимуме по Парето в линейно-квадратичных динамических играх, но отдельные детали желательно сверить с книгами по векторной оптимизации динамических систем.
\zh{第18题的主要来源是关于矩阵常微分方程中一个未解决问题的文章；第19题的主要来源是 Berge--Vaisman equilibrium for one linear-quadratic differential game；第20题使用关于线性二次动态博弈中 Pareto 最大值的公开资料，但个别细节仍建议同动态系统向量优化著作核对。}

\ticket{18}{Условия существования и явный вид ситуации равновесия по Нэшу в дифференциальной линейно-квадратичной дифференциальной игре трех лиц.}{三人线性二次微分博弈中 Nash 均衡情形的存在条件和显式形式。}


Рассмотрим дифференциальную позиционную линейно-квадратичную игру трех лиц
\[
\Gamma=\langle \{1,2,3\},\Sigma,\{\mathcal A_i\}_{i=1}^3,\{J_i\}_{i=1}^3\rangle .
\]
Фазовый вектор и управляющие воздействия имеют размерность
\[
x(t)\in\R^n,\qquad u_i(t)\in\R^n,\quad i=1,2,3.
\]
Динамика задается системой
\[
\dot x(t)=A(t)x(t)+u_1(t)+u_2(t)+u_3(t),\qquad x(t_0)=x_0,
\]
где \(A(\cdot)\in C^{n\times n}[0,\vartheta]\), а \(\vartheta>t_0\) -- фиксированный момент окончания игры.

Игроки используют линейные позиционные стратегии
\[
U_i:\quad u_i(t,x)=Q_i(t)x,\qquad Q_i(\cdot)\in C^{n\times n}[0,\vartheta].
\]
Множество стратегий \(i\)-го игрока:
\[
\mathcal A_i=\{U_i:\ u_i(t,x)=Q_i(t)x,\ Q_i(\cdot)\in C^{n\times n}[0,\vartheta]\}.
\]
Ситуация игры есть тройка
\[
U=(U_1,U_2,U_3)\in \mathcal A_1\times\mathcal A_2\times\mathcal A_3.
\]
При фиксированной ситуации \(U\) движение удовлетворяет
\[
\dot x(t)=\bigl(A(t)+Q_1(t)+Q_2(t)+Q_3(t)\bigr)x(t),\qquad x(t_0)=x_0.
\]

Функция выигрыша \(i\)-го игрока задается квадратичным функционалом
\[
J_i(U,t_0,x_0)=x^{\top}(\vartheta)C_i x(\vartheta)+
\int_{t_0}^{\vartheta}\sum_{j=1}^3 u_j^{\top}(t)D_{ij}u_j(t)\,dt,
\qquad i=1,2,3,
\]
где \(C_i,D_{ij}\) -- постоянные симметричные \(n\times n\)-матрицы. Каждый игрок стремится максимизировать свой выигрыш.
\zh{考虑三人线性二次位置微分博弈。相变量 \(x(t)\in\R^n\)，各玩家控制 \(u_i(t)\in\R^n\)。系统动力学为 \(\dot x=A(t)x+u_1+u_2+u_3\)。玩家采用线性位置策略 \(u_i(t,x)=Q_i(t)x\)。固定策略组后，闭环运动由 \(\dot x=(A+Q_1+Q_2+Q_3)x\) 给出。第 \(i\) 个玩家的收益为二次泛函 \(J_i\)，其中 \(C_i,D_{ij}\) 为对称矩阵，每个玩家都最大化自己的收益。}

\qpart{1}{Условия существования равновесия по Нэшу}{Nash 均衡的存在条件}

\subsection*{Равновесие по Нэшу}
\zh{Nash 均衡}
Ситуация \(U^e=(U_1^e,U_2^e,U_3^e)\) называется равновесием по Нэшу, если для любой начальной позиции
\[
(t_0,x_0)\in [0,\vartheta)\times\R^n,\qquad x_0\ne 0,
\]
выполнены условия
\[
\max_{U_1\in\mathcal A_1}J_1(U_1,U_2^e,U_3^e,t_0,x_0)=J_1(U^e,t_0,x_0),
\]
\[
\max_{U_2\in\mathcal A_2}J_2(U_1^e,U_2,U_3^e,t_0,x_0)=J_2(U^e,t_0,x_0),
\]
\[
\max_{U_3\in\mathcal A_3}J_3(U_1^e,U_2^e,U_3,t_0,x_0)=J_3(U^e,t_0,x_0).
\]
Смысл: ни один игрок не может увеличить свой выигрыш, если остальные два сохраняют свои равновесные стратегии.
\zh{若对任意初始位置 \((t_0,x_0)\) 和每个玩家 \(i\)，在其他两个玩家保持均衡策略时，第 \(i\) 个玩家不能通过单方面改变自己的策略来提高收益，则 \(U^e\) 称为 Nash 均衡。直观地说，Nash 均衡中没有任何一个玩家有单边偏离的动机。}

\qpart{2}{Явный вид ситуации равновесия в трехлицевой LQ-игре}{三人 LQ 博弈中均衡情形的显式形式}

\subsection*{Метод динамического программирования}
\zh{动态规划方法}
Для построения равновесия вводятся функции Беллмана
\[
V_i(t,x),\qquad i=1,2,3.
\]
Определим
\[
W_i(t,x,u_1,u_2,u_3,V_i)=
\frac{\partial V_i}{\partial t}+
\left(\frac{\partial V_i}{\partial x}\right)^{\top}(A(t)x+u_1+u_2+u_3)+
\sum_{j=1}^3u_j^{\top} D_{ij}u_j .
\]
В равновесии по Нэшу \(i\)-й игрок максимизирует \(W_i\) только по собственному управлению \(u_i\), считая управления остальных игроков фиксированными. Для конечного максимума по \(u_i\) требуется вогнутость по \(u_i\), то есть
\[
D_{ii}<0,
\qquad i=1,2,3.
\]
Условие первого порядка дает
\[
\frac{\partial W_i}{\partial u_i}=\frac{\partial V_i}{\partial x}+2D_{ii}u_i=0.
\]
Следовательно,
\[
u_i^e(t,x,V_i)=-\frac12D_{ii}^{-1}\frac{\partial V_i}{\partial x}.
\]
\zh{构造均衡时引入各玩家的 Bellman 函数 \(V_i(t,x)\) 和表达式 \(W_i\)。在 Nash 均衡中，第 \(i\) 个玩家只对自己的控制 \(u_i\) 最大化 \(W_i\)，其余玩家控制视为固定。为了关于 \(u_i\) 的最大值有限，需要 \(D_{ii}<0\)，即关于自身控制的二次项为凹。由一阶条件 \(\frac{\partial W_i}{\partial u_i}=0\) 得到 \(u_i^e=-\frac12D_{ii}^{-1}V_{i,x}\)。}

В линейно-квадратичном случае функции Беллмана ищутся в виде
\[
V_i(t,x)=x^{\top}\Theta_i(t)x,
\qquad \Theta_i^{\top}(t)=\Theta_i(t).
\]
Тогда
\[
\frac{\partial V_i}{\partial x}=2\Theta_i(t)x,
\]
и равновесные стратегии имеют вид
\[
u_i^e(t,x)=-D_{ii}^{-1}\Theta_i(t)x,
\qquad i=1,2,3.
\]
Итак,
\[
Q_i^e(t)=-D_{ii}^{-1}\Theta_i(t).
\]
\zh{在线性二次情形下，Bellman 函数取二次型 \(V_i(t,x)=x^\top\Theta_i(t)x\)，其中 \(\Theta_i\) 对称。于是 \(V_{i,x}=2\Theta_i x\)，均衡策略具有线性位置形式 \(u_i^e(t,x)=-D_{ii}^{-1}\Theta_i(t)x\)，即 \(Q_i^e(t)=-D_{ii}^{-1}\Theta_i(t)\)。}

\subsection*{Система матричных уравнений типа Риккати}
\zh{Riccati 型矩阵方程组}
Подстановка \(V_i=x^{\top}\Theta_i x\) и \(u_i^e=-D_{ii}^{-1}\Theta_i x\) в уравнения динамического программирования приводит к связанной системе матричных уравнений типа Риккати. В симметричной компактной записи для \(i=1,2,3\):
\[
\begin{aligned}
0={}&\dot\Theta_i+\Theta_iA+A^{\top}\Theta_i-\Theta_iD_{ii}^{-1}\Theta_i \\
&-\sum_{\substack{j=1\\ j\ne i}}^3
\left(\Theta_iD_{jj}^{-1}\Theta_j+\Theta_jD_{jj}^{-1}\Theta_i\right)
+\sum_{\substack{j=1\\ j\ne i}}^3
\Theta_jD_{jj}^{-1}D_{ij}D_{jj}^{-1}\Theta_j .
\end{aligned}
\]
Терминальные условия:
\[
\Theta_i(\vartheta)=C_i,\qquad i=1,2,3.
\]

Эта связанная система является главным техническим объектом билета. В общем случае построение ее явного решения для трех игроков является нетривиальной задачей; поэтому корректная экзаменационная формулировка должна звучать условно: если такая система имеет продолжимое решение, то соответствующие стратегии образуют равновесие по Нэшу.
\zh{把二次型 \(V_i=x^\top\Theta_i x\) 和线性均衡控制代入动态规划方程，会得到耦合的 Riccati 型矩阵方程组，并带终端条件 \(\Theta_i(\vartheta)=C_i\)。这是本题的主要技术对象。三人情形中显式求解该耦合系统通常并不简单，因此考试表述应带条件：如果该系统在整个区间上存在可延拓的对称矩阵解，则相应线性位置策略构成 Nash 均衡。}

\subsection*{Достаточное условие существования}
\zh{存在性的充分条件}
Если выполнены условия
\[
D_{ii}<0,
\qquad i=1,2,3,
\]
и система матричных уравнений типа Риккати с терминальными условиями \(\Theta_i(\vartheta)=C_i\) имеет продолжимое на \([0,\vartheta]\) решение в классе симметричных матриц
\[
\Theta_i(t)=\Theta_i^{\top}(t),
\qquad i=1,2,3,
\]
то ситуация
\[
U^e=(U_1^e,U_2^e,U_3^e),
\qquad
U_i^e:\ u_i^e(t,x)=-D_{ii}^{-1}\Theta_i(t)x,
\]
является равновесием по Нэшу. Равновесные выигрыши равны
\[
J_i^e(t_0,x_0)=J_i(U^e,t_0,x_0)=x_0^{\top}\Theta_i(t_0)x_0,
\qquad i=1,2,3.
\]
\zh{若 \(D_{ii}<0\) 且带终端条件的 Riccati 型矩阵方程组在 \([0,\vartheta]\) 上存在对称可延拓解，则由 \(u_i^e(t,x)=-D_{ii}^{-1}\Theta_i(t)x\) 给出的策略组 \(U^e\) 是 Nash 均衡。对应的均衡收益为 \(J_i^e(t_0,x_0)=x_0^\top\Theta_i(t_0)x_0\)。}

\subsection*{Случай отсутствия равновесия}
\zh{不存在均衡的情形}
Если хотя бы для одного игрока
\[
D_{ii}>0,
\]
то максимум по \(u_i\) в определении равновесия по Нэшу не существует: игрок может увеличивать выигрыш за счет неограниченного усиления своей линейной стратегии. Следовательно, в рассматриваемом классе линейных позиционных стратегий равновесие по Нэшу отсутствует.
\zh{若至少存在一个玩家满足 \(D_{ii}>0\)，则关于该玩家控制的最大化不存在：该玩家可以通过无限增大线性策略强度来提高收益。因此在所考虑的线性位置策略类中 Nash 均衡不存在。}

\memory
Модель:
\[
\dot x=A(t)x+u_1+u_2+u_3,
\qquad u_i(t,x)=Q_i(t)x.
\]
Выигрыши:
\[
J_i=x^{\top}(\vartheta)C_ix(\vartheta)+\int_{t_0}^{\vartheta}\sum_{j=1}^3u_j^{\top}D_{ij}u_j\,dt.
\]
Равновесие по Нэшу:
\[
J_i(U^e)\ge J_i(U_i,U_{-i}^e)\quad \forall U_i,\\quad i=1,2,3.
\]
Ищем
\[
V_i(t,x)=x^{\top}\Theta_i(t)x.
\]
Условие максимума по \(u_i\):
\[
V_{i,x}+2D_{ii}u_i=0.
\]
Если \(D_{ii}<0\), то
\[
u_i^e(t,x)=-D_{ii}^{-1}\Theta_i(t)x.
\]
Матрицы \(\Theta_i\) удовлетворяют связанной системе Риккати с условиями
\[
\Theta_i(\vartheta)=C_i.
\]
Если система имеет продолжимое симметричное решение на \([0,\vartheta]\), то найденные стратегии образуют равновесие по Нэшу, а
\[
J_i^e(t_0,x_0)=x_0^{\top}\Theta_i(t_0)x_0.
\]
Если хотя бы одно \(D_{ii}>0\), равновесие по Нэшу отсутствует.
\zh{记忆版：模型为 \(\dot x=A(t)x+u_1+u_2+u_3\)，策略为 \(u_i=Q_i x\)。取 \(V_i=x^\top\Theta_i x\)，若 \(D_{ii}<0\)，一阶条件给出 \(u_i^e=-D_{ii}^{-1}\Theta_i x\)。矩阵 \(\Theta_i\) 满足耦合 Riccati 系统和终端条件 \(\Theta_i(\vartheta)=C_i\)。若该系统有可延拓对称解，则得到 Nash 均衡；若某个 \(D_{ii}>0\)，Nash 均衡不存在。}

\newpage
\ticket{19}{Условия существования и явный вид ситуации равновесия по Бержу в дифференциальной линейно-квадратичной игре двух лиц с малым влиянием одного из игроков на скорость изменения фазового вектора.}{两人线性二次微分博弈中 Berge 均衡情形的存在条件和显式形式，其中一个玩家对相向量变化速度的影响较小。}


Рассмотрим дифференциальную позиционную линейно-квадратичную игру двух лиц
\[
\Gamma=\langle \{1,2\},\Sigma,\{\mathcal A_i\}_{i=1}^2,\{J_i\}_{i=1}^2\rangle .
\]
Динамика фазового вектора \(x(t)\in\R^n\) задается системой
\[
\dot x(t)=A(t)x(t)+u_1(t)+\varepsilon u_2(t),
\qquad x(t_0)=x_0,
\]
где \(A(\cdot)\in C^{n\times n}[0,\vartheta]\), а \(\varepsilon\ge 0\) -- малый параметр. Содержательно \(\varepsilon u_2\) означает малое влияние второго игрока на скорость изменения фазового вектора.
\zh{考虑两人线性二次位置微分博弈。相变量满足 \(\dot x=A(t)x+u_1+\varepsilon u_2\)，其中 \(\varepsilon\ge0\) 是小参数。项 \(\varepsilon u_2\) 表示第二个玩家对相向量变化速度的影响较小。}

Игроки используют линейные позиционные стратегии
\[
U_i:\quad u_i(t,x)=Q_i(t)x,
\qquad Q_i(\cdot)\in C^{n\times n}[0,\vartheta],
\quad i=1,2.
\]
При фиксированной ситуации \(U=(U_1,U_2)\) движение имеет вид
\[
\dot x(t)=\bigl(A(t)+Q_1(t)+\varepsilon Q_2(t)\bigr)x(t).
\]
\zh{玩家采用线性位置策略 \(u_i(t,x)=Q_i(t)x\)。固定策略组 \(U=(U_1,U_2)\) 后，闭环运动为 \(\dot x=(A+Q_1+\varepsilon Q_2)x\)。}

Функции выигрыша:
\[
J_1(U,t_0,x_0)=x^{\top}(\vartheta)C_1x(\vartheta)+
\int_{t_0}^{\vartheta}\left(u_1^{\top}D_{11}u_1-u_2^{\top}D_{12}u_2+x^{\top}G_1x\right)dt,
\]
\[
J_2(U,t_0,x_0)=x^{\top}(\vartheta)C_2x(\vartheta)+
\int_{t_0}^{\vartheta}\left(-u_1^{\top}D_{21}u_1+u_2^{\top}D_{22}u_2+x^{\top}G_2x\right)dt.
\]
Матрицы \(C_i,G_i,D_{ij}\) симметричны. В данной записи отрицательные знаки перед чужими управлениями вынесены явно, поэтому условия вогнутости формулируются как \(D_{12}>0\), \(D_{21}>0\).
\zh{两个玩家的收益函数为线性二次型，矩阵 \(C_i,G_i,D_{ij}\) 均为对称矩阵。这里把“对方控制”前的负号显式写出，所以凹性条件写成 \(D_{12}>0\)、\(D_{21}>0\)。}

\qpart{1}{Условия существования равновесия по Бержу}{Berge 均衡的存在条件}

\subsection*{Равновесие по Бержу}
\zh{Berge 均衡}
Равновесие по Бержу отличается от равновесия по Нэшу. В Нэше каждый игрок максимизирует свой выигрыш по собственной стратегии. В равновесии по Бержу выигрыш каждого игрока максимизируется за счет стратегии другого игрока.
\zh{Berge 均衡不同于 Nash 均衡。在 Nash 均衡中，每个玩家通过选择自己的策略来最大化自己的收益；在 Berge 均衡中，每个玩家的收益是由另一个玩家的策略来最大化的。}

Ситуация \(U^B=(U_1^B,U_2^B)\) называется равновесием по Бержу, если
\[
J_1(U_1^B,U_2,t_0,x_0)\le J_1(U^B,t_0,x_0)
\qquad \forall U_2\in\mathcal A_2,
\]
\[
J_2(U_1,U_2^B,t_0,x_0)\le J_2(U^B,t_0,x_0)
\qquad \forall U_1\in\mathcal A_1.
\]
В строгой версии источника используется равновесие по Бержу--Вайсману, где дополнительно учитывается индивидуальная рациональность.
\zh{若策略组 \(U^B=(U_1^B,U_2^B)\) 满足：在 \(U_1^B\) 固定时，\(U_2^B\) 使第1个玩家收益最大；在 \(U_2^B\) 固定时，\(U_1^B\) 使第2个玩家收益最大，则称为 Berge 均衡。严格来源中使用 Berge--Vaisman 均衡，还额外考虑个体理性条件。}

\qpart{2}{Явный вид ситуации равновесия по Бержу}{Berge 均衡情形的显式形式}

\subsection*{Динамическое программирование}
\zh{动态规划}
Введем функции Беллмана \(V_1(t,x)\), \(V_2(t,x)\) и выражения
\[
\begin{aligned}
W_1={}&\frac{\partial V_1}{\partial t}+\left(\frac{\partial V_1}{\partial x}\right)^{\top}(Ax+u_1+\varepsilon u_2)
+u_1^{\top}D_{11}u_1-u_2^{\top}D_{12}u_2+x^{\top}G_1x,\\
W_2={}&\frac{\partial V_2}{\partial t}+\left(\frac{\partial V_2}{\partial x}\right)^{\top}(Ax+u_1+\varepsilon u_2)
-u_1^{\top}D_{21}u_1+u_2^{\top}D_{22}u_2+x^{\top}G_2x.
\end{aligned}
\]
Для равновесия по Бержу условия максимума имеют перекрестный вид:
\[
\max_{u_2}W_1(t,x,u_1^B,u_2,V_1)=W_1(t,x,u_1^B,u_2^B,V_1),
\]
\[
\max_{u_1}W_2(t,x,u_1,u_2^B,V_2)=W_2(t,x,u_1^B,u_2^B,V_2).
\]
Условия первого порядка:
\[
\frac{\partial W_1}{\partial u_2}=\varepsilon\frac{\partial V_1}{\partial x}-2D_{12}u_2=0,
\]
\[
\frac{\partial W_2}{\partial u_1}=\frac{\partial V_2}{\partial x}-2D_{21}u_1=0.
\]
Отсюда
\[
u_2^B(t,x,V_1)=\frac{\varepsilon}{2}D_{12}^{-1}\frac{\partial V_1}{\partial x},
\qquad
u_1^B(t,x,V_2)=\frac12D_{21}^{-1}\frac{\partial V_2}{\partial x}.
\]
\zh{引入 Bellman 函数 \(V_1,V_2\) 和表达式 \(W_1,W_2\)。对于 Berge 均衡，最大化条件呈交叉形式：\(J_1\) 对 \(u_2\) 最大化，\(J_2\) 对 \(u_1\) 最大化。一阶条件给出 \(\partial W_1/\partial u_2=0\)、\(\partial W_2/\partial u_1=0\)，由此得到 \(u_2^B=\frac{\varepsilon}{2}D_{12}^{-1}V_{1,x}\)、\(u_1^B=\frac12D_{21}^{-1}V_{2,x}\)。}

Ищем функции Беллмана в виде
\[
V_i(t,x)=x^{\top}\Theta_i(t,\varepsilon)x,
\qquad \Theta_i^{\top}(t,\varepsilon)=\Theta_i(t,\varepsilon).
\]
Тогда
\[
\frac{\partial V_i}{\partial x}=2\Theta_i(t,\varepsilon)x,
\]
и равновесные стратегии имеют вид
\[
u_1^B(t,x)=D_{21}^{-1}\Theta_2(t,\varepsilon)x,
\]
\[
u_2^B(t,x)=\varepsilon D_{12}^{-1}\Theta_1(t,\varepsilon)x.
\]
Следовательно,
\[
Q_1^B(t,\varepsilon)=D_{21}^{-1}\Theta_2(t,\varepsilon),
\qquad
Q_2^B(t,\varepsilon)=\varepsilon D_{12}^{-1}\Theta_1(t,\varepsilon).
\]
\zh{取二次型 Bellman 函数 \(V_i(t,x)=x^\top\Theta_i(t,\varepsilon)x\)。由于 \(V_{i,x}=2\Theta_i x\)，均衡策略变为 \(u_1^B(t,x)=D_{21}^{-1}\Theta_2(t,\varepsilon)x\)，\(u_2^B(t,x)=\varepsilon D_{12}^{-1}\Theta_1(t,\varepsilon)x\)。因此 \(Q_1^B=D_{21}^{-1}\Theta_2\)，\(Q_2^B=\varepsilon D_{12}^{-1}\Theta_1\)。}

\qpart{3}{Малое влияние одного игрока на скорость изменения фазового вектора}{一个玩家对相向量速度的小影响}

\subsection*{Матричная система и малый параметр}
\zh{矩阵系统与小参数}
Подстановка найденных управлений в уравнения динамического программирования приводит к системе матричных уравнений для \(\Theta_1\), \(\Theta_2\). В форме, удобной для экзаменационного изложения, ее можно записать так:
\[
\begin{aligned}
0={}&\dot\Theta_1+
\Theta_1(A+D_{21}^{-1}\Theta_2)+(A^{\top}+\Theta_2D_{21}^{-1})\Theta_1+G_1 \\
&+\Theta_2D_{21}^{-1}D_{11}D_{21}^{-1}\Theta_2
+\varepsilon^2\Psi_1(\Theta_1,\Theta_2),
\qquad \Theta_1(\vartheta,\varepsilon)=C_1,
\end{aligned}
\]
\[
\begin{aligned}
0={}&\dot\Theta_2+
\Theta_2A+A^{\top}\Theta_2+G_2+\Theta_2D_{21}^{-1}\Theta_2
+\varepsilon^2\Psi_2(\Theta_1,\Theta_2),\\
&\hspace{8cm}\Theta_2(\vartheta,\varepsilon)=C_2.
\end{aligned}
\]
Здесь \(\Psi_1,\Psi_2\) обозначают слагаемые, квадратичные по матрицам \(\Theta_1,\Theta_2\), возникающие из членов с \(\varepsilon u_2\). Главное: при \(\varepsilon=0\) система упрощается, а при достаточно малом \(\varepsilon\) решение продолжается по теореме о зависимости решений от параметра.
\zh{把得到的控制代入动态规划方程，可得关于 \(\Theta_1,\Theta_2\) 的矩阵方程组。其中 \(\Psi_1,\Psi_2\) 表示由 \(\varepsilon u_2\) 项产生的、关于 \(\Theta_1,\Theta_2\) 的二次项。重点是：当 \(\varepsilon=0\) 时系统简化；当 \(\varepsilon\) 足够小时，可用解对参数连续依赖的定理把解延拓出来。}

\subsection*{Условия существования}
\zh{存在条件}
Если \(\varepsilon\ge 0\) достаточно мало и выполнены условия
\[
C_2\le 0,
\qquad G_2\le 0,
\qquad D_{ij}>0,
\quad i,j=1,2,
\]
то ситуация равновесия по Нэшу в рассматриваемой игре не существует, но существует ситуация равновесия по Бержу--Вайсману
\[
U^B=(U_1^B,U_2^B),
\]
причем она имеет линейный позиционный вид
\[
u_1^B(t,x)=D_{21}^{-1}\Theta_2(t,\varepsilon)x,
\qquad
u_2^B(t,x)=\varepsilon D_{12}^{-1}\Theta_1(t,\varepsilon)x.
\]
Равновесные выигрыши:
\[
J_i(U^B,t_0,x_0)=V_i(t_0,x_0)=x_0^{\top}\Theta_i(t_0,\varepsilon)x_0,
\qquad i=1,2.
\]
\zh{若 \(\varepsilon\ge0\) 足够小，且 \(C_2\le0\)、\(G_2\le0\)、\(D_{ij}>0\)，则所考虑博弈中 Nash 均衡不存在，但存在 Berge--Vaisman 均衡。它具有线性位置形式 \(u_1^B=D_{21}^{-1}\Theta_2x\)、\(u_2^B=\varepsilon D_{12}^{-1}\Theta_1x\)。均衡收益为 \(J_i(U^B,t_0,x_0)=x_0^\top\Theta_i(t_0,\varepsilon)x_0\)。}

Малость \(\varepsilon\) важна не только содержательно, но и математически: при \(\varepsilon=0\) система решается проще; затем применяется метод малого параметра Пуанкаре, позволяющий представить решение в виде ряда по степеням \(\varepsilon\):
\[
\Theta_i(t,\varepsilon)=\Theta_i^{(0)}(t)+\varepsilon\Theta_i^{(1)}(t)+\varepsilon^2\Theta_i^{(2)}(t)+\cdots .
\]
\zh{小参数 \(\varepsilon\) 不仅有建模意义，也有数学意义：在 \(\varepsilon=0\) 时系统较易求解，随后可使用 Poincare 小参数方法，把解表示为关于 \(\varepsilon\) 的幂级数。}

\memory
Модель:
\[
\dot x=A(t)x+u_1+\varepsilon u_2,
\qquad 0\le \varepsilon\ll 1.
\]
Стратегии:
\[
u_i(t,x)=Q_i(t)x.
\]
Равновесие по Бержу:
\[
J_1(U_1^B,U_2)\le J_1(U^B)\quad \forall U_2,
\]
\[
J_2(U_1,U_2^B)\le J_2(U^B)\quad \forall U_1.
\]
Главное отличие от Нэша:
\[
J_1 \text{ максимизируется по } u_2,
\qquad
J_2 \text{ максимизируется по } u_1.
\]
Функции Беллмана:
\[
V_i(t,x)=x^{\top}\Theta_i(t,\varepsilon)x.
\]
Условия первого порядка:
\[
\frac{\partial W_1}{\partial u_2}=0,
\qquad
\frac{\partial W_2}{\partial u_1}=0.
\]
Отсюда
\[
u_1^B(t,x)=D_{21}^{-1}\Theta_2(t,\varepsilon)x,
\]
\[
u_2^B(t,x)=\varepsilon D_{12}^{-1}\Theta_1(t,\varepsilon)x.
\]
При достаточно малом \(\varepsilon\) и условиях
\[
C_2\le 0,
\quad G_2\le 0,
\quad D_{ij}>0,
\]
существует равновесие по Бержу--Вайсману. Выигрыши:
\[
J_i^B(t_0,x_0)=x_0^{\top}\Theta_i(t_0,\varepsilon)x_0.
\]
\zh{记忆版：模型为 \(\dot x=A(t)x+u_1+\varepsilon u_2\)，\(0\le\varepsilon\ll1\)。Berge 与 Nash 的关键区别是交叉最大化：\(J_1\) 对 \(u_2\) 最大化，\(J_2\) 对 \(u_1\) 最大化。因此 \(u_1^B=D_{21}^{-1}\Theta_2x\)，\(u_2^B=\varepsilon D_{12}^{-1}\Theta_1x\)。在相应符号条件和小参数条件下存在 Berge--Vaisman 均衡，收益为 \(J_i^B=x_0^\top\Theta_i(t_0,\varepsilon)x_0\)。}

\newpage
\ticket{20}{Максимум по Парето в двухкритериальной линейно-квадратичной динамической задаче.}{双准则线性二次动态问题中的 Pareto 最大值。}


Рассматривается двухкритериальная задача принятия решения при неопределенности или ее динамический линейно-квадратичный аналог. Лицо, принимающее решение, выбирает стратегию или управление, а неопределенный фактор заранее неизвестен.

В статической записи стратегия обозначается
\[
x\in X\subset\R^n,
\]
а неопределенность
\[
y\in Y\subset\R^m.
\]
Есть два критерия:
\[
f(x,y)\quad \text{-- выигрыш},
\]
\[
-R(x,y)\quad \text{-- отрицательный риск}.
\]
Поэтому задача записывается как
\[
G_2=\langle X,Y,\{f(x,y),-R(x,y)\}\rangle.
\]
Содержательно нужно одновременно увеличить выигрыш и уменьшить риск:
\[
f(x,y)\to\max,
\qquad
R(x,y)\to\min.
\]
Так как векторная оптимизация обычно формулируется как максимум всех критериев, второй критерий берется в виде \(-R\).
\zh{这里考虑不确定性下的双准则决策问题，或其动态线性二次类似物。决策者选择策略或控制，而不确定因素事先未知。静态写法中，策略为 \(x\in X\)，不确定性为 \(y\in Y\)。两个准则是收益 \(f(x,y)\) 和负风险 \(-R(x,y)\)。实际含义是同时增大收益并减小风险；由于向量优化通常写成所有准则都最大化，所以把第二个准则写成 \(-R\)。}

\qpart{1}{Максимум по Парето в двухкритериальной линейно-квадратичной динамической задаче}{双准则线性二次动态问题中的 Pareto 最大值}

\subsection*{Линейно-квадратичный выигрыш}
\zh{线性二次收益}
Функция выигрыша имеет вид
\[
f(x,y)=x^{\top}Ax+2x^{\top}By+y^{\top}Dy+2a^{\top}x+2d^{\top}y,
\]
где
\[
A=A^{\top},
\qquad
D=D^{\top}.
\]
Матрицы и векторы имеют согласованные размерности.
\zh{收益函数取线性二次形式 \(f(x,y)=x^\top Ax+2x^\top By+y^\top Dy+2a^\top x+2d^\top y\)，其中 \(A,D\) 对称，所有矩阵和向量维数相容。}

\subsection*{Функция риска}
\zh{风险函数}
Риск, или сожаление, определяется как потеря выигрыша по сравнению с наилучшим выбором стратегии при уже известной неопределенности \(y\):
\[
R(x,y)=\max_{z\in X}f(z,y)-f(x,y).
\]
Очевидно,
\[
R(x,y)\ge 0.
\]
Если \(R(x,y)=0\), то выбранная стратегия \(x\) оптимальна при данном \(y\).
\zh{风险或后悔值定义为：在不确定性 \(y\) 已知时可取得的最佳收益，与当前策略 \(x\) 的收益之差，即 \(R(x,y)=\max_{z\in X}f(z,y)-f(x,y)\)。显然 \(R(x,y)\ge0\)。若 \(R(x,y)=0\)，说明策略 \(x\) 在该 \(y\) 下就是最优策略。}

\subsection*{Минимальная по Парето неопределенность}
\zh{Pareto 极小不确定性}
Для фиксированной стратегии \(x\) рассматривается двухкритериальная задача по неопределенности
\[
G_2(x)=\langle Y,\{f(x,y),-R(x,y)\}\rangle.
\]
Неопределенность \(y^P(x)\in Y\) называется минимальной по Парето для фиксированного \(x\), если не существует такого \(y\in Y\), которое одновременно ухудшает оба критерия для лица, принимающего решение, то есть нельзя одновременно уменьшить выигрыш и увеличить риск. Формально несовместна система
\[
f(x,y)\le f(x,y^P(x)),
\]
\[
-R(x,y)\le -R(x,y^P(x)),
\]
причем хотя бы одно из неравенств строгое.

Эквивалентно: нельзя найти другую неопределенность, при которой выигрыш не больше, риск не меньше и хотя бы одно ухудшение строгое.
\zh{对固定策略 \(x\)，考虑关于不确定性 \(y\) 的双准则问题 \(G_2(x)=\langle Y,\{f(x,y),-R(x,y)\}\rangle\)。若不存在另一个 \(y\) 同时使决策者的两个准则变坏，也就是不能同时降低收益并增加风险，则称 \(y^P(x)\) 为固定 \(x\) 下的 Pareto 极小不确定性。等价地，不存在另一个不确定性使收益不更大、风险不更小，并且至少一个变坏严格成立。}

\subsection*{Скаляризация для построения неопределенности}
\zh{构造不确定性的标量化方法}
Обозначим
\[
f[y]=\max_{x\in X}f(x,y).
\]
Пусть существует \(\sigma\in(0,1)\) и функция \(y^P(x):X\to Y\), такая что при каждом \(x\)
\[
y^P(x)\in \arg\min_{y\in Y}\bigl(f(x,y)-\sigma f[y]\bigr).
\]
Тогда \(y^P(x)\) является минимальной по Парето неопределенностью для задачи \(G_2(x)\). Эта конструкция переводит двухкритериальное сравнение в скалярную вспомогательную задачу.
\zh{记 \(f[y]=\max_{x\in X}f(x,y)\)。如果存在 \(\sigma\in(0,1)\) 和映射 \(y^P(x)\)，使得对每个 \(x\)，\(y^P(x)\) 最小化 \(f(x,y)-\sigma f[y]\)，则 \(y^P(x)\) 是问题 \(G_2(x)\) 的 Pareto 极小不确定性。这个构造把双准则比较转化为一个辅助的标量优化问题。}

\subsection*{Максимальная по Парето стратегия}
\zh{Pareto 最大策略}
После построения минимальной по Парето неопределенности \(y^P(x)\) рассматривается задача по стратегии:
\[
\langle X,\{f(x,y^P(x)),-R(x,y^P(x))\}\rangle.
\]
Стратегия \(x^P\in X\) называется максимальной по Парето, если не существует \(x\in X\), такого что
\[
f(x,y^P(x))\ge f(x^P,y^P(x^P)),
\]
\[
-R(x,y^P(x))\ge -R(x^P,y^P(x^P)),
\]
и хотя бы одно из этих неравенств строгое. Иначе говоря, нельзя одновременно получить больший выигрыш и меньший риск.

Тройку
\[
(x^P,f^P,R^P)
\]
можно назвать гарантированным по исходам и рискам решением, где
\[
f^P=f(x^P,y^P(x^P)),
\qquad
R^P=R(x^P,y^P(x^P)).
\]
\zh{构造出 Pareto 极小不确定性 \(y^P(x)\) 后，再考虑关于策略 \(x\) 的双准则问题。若不存在另一个策略 \(x\) 同时使收益不低于 \(x^P\)、风险不高于 \(x^P\)，且至少一项严格改善，则 \(x^P\) 称为 Pareto 最大策略。换言之，不能同时获得更大收益和更小风险。三元组 \((x^P,f^P,R^P)\) 可称为按结果和风险保证的解。}

\subsection*{Динамическая линейно-квадратичная версия}
\zh{动态线性二次版本}
Официальный билет говорит именно о динамической задаче. В динамическом варианте вместо \(f(x,y)\) рассматриваются функционалы
\[
J_1(U,t_0,x_0),
\qquad
J_2(U,t_0,x_0),
\]
заданные на управляемой линейной системе. Типичная форма:
\[
\dot x=A(t)x+\sum_j u_j,
\]
\[
J_i(U,t_0,x_0)=x^{\top}(\vartheta)C_i x(\vartheta)+
\int_{t_0}^{\vartheta}\sum_j u_j^{\top} D_{ij}u_j\,dt,
\qquad i=1,2.
\]
Альтернатива \(U^P\) называется максимальной по Парето, если не существует другой альтернативы \(U\), такой что
\[
J_1(U,t_0,x_0)\ge J_1(U^P,t_0,x_0),
\]
\[
J_2(U,t_0,x_0)\ge J_2(U^P,t_0,x_0),
\]
и хотя бы одно неравенство строгое.
\zh{官方题目说的是动态问题。动态版本中，不再用静态函数 \(f(x,y)\)，而是在受控线性系统上考虑泛函 \(J_1(U,t_0,x_0)\)、\(J_2(U,t_0,x_0)\)。若不存在另一个备选控制 \(U\) 同时使两个收益都不小于 \(U^P\)，且至少一个严格更大，则 \(U^P\) 称为 Pareto 最大备选方案。}

Для построения Pareto-решения используется положительная скаляризация:
\[
J_\alpha(U)=J_1(U)+\alpha J_2(U),
\qquad \alpha>0.
\]
Если \(U^P\) максимизирует \(J_\alpha\), то \(U^P\) является максимальной по Парето альтернативой. Это стандартное достаточное условие Pareto-оптимальности: максимум положительной линейной комбинации критериев дает Pareto-максимальную альтернативу.
\zh{构造 Pareto 解时常用正权标量化 \(J_\alpha(U)=J_1(U)+\alpha J_2(U)\)，其中 \(\alpha>0\)。若 \(U^P\) 最大化 \(J_\alpha\)，则 \(U^P\) 是 Pareto 最大备选方案。这是 Pareto 最优性的标准充分条件：正线性组合的最大值给出 Pareto 最大解。}

В линейно-квадратичном случае после скаляризации снова получается LQ-задача. Вводятся взвешенные матрицы
\[
C(\alpha)=C_1+\alpha C_2,
\]
\[
D_j(\alpha)=D_{1j}+\alpha D_{2j}.
\]
Ищется функция Беллмана
\[
V^P(t,x)=x^{\top}\Theta^P(t)x.
\]
Если соответствующие матрицы имеют нужную знакоопределенность, то Pareto-максимальные управления имеют линейный позиционный вид
\[
u_j^P(t,x)=-D_j^{-1}(\alpha)\Theta^P(t)x.
\]
Матрица \(\Theta^P(t)\) удовлетворяет уравнению типа Риккати
\[
\dot\Theta^P+
\Theta^P A+A^{\top}\Theta^P-
\Theta^P\left(\sum_jD_j^{-1}(\alpha)\right)\Theta^P=0,
\]
с терминальным условием
\[
\Theta^P(\vartheta)=C(\alpha).
\]
\zh{在线性二次情形中，标量化后仍得到 LQ 问题。引入加权矩阵 \(C(\alpha)=C_1+\alpha C_2\)、\(D_j(\alpha)=D_{1j}+\alpha D_{2j}\)，寻找二次型 Bellman 函数 \(V^P=x^\top\Theta^P x\)。在符号定性条件合适时，Pareto 最大控制具有线性位置形式 \(u_j^P=-D_j^{-1}(\alpha)\Theta^P x\)，矩阵 \(\Theta^P\) 满足相应 Riccati 型方程和终端条件 \(\Theta^P(\vartheta)=C(\alpha)\)。}

Таким образом, максимум по Парето в двухкритериальной LQ-динамической задаче строится как неулучшаемый компромисс между двумя критериями, а технически может быть найден через скаляризацию и решение соответствующей задачи динамического программирования.
\zh{因此，双准则 LQ 动态问题中的 Pareto 最大值，是两个准则之间不可再共同改进的折中。从技术上看，可通过正权标量化并求解相应的动态规划问题来构造。}

\memory
Два критерия:
\[
f(x,y)\to\max,
\qquad R(x,y)\to\min.
\]
Пишем как Pareto-максимизацию:
\[
\{f(x,y),-R(x,y)\}.
\]
ЛК-выигрыш:
\[
f(x,y)=x^{\top}Ax+2x^{\top}By+y^{\top}Dy+2a^{\top}x+2d^{\top}y.
\]
Риск:
\[
R(x,y)=\max_{z\in X}f(z,y)-f(x,y)\ge0.
\]
Минимальная по Парето неопределенность \(y^P(x)\): нельзя одновременно уменьшить \(f\) и увеличить \(R\).

Достаточная конструкция:
\[
f[y]=\max_{x\in X}f(x,y),
\]
\[
y^P(x)\in\arg\min_{y\in Y}\bigl(f(x,y)-\sigma f[y]\bigr),
\qquad \sigma\in(0,1).
\]
Стратегия \(x^P\) максимальна по Парето, если не существует \(x\), такого что
\[
f(x,y^P(x))\ge f(x^P,y^P(x^P)),
\]
\[
-R(x,y^P(x))\ge -R(x^P,y^P(x^P)),
\]
и хотя бы одно неравенство строгое.

Динамическая версия:
\[
J_1(U),\quad J_2(U).
\]
Скаляризация:
\[
J_\alpha(U)=J_1(U)+\alpha J_2(U),\qquad \alpha>0.
\]
Если \(U^P\) максимизирует \(J_\alpha\), то \(U^P\) Pareto-максимальна.
В LQ-случае:
\[
V^P(t,x)=x^{\top}\Theta^P(t)x,
\qquad u^P(t,x)=Q^P(t)x,
\]
а \(\Theta^P\) находится из уравнения Риккати.
\zh{记忆版：两个准则是收益最大化和风险最小化，因此写成 Pareto 最大化 \(\{f,-R\}\)，其中 \(R=\max_z f(z,y)-f(x,y)\ge0\)。先构造 Pareto 极小不确定性 \(y^P(x)\)，再找 Pareto 最大策略。动态版本中用 \(J_1,J_2\)，通过 \(J_\alpha=J_1+\alpha J_2\) 标量化；LQ 情形下仍取 \(V^P=x^\top\Theta^P x\)，\(\Theta^P\) 由 Riccati 方程确定。}

\section*{Итоговая краткая схема блока F}
\begin{center}
\begin{tabular}{p{2cm}p{12cm}}
Билет & Что помнить \\
18 & Nash в трехлицевой LQ-игре: \(\dot x=Ax+u_1+u_2+u_3\). Ищем \(V_i=x^{\top}\Theta_i x\). При \(D_{ii}<0\): \(u_i^e=-D_{ii}^{-1}\Theta_i x\). Матрицы \(\Theta_i\) решают связанную систему Риккати. \\
19 & Berge в двухлицевой LQ-игре: \(\dot x=Ax+u_1+\varepsilon u_2\). Главное отличие: \(J_1\) максимизируется по \(u_2\), \(J_2\) по \(u_1\). \(u_1^B=D_{21}^{-1}\Theta_2x\), \(u_2^B=\varepsilon D_{12}^{-1}\Theta_1x\). \\
20 & Pareto максимум: одновременно увеличить выигрыш и уменьшить риск. Используем критерии \(\{f,-R\}\), где \(R=\max_z f(z,y)-f(x,y)\). В динамической LQ-задаче строится через скаляризацию \(J_\alpha=J_1+\alpha J_2\) и Riccati. \\
\end{tabular}
\end{center}
\zh{F 板块总结：第18题是三人 LQ 博弈中的 Nash 均衡，核心是 \(D_{ii}<0\)、线性策略 \(u_i^e=-D_{ii}^{-1}\Theta_i x\) 和耦合 Riccati 系统；第19题是两人 LQ 博弈中的 Berge 均衡，关键区别是交叉最大化，得到 \(u_1^B=D_{21}^{-1}\Theta_2x\)、\(u_2^B=\varepsilon D_{12}^{-1}\Theta_1x\)；第20题是 Pareto 最大值，同时考虑收益和风险，用 \(\{f,-R\}\) 与正权标量化 \(J_\alpha=J_1+\alpha J_2\) 构造。}

\section*{Список использованных материалов}
\begin{enumerate}
\item Программа государственного экзамена кафедры оптимального управления ВМК МГУ, 2026.
\item Жуковский В.И., Смирнова Л.В., Высокос М.И. \emph{Об одной нерешенной задаче в матричных обыкновенных дифференциальных уравнениях}.
\item Zhukovskiy V.I., Smirnova L.V. \emph{Berge--Vaisman equilibrium for one linear-quadratic differential game}.
\item Жуковский В.И., Мухина Ю.С., Романова В.Э. \emph{Дифференциальная игра N лиц, в которой существует паретовское равновесие угроз и контругроз, но отсутствует равновесие по Нэшу}.
\item Студенческие материалы 2017 г. и краткая версия 2026 г.; использовались как вспомогательные материалы, а не как основной источник.
\end{enumerate}
\zh{使用材料包括官方考试大纲、Жуковский 及合作者关于矩阵微分方程、Berge--Vaisman 均衡和无 Nash 均衡情形的文章，以及 2017 年学生材料和 2026 年简版材料。}

\blocktitle{G}{Риск, гарантированные решения и равновесия при неопределенности}{Билеты 21--23}
\section*{Общая логика блока G}
Билеты 21--23 объединены одной схемой:
\[
\text{неопределенность} \;\longrightarrow\; \text{риск / гарантия} \;\longrightarrow\; \text{равновесие по Нэшу или Парето}.
\]

Основные понятия блока:
\begin{itemize}
  \item функция выигрыша $f(x,y)$ или $f_i(x,y)$;
  \item неопределенность $y\in Y$;
  \item риск, или сожаление,
  \[
  R(x,y)=\max_{z\in X} f(z,y)-f(x,y);
  \]
  \item гарантированное решение: выбор стратегии, учитывающий худшую реализацию неопределенности;
  \item Парето-подход: одновременно учитываются несколько критериев, например $f_1,-R_1,f_2,-R_2$.
\end{itemize}

В этом блоке важно не менять направления оптимизации: выигрыши максимизируются, риски минимизируются, поэтому в многокритериальной записи риски входят как $-R_i$.

\ticket{21}{Необходимые и достаточные условия существования стратегий с нулевым риском.}{零风险策略存在的必要充分条件。}

\qpart{1}{Необходимые условия существования стратегий с нулевым риском}{零风险策略存在的必要条件}

Рассмотрим однокритериальную задачу принятия решений при неопределенности
\[
\langle X,Y,f(x,y)\rangle,
\]
где $x\in X\subset\R^n$ --- стратегия лица, принимающего решение, $y\in Y\subset\R^m$ --- неопределенный фактор, а $f(x,y)$ --- функция выигрыша. Лицо, принимающее решение, стремится выбрать $x$, обеспечивая большой выигрыш при возможной реализации любого $y\in Y$.

\zh{考虑不确定性下的单准则决策问题，其中 $x$ 是决策者选择的策略，$y$ 是不确定因素，$f(x,y)$ 是收益函数。决策者希望选择一个策略，使得在任何可能的不确定性实现下都能获得较高收益。}

\subsubsection{Функция риска}
Функция риска, или сожаления, определяется формулой
\[
\Phi(x,y)=\max_{z\in X} f(z,y)-f(x,y).
\]
Здесь $\max_{z\in X}f(z,y)$ --- наилучший выигрыш, который можно было бы получить, если бы неопределенность $y$ была известна заранее. Поэтому $\Phi(x,y)$ показывает потерю выигрыша из-за выбора стратегии $x$ при реализации $y$.

\zh{风险函数，也称为后悔值，定义为事后已知不确定性时可获得的最优收益与实际选择策略所得收益之差。因此，它衡量的是由于选择策略 $x$ 而在不确定性 $y$ 实现时损失的收益。}

Всегда
\[
\Phi(x,y)\ge 0,
\]
так как
\[
\max_{z\in X}f(z,y)\ge f(x,y).
\]
Риск равен нулю тогда и только тогда, когда $x$ оптимальна при данном $y$:
\[
\Phi(x,y)=0
\quad\Longleftrightarrow\quad
f(x,y)=\max_{z\in X}f(z,y).
\]

\zh{风险总是非负的，因为最优收益不小于任意给定策略的收益。风险等于零，当且仅当策略 $x$ 在给定的不确定性 $y$ 下本身就是最优策略。}

\subsubsection{Гарантированное по риску решение}
Для фиксированной стратегии $x$ естественно рассматривать наихудший риск:
\[
\max_{y\in Y}\Phi(x,y).
\]
Гарантированная по риску стратегия определяется задачей
\[
\varphi^r=\min_{x\in X}\max_{y\in Y}\Phi(x,y),
\]
\[
x^r\in \argmin_{x\in X}\max_{y\in Y}\Phi(x,y).
\]
Число $\varphi^r$ называется гарантированным риском. При компактности $X,Y$ и непрерывности $f$ соответствующие максимумы и минимумы существуют.

\zh{对固定策略 $x$，自然要考察其最坏情况下的风险。按风险保证的策略就是使最大风险最小的策略；相应的最小值称为保证风险。在 $X,Y$ 紧且 $f$ 连续时，相关的最大值和最小值都存在。}

\subsubsection{Стратегия с нулевым риском}
Стратегия $x^0\in X$ называется стратегией с нулевым риском, если
\[
\Phi(x^0,y)=0\qquad \forall y\in Y.
\]
Эквивалентно,
\[
f(x^0,y)=\max_{z\in X}f(z,y)\qquad \forall y\in Y.
\]
То есть одна и та же стратегия $x^0$ должна быть оптимальной при всех возможных реализациях неопределенности.

\zh{如果某个策略 $x^0$ 对所有不确定性 $y$ 都使风险为零，则它称为零风险策略。等价地，同一个策略 $x^0$ 必须在所有可能的不确定性实现下都达到对应的最大收益。}

Главное необходимое и достаточное условие:
\[
\bigcap_{y\in Y}\Argmax_{z\in X} f(z,y)\ne\varnothing.
\]
Если это пересечение непусто, то любой его элемент является стратегией с нулевым риском.

\zh{主要的必要充分条件是：所有不确定性情形下最优策略集合的交非空。如果这个交集非空，那么其中任意元素都是零风险策略。}

\qpart{2}{Достаточные условия существования стратегий с нулевым риском}{零风险策略存在的充分条件}

\subsubsection{Теорема об эквивалентных условиях}
Пусть для каждого $y\in Y$ максимум $\max_{z\in X}f(z,y)$ достигается. Тогда следующие условия эквивалентны:
\begin{enumerate}
  \item существует стратегия с нулевым риском:
  \[
  \exists x^0\in X:\quad \Phi(x^0,y)=0\quad \forall y\in Y;
  \]
  \item пересечение множеств оптимальных стратегий непусто:
  \[
  \bigcap_{y\in Y}\Argmax_{z\in X} f(z,y)\ne\varnothing;
  \]
  \item минимаксное значение риска равно нулю:
  \[
  \min_{x\in X}\max_{y\in Y}\Phi(x,y)=0;
  \]
  \item функция риска имеет седловую точку нулевого значения:
  \[
  \max_{y\in Y}\Phi(x^0,y)=\Phi(x^0,y^0)=\min_{x\in X}\Phi(x,y^0)=0.
  \]
\end{enumerate}

\zh{等价条件定理说明：若对每个 $y$ 最大值都能达到，则零风险策略存在、所有最优策略集合交非空、风险的极小极大值为零，以及风险函数存在零值鞍点，这些条件彼此等价。}

Стандартное минимаксное неравенство имеет направление
\[
\max_{y\in Y}\min_{x\in X}\Phi(x,y)\le
\min_{x\in X}\max_{y\in Y}\Phi(x,y).
\]
Это направление важно: обратное неравенство в общем случае неверно.

\zh{标准极小极大不等式的方向必须写正确：先对 $x$ 取最小再对 $y$ 取最大，一般大于等于相反顺序的值。反向不等式通常不成立。}

\subsubsection{Доказательство}
Если существует $x^0$ с нулевым риском, то
\[
\Phi(x^0,y)=0\quad \forall y\in Y.
\]
Так как $\Phi\ge 0$, имеем
\[
\max_{y\in Y}\Phi(x^0,y)=0,
\]
следовательно,
\[
\min_{x\in X}\max_{y\in Y}\Phi(x,y)=0.
\]
Кроме того,
\[
f(x^0,y)=\max_{z\in X}f(z,y)\quad \forall y,
\]
поэтому $x^0$ принадлежит всем множествам $\Argmax_{z\in X}f(z,y)$.

\zh{证明的正向部分如下：如果存在零风险策略 $x^0$，则它对所有 $y$ 都使风险为零。由于风险非负，最大风险也为零，从而极小极大风险为零；同时 $x^0$ 属于每个不确定性下的最优策略集合。}

Обратно, если
\[
\bigcap_{y\in Y}\Argmax_{z\in X}f(z,y)\ne\varnothing,
\]
возьмем $x^0$ из этого пересечения. Тогда для любого $y\in Y$
\[
f(x^0,y)=\max_{z\in X}f(z,y),
\]
значит
\[
\Phi(x^0,y)=0\quad \forall y\in Y.
\]
Следовательно, $x^0$ имеет нулевой риск.

\zh{反过来，如果所有最优策略集合的交非空，取其中一个元素 $x^0$。则对任意 $y$，它都达到该 $y$ 下的最大收益，因此风险恒为零，所以 $x^0$ 是零风险策略。}

\newpage
\ticket{22}{Гарантированное по Парето равновесие в многошаговом варианте дуополии Курно с учетом импорта.}{考虑进口的多阶段 Cournot 双寡头模型中的 Pareto 保证均衡。}

\qpart{1}{Гарантированное по Парето равновесие}{Pareto 保证均衡}

Рассматривается бескоалиционная игра двух фирм на рынке одного продукта. Игроки выбирают объемы выпуска или корректирующие управления, а импорт выступает как неопределенный фактор.

\zh{考虑一个由两家企业组成的非合作博弈，市场上只有一种产品。企业选择产量或修正控制，而进口量作为不确定因素进入模型。}

\subsubsection{Статическая модель Курно с импортом}
В статической модели объемы выпуска фирм обозначаются
\[
q_1,q_2\ge 0,
\]
а объем импорта ---
\[
y\in Y=[0,+\infty).
\]
Суммарный объем товара на рынке:
\[
\bar q=q_1+q_2+y.
\]
Цена линейно зависит от совокупного предложения:
\[
p(\bar q)=a-b\bar q,
\qquad a>0,\quad b>0.
\]
Издержки $i$-й фирмы:
\[
C_i(q_i)=cq_i+d.
\]
Прибыль первой фирмы:
\[
\psi_1(q_1,q_2,y)=\bigl[a-b(q_1+q_2+y)\bigr]q_1-(cq_1+d),
\]
то есть
\[
\psi_1=aq_1-bq_1^2-bq_1q_2-byq_1-cq_1-d.
\]
Прибыль второй фирмы:
\[
\psi_2(q_1,q_2,y)=\bigl[a-b(q_1+q_2+y)\bigr]q_2-(cq_2+d),
\]
\[
\psi_2=aq_2-bq_1q_2-bq_2^2-byq_2-cq_2-d.
\]

\zh{在含进口的静态 Cournot 模型中，两家企业的产量记为 $q_1,q_2$，进口量记为 $y$。市场总供给等于两家企业产量与进口量之和，价格是总供给的线性函数；利润由销售收入减去成本得到。}

Для учета неопределенности вводятся гарантированные функции
\[
\Phi_i(q_1,q_2,y)=\psi_i(q_1,q_2,y)+\frac{y^2}{2},
\qquad i=1,2.
\]

\zh{为了把不确定性纳入保证框架，引入保证函数：它等于企业利润加上关于进口不确定性的二次项。}

\subsubsection{Общая схема Парето-гарантированного равновесия}
Для фиксированной ситуации $q=(q_1,q_2)$ рассматривается двухкритериальная задача по неопределенности
\[
\left\langle Y,\{\Phi_1(q,y),\Phi_2(q,y)\}\right\rangle.
\]
Сначала строится Парето-минимальная неопределенность
\[
y^P(q).
\]
Затем после подстановки $y=y^P(q)$ получается игра гарантий без неопределенности:
\[
\left\langle \{1,2\},\{X_i\}_{i=1,2},\{\Phi_i(q,y^P(q))\}_{i=1,2}\right\rangle.
\]
Ситуация $q^e=(q_1^e,q_2^e)$ является равновесием Нэша в этой игре, если
\[
\Phi_1(q_1^e,q_2^e,y^P(q^e))\ge
\Phi_1(q_1,q_2^e,y^P(q_1,q_2^e))\quad \forall q_1\ge 0,
\]
\[
\Phi_2(q_1^e,q_2^e,y^P(q^e))\ge
\Phi_2(q_1^e,q_2,y^P(q_1^e,q_2))\quad \forall q_2\ge 0.
\]

\zh{Pareto 保证均衡的一般构造是：对固定情形 $q$，先在不确定性变量上求 Pareto 极小的不确定性 $y^P(q)$；再把它代回去，得到一个没有不确定性的保证博弈；最后在该博弈中求 Nash 均衡。}

\subsubsection{Явный вид в статической модели}
Для модели Курно с импортом получается
\[
y^P(q_1,q_2)=\frac{b(q_1+q_2)}{2}.
\]
Условия первого порядка в игре гарантий дают симметричное равновесие
\[
q_1^e=q_2^e=\frac{a-c}{b(3+b)}.
\]
Соответствующая Парето-минимальная неопределенность:
\[
y^P(q^e)=\frac{a-c}{3+b}.
\]
Прибыль каждой фирмы:
\[
\psi_i^e=\frac{(a-c)^2}{b(3+b)^2}-d,
\qquad i=1,2.
\]
Гарантированные выигрыши:
\[
\Phi_i^e=\psi_i^e+\frac{(y^P(q^e))^2}{2}.
\]

\zh{在静态含进口 Cournot 模型中，可以显式得到 Pareto 极小不确定性、对称均衡产量、相应的不确定性值以及各企业的利润和保证收益。}

\qpart{2}{Многошаговый вариант дуополии Курно с учетом импорта}{考虑进口的多阶段 Cournot 双寡头模型}

\subsubsection{Многошаговый вариант}
В многошаговой версии состояние системы в момент $k$ задается вектором
\[
q(k)=(q_1(k),q_2(k)).
\]
В двухшаговой позиционной модели рассматриваются моменты $k=0,1$:
\[
q_1(k+1)=\frac{a-c}{2b}-\frac{q_2(k)}{2}-\frac{z(k)}{2}+u_1(k,q(k)),
\]
\[
q_2(k+1)=\frac{a-c}{2b}-\frac{q_1(k)}{2}-\frac{z(k)}{2}+u_2(k,q(k)),
\]
\[
q_1(0)=q_{10},\qquad q_2(0)=q_{20}.
\]
Здесь $u_i(k,q(k))$ --- позиционное управление $i$-й фирмы, $z(k)$ --- неопределенность, связанная с импортом.

\zh{在多阶段版本中，系统在时刻 $k$ 的状态由两家企业产量构成。两步位置模型在 $k=0,1$ 两个时刻上给出状态递推式，其中 $u_i$ 是第 $i$ 家企业的位置控制，$z(k)$ 是与进口有关的不确定性。}

Функции выигрыша:
\[
J_i(U,Z,q_0)=[q_i(2)]^2+\sum_{k=0}^{1}\left(q_i(k)-2u_i^2(k)+\frac{z^2(k)}{2}\right),
\qquad i=1,2.
\]

\zh{收益函数把终端产量平方和各阶段收益相加，并包含控制成本项和不确定性二次项。}

\subsubsection{Определение в многошаговой игре}
Тройка
\[
(U^e,J_1^e[q_0],J_2^e[q_0])
\]
называется гарантированным по Парето равновесием, если существует Парето-минимальная неопределенность $Z^P$, такая что:
\begin{enumerate}
  \item $Z^P$ минимальна по Парето в двухкритериальной задаче по неопределенности;
  \item $U^e=(U_1^e,U_2^e)$ является равновесием Нэша в игре без неопределенности после подстановки $Z=Z^P$;
  \item гарантированные выигрыши равны
  \[
  J_i^e[q_0]=J_i(U^e,Z^P,q_0),\qquad i=1,2.
  \]
\end{enumerate}

\zh{在多阶段博弈中，若存在 Pareto 极小的不确定性 $Z^P$，使得在代入 $Z^P$ 后的无不确定性博弈中 $U^e$ 是 Nash 均衡，并且收益按该情形计算，则三元组称为 Pareto 保证均衡。}

\subsubsection{Алгоритм построения}
Многошаговая игра решается обратной индукцией:
\[
k=2\rightarrow k=1\rightarrow k=0.
\]
В конечный момент задаются терминальные функции
\[
V_i^{(2)}(q)=q_i^2,
\qquad i=1,2.
\]
На шаге $k=1$ для фиксированных $q$ и $u$ сначала находится Парето-минимальное значение $z^P(1,q,u)$, затем при нем строятся равновесные управления
\[
u_1^e(1,q),\qquad u_2^e(1,q),
\]
после чего вычисляются функции продолжения $V_i^{(1)}(q)$.

\zh{构造算法采用反向归纳：从终端时刻开始，给出终端函数；然后在上一阶段先求 Pareto 极小不确定性，再求均衡控制，并由此计算继续价值函数。}

На шаге $k=0$ процедура повторяется с учетом уже найденных функций $V_i^{(1)}(q)$: сначала строится $z^P(0,q,u)$, затем равновесные управления $u_i^e(0,q)$.

\zh{在 $k=0$ 时重复同样步骤，但要利用已经求出的继续价值函数；先确定 $z^P(0,q,u)$，再确定均衡控制。}

Итоговые стратегии:
\[
U_i^e=\{u_i^e(0,q),u_i^e(1,q)\},
\]
а равновесные гарантированные выигрыши равны
\[
J_i^e[q_0]=V_i^{(0)}(q_0),\qquad i=1,2.
\]

\zh{最终策略由各阶段的均衡位置控制组成，均衡保证收益等于初始状态处的价值函数。}

\subsubsection{Экономический смысл}
Импорт увеличивает совокупное предложение и снижает цену:
\[
p=a-b(q_1+q_2+y).
\]
Поэтому фирмы должны учитывать не только конкурента, но и внешний неблагоприятный фактор. Парето-гарантированное равновесие означает: сначала выбирается Парето-минимальная неопределенность импорта, затем при ней строится устойчивое по Нэшу поведение фирм.

\zh{经济含义是：进口增加市场总供给并降低价格，因此企业不仅要考虑竞争对手，还要考虑外部不利因素。Pareto 保证均衡表示先选取进口不确定性的 Pareto 极小情形，再在该情形下构造 Nash 稳定行为。}

\newpage
\ticket{23}{Гарантированное по выигрышам и рискам равновесие в линейно-квадратичной игре двух лиц.}{两人线性二次游戏中按收益和风险保证的均衡。}

\qpart{1}{Гарантированное по выигрышам равновесие}{按收益保证的均衡}

Рассматривается бескоалиционная игра двух лиц при неопределенности
\[
\Gamma_2=\left\langle \{1,2\},X_1,X_2,Y,\{f_i(x,y)\}_{i=1}^2\right\rangle,
\]
где
\[
X_1=X_2=\R^n,
\qquad Y=\R^m.
\]
Стратегии игроков: $x_1\in X_1$, $x_2\in X_2$; неопределенность: $y\in Y$.

\zh{考虑不确定性下的两人非合作博弈。两名玩家的策略空间为欧氏空间，不确定性也取值于欧氏空间；策略分别为 $x_1,x_2$，不确定性为 $y$。}

\subsubsection{Линейно-квадратичные выигрыши}
Функции выигрыша имеют вид
\[
f_1(x_1,x_2,y)=x_1^\top A_1x_1+2x_1^\top x_2+y^\top D_1y+2a_1^\top x_1+2d_1^\top y,
\]
\[
f_2(x_1,x_2,y)=x_2^\top A_2x_2-2x_2^\top x_1+y^\top D_2y+2a_2^\top x_2+2d_2^\top y.
\]
Предполагается
\[
A_i=A_i^\top<0,
\qquad D_i=D_i^\top>0,
\qquad i=1,2.
\]
Условие $A_i<0$ обеспечивает существование максимума по собственной стратегии игрока, а условие $D_i>0$ используется при построении Парето-минимальной неопределенности.

\zh{收益函数具有线性二次形式。假设 $A_i$ 为负定对称矩阵，保证玩家关于自身策略的最大值存在；$D_i$ 为正定对称矩阵，用于构造 Pareto 极小不确定性。}

\subsubsection{Лучшие выигрыши игроков}
Лучший выигрыш первого игрока при фиксированных $x_2,y$:
\[
f_1[x_2,y]=\max_{x_1\in X_1}f_1(x_1,x_2,y).
\]
Так как $A_1<0$, максимум находится из условия
\[
\frac{\partial f_1}{\partial x_1}=2A_1x_1+2x_2+2a_1=0,
\]
откуда
\[
x_1^*(x_2)=-A_1^{-1}(x_2+a_1).
\]
Следовательно,
\[
f_1[x_2,y]=-(x_2+a_1)^\top A_1^{-1}(x_2+a_1)+y^\top D_1y+2d_1^\top y.
\]

\zh{第一名玩家在给定对方策略和不确定性时的最优收益，通过对 $x_1$ 最大化得到。由于 $A_1<0$，一阶条件给出最优反应，并由此得到最优收益的显式表达式。}

Аналогично
\[
f_2[x_1,y]=\max_{x_2\in X_2}f_2(x_1,x_2,y),
\]
\[
\frac{\partial f_2}{\partial x_2}=2A_2x_2-2x_1+2a_2=0,
\]
\[
x_2^*(x_1)=A_2^{-1}(x_1-a_2),
\]
\[
f_2[x_1,y]=-(x_1-a_2)^\top A_2^{-1}(x_1-a_2)+y^\top D_2y+2d_2^\top y.
\]

\zh{第二名玩家的最优收益类似得到：先写出关于 $x_2$ 的一阶条件，再求出最优反应和相应的最优收益表达式。}

\qpart{2}{Гарантированное по рискам равновесие в линейно-квадратичной игре двух лиц}{两人线性二次游戏中按风险保证的均衡}

\subsubsection{Функции риска}
Риск игрока --- это потеря относительно лучшей реакции:
\[
R_1(x_1,x_2,y)=f_1[x_2,y]-f_1(x_1,x_2,y),
\]
\[
R_2(x_1,x_2,y)=f_2[x_1,y]-f_2(x_1,x_2,y).
\]
Всегда
\[
R_i(x_1,x_2,y)\ge 0.
\]
Равенство $R_i=0$ означает, что игрок $i$ выбрал свой лучшую реакцию.

\zh{玩家的风险定义为最优反应收益与实际收益之间的差。风险始终非负；当 $R_i=0$ 时，表示第 $i$ 名玩家选择了自己的最优反应。}

\subsubsection{Определение гарантированного по выигрышам и рискам равновесия}
Пятерка
\[
(x^e,f_1^P,f_2^P,R_1^P,R_2^P)\in X\times\R^4
\]
называется гарантированным по выигрышам и рискам равновесием, если существует неопределенность $y^P\in Y$, такая что
\[
f_i^P=f_i(x^e,y^P),
\qquad R_i^P=R_i(x^e,y^P),
\qquad i=1,2,
\]
и выполнены два условия:
\begin{enumerate}
  \item $x^e=(x_1^e,x_2^e)$ является равновесием Нэша в игре при фиксированном $y=y^P$:
  \[
  f_1(x_1,x_2^e,y^P)\le f_1(x^e,y^P)\quad \forall x_1,
  \]
  \[
  f_2(x_1^e,x_2,y^P)\le f_2(x^e,y^P)\quad \forall x_2;
  \]
  \item $y^P$ минимальна по Парето в четырехкритериальной задаче
  \[
  \left\langle Y,\{f_1(x^e,y),-R_1(x^e,y),f_2(x^e,y),-R_2(x^e,y)\}\right\rangle.
  \]
\end{enumerate}

\zh{按收益和风险保证的均衡由均衡策略、两个保证收益和两个保证风险组成。它要求存在一个 Pareto 极小不确定性，使得在该不确定性下策略是 Nash 均衡，并且不确定性在由收益和负风险组成的四准则问题中为 Pareto 极小。}

Знак минус перед рисками обязателен: выигрыши максимизируются, а риски минимизируются.

\zh{风险前面的负号是必要的：收益要最大化，而风险要最小化。}

\subsubsection{Построение $y^P$}
В данной модели неопределенность входит в выигрыши только через
\[
y^\top D_i y+2d_i^\top y.
\]
Скаляризация четырех критериев приводит к квадратичной задаче по $y$ и дает
\[
y^P=-(D_1+D_2)^{-1}(d_1+d_2).
\]
В данной упрощенной модели $y^P$ не зависит от $x$, поскольку отсутствуют смешанные члены вида $x_i^\top C_i y$.

\zh{在该模型中，不确定性只通过关于 $y$ 的二次项和线性项进入收益。对四个准则进行标量化会得到关于 $y$ 的二次问题，从而得到 $y^P$ 的显式公式；在简化模型中它与 $x$ 无关。}

\subsubsection{Построение $x^e$}
При фиксированном $y=y^P$ условия Нэша имеют вид
\[
2A_1x_1^e+2x_2^e+2a_1=0,
\]
\[
2A_2x_2^e-2x_1^e+2a_2=0.
\]
То есть
\[
A_1x_1^e+x_2^e=-a_1,
\]
\[
-x_1^e+A_2x_2^e=-a_2.
\]
Эта система является самым надежным способом записи результата, поскольку явные формулы зависят от выбранной конвенции знаков.

\zh{给定 $y=y^P$ 后，Nash 条件化为关于 $x_1^e,x_2^e$ 的线性方程组。把结果写成这个方程组最稳妥，因为显式公式会随符号约定而改变。}

При используемой здесь конвенции знаков:
\[
x_1^e=(A_2A_1+E_n)^{-1}(a_2-A_2a_1),
\]
\[
x_2^e=-(A_1A_2+E_n)^{-1}(a_1+A_1a_2).
\]

\zh{在本文采用的符号约定下，可以把均衡策略写成矩阵逆形式的显式公式。}

\subsubsection{Гарантированные выигрыши и риски}
Гарантированные выигрыши:
\[
f_i^P=f_i(x^e,y^P),
\qquad i=1,2.
\]
В развернутом виде:
\[
f_1^P=(x_1^e)^\top A_1x_1^e+2(x_1^e)^\top x_2^e+(y^P)^\top D_1y^P+2a_1^\top x_1^e+2d_1^\top y^P,
\]
\[
f_2^P=(x_2^e)^\top A_2x_2^e-2(x_2^e)^\top x_1^e+(y^P)^\top D_2y^P+2a_2^\top x_2^e+2d_2^\top y^P.
\]
Используя условия Нэша, можно упростить:
\[
f_1^P=-(x_1^e)^\top A_1x_1^e+(y^P)^\top D_1y^P+2d_1^\top y^P,
\]
\[
f_2^P=-(x_2^e)^\top A_2x_2^e+(y^P)^\top D_2y^P+2d_2^\top y^P.
\]
Риски:
\[
R_1^P=f_1[x_2^e,y^P]-f_1(x^e,y^P),
\]
\[
R_2^P=f_2[x_1^e,y^P]-f_2(x^e,y^P).
\]
Так как $x^e$ является равновесием Нэша при $y=y^P$, каждый игрок выбирает лучшую реакцию. Поэтому в точной равновесной ситуации
\[
R_1^P=R_2^P=0.
\]

\zh{保证收益通过把均衡策略和 Pareto 极小不确定性代入收益函数得到。利用 Nash 条件可以化简表达式。保证风险由最优反应收益减去均衡收益得到；在精确均衡中两名玩家都采用最优反应，因此风险为零。}

\newpage
\section*{Итог блока G}
\begin{tabular}{p{0.12\textwidth}p{0.78\textwidth}}
\toprule
Билет & Ключевой смысл \\
\midrule
21 & Нулевой риск существует тогда и только тогда, когда есть одна стратегия, оптимальная при всех неопределенностях. \\
22 & В дуополии Курно с импортом сначала строится Парето-минимальная неопределенность импорта, затем равновесие Нэша в игре гарантий. \\
23 & В ЛК-игре двух лиц учитываются одновременно выигрыши и риски: критерии имеют вид $\{f_1,-R_1,f_2,-R_2\}$; после выбора $y^P$ строится Nash. \\
\bottomrule
\end{tabular}

\zh{G 板块总结：第21题说明零风险存在的充要条件；第22题先求进口的不确定性，再求保证博弈中的 Nash 均衡；第23题同时处理收益和风险，把准则写成收益与负风险的组合。}

\section*{Использованные материалы}
\begin{enumerate}
  \item Официальный список вопросов кафедры оптимального управления ВМК МГУ на 2026 год.
  \item Студенческий конспект 2017 года по билетам государственного экзамена.
  \item Краткая версия материалов 2026 года.
  \item Жуковский В.И., Кудрявцев К.Н. \textit{Уравновешивание конфликтов при неопределенности. II. Аналог максимина}.
  \item Болтянский В.Г. \textit{Оптимальное управление дискретными системами}. М.: Наука, 1973.
  \item Пропой А.И. \textit{Элементы теории оптимальных дискретных процессов}. М.: Наука, 1973.
\end{enumerate}

\zh{使用的材料包括官方题目清单、学生笔记、2026 年简版材料，以及与不确定性下冲突均衡、离散最优控制和离散过程理论相关的参考书。}

\blocktitle{H}{Экономические модели, природные ресурсы и DICE}{Билеты 24--26}
\section*{Замечание об источниках и надежности}
\zh{可靠性说明}
Для блока H использованы официальная программа 2026 года, студенческий конспект 2017 года, краткая версия 2026 года и следующие основные источники: Hritonenko--Yatsenko \textit{Mathematical Modeling in Economics, Ecology and the Environment}; Barro--Sala-i-Martin \textit{Economic Growth}; Nordhaus--Sztorc \textit{DICE 2013R: Introduction and User's Manual}; Понтрягин и др. \textit{Математическая теория оптимальных процессов}; Асеев--Кряжимский и Aseev по бесконечному горизонту. Книга Nordhaus \textit{Managing the Global Commons} в текущем наборе материалов отсутствует; для конкретной структуры модели DICE используется руководство Nordhaus--Sztorc DICE 2013R. Книга Seierstad--Sydsaeter в полном виде также отсутствует; для билета 25 основным предметным источником является Hritonenko--Yatsenko, глава о невозобновляемых ресурсах, а оптимизационная часть проверяется по Понтрягину и Асееву.
\\zh{H 板块使用的主要来源包括 2026 年官方大纲、2017 年学生笔记、2026 年简版材料，以及 Hritonenko--Yatsenko、Barro--Sala-i-Martin、Nordhaus--Sztorc、Pontryagin 等关于经济增长、环境经济与最优控制的文献。当前材料中缺少 Nordhaus《Managing the Global Commons》和 Seierstad--Sydsaeter 的完整版本，因此 DICE 的具体结构以 Nordhaus--Sztorc DICE 2013R 手册为准，第25题则主要依据 Hritonenko--Yatsenko 中关于不可再生资源的一章，并辅以 Pontryagin 与 Aseev 的最优控制理论核对。}

\zh{H 板块关于经济模型、自然资源和 DICE。该板块使用了官方大纲、学生笔记、2026 年简版材料以及经济增长、环境经济模型、DICE 和最优控制方面的主要参考文献。}

\ticket{24}{Математические модели технического прогресса (экзогенный и эндогенный прогресс). Задачи оптимизации, подходы к решению.}{技术进步的数学模型（外生和内生进步）。优化问题与求解方法。}

\qpart{1}{Экзогенный и эндогенный технический прогресс}{外生与内生技术进步}

Технический прогресс в экономических моделях означает рост выпуска, который нельзя объяснить только увеличением количества традиционных факторов производства: капитала $K$, труда $L$ и природных ресурсов. Формально это отражается либо явной зависимостью производственной функции от времени, либо введением отдельного технологического уровня $A(t)$, либо включением человеческого капитала $H(t)$, знаний, опыта, НИОКР и образовательного сектора.

\zh{经济模型中的技术进步是指产出增长不能仅由资本、劳动和自然资源等传统生产要素数量增加来解释。形式上，它可以通过生产函数显含时间、引入技术水平 $A(t)$，或加入人力资本、知识、经验、研发和教育部门来表示。}

Обычно различают
\[
\text{экзогенный технический прогресс}
\]
и
\[
\text{эндогенный технический прогресс}.
\]
Экзогенный прогресс задается вне модели. Его динамика не объясняется поведением агентов. Эндогенный прогресс возникает внутри модели как результат целенаправленной деятельности: инвестиций в образование, НИОКР, накопление человеческого капитала, обучение, распространение знаний и государственной политики.

\zh{通常区分外生技术进步和内生技术进步。外生进步由模型外给定，其动态不由经济主体行为解释；内生进步则是在模型内部由教育、研发、人力资本积累、学习、知识扩散和政策等有目的活动产生。}

\subsubsection*{1. Экзогенный технический прогресс}
В простейшей форме выпуск задается производственной функцией
\[
Y(t)=F(K(t),L(t),t),
\]
или, что более удобно в теории роста,
\[
Y(t)=F(K(t),A(t)L(t)),
\]
где $A(t)$ -- уровень технологии, а $A(t)L(t)$ -- эффективный труд. В модели экзогенного прогресса $A(t)$ задана извне, например
\[
\dot A(t)=gA(t), \qquad A(0)=A_0,
\]
то есть
\[
A(t)=A_0e^{gt}.
\]
Труд также часто считается экзогенным:
\[
\dot L(t)=nL(t), \qquad L(0)=L_0.
\]
Капитал накапливается по уравнению
\[
\dot K(t)=I(t)-\delta K(t),
\]
где $I(t)$ -- инвестиции, $\delta>0$ -- норма амортизации.

\zh{外生技术进步模型中，产出可以写成依赖资本、劳动和时间的生产函数，也常写成资本与有效劳动的函数。技术水平 $A(t)$ 外生给定，例如按指数规律增长；劳动也常外生增长，资本则由投资减去折旧而积累。}

Если используется постоянная норма накопления
\[
s=\frac{I(t)}{Y(t)}, \qquad 0<s<1,
\]
то получаем модель Солоу--Свона с экзогенным техническим прогрессом:
\[
\dot K(t)=sF(K(t),A(t)L(t))-\delta K(t).
\]
В переменных на единицу эффективного труда
\[
k(t)=\frac{K(t)}{A(t)L(t)}
\]
при постоянной отдаче от масштаба получается фундаментальное уравнение:
\[
\dot k(t)=sf(k(t))-(\delta+n+g)k(t),
\]
где $f(k)=F(k,1)$. Экономический смысл: технический прогресс $g$ повышает эффективность труда, но одновременно увеличивает ``эффективную амортизацию'' капитала на единицу эффективного труда:
\[
\delta+n+g.
\]
В этой модели $s$, $n$, $g$, $\delta$ заданы, поэтому модель Солоу--Свона является главным образом моделью динамики и сравнительной статики, а не полной задачей оптимального управления.

\zh{若采用常数积累率，就得到含外生技术进步的 Solow-Swan 模型。转为单位有效劳动变量后得到基本方程，其中技术进步提高劳动效率，同时也提高单位有效劳动资本的“有效折旧”。在此模型中参数给定，因此主要是动态与比较静态模型，而不是完整的最优控制问题。}

\qpart{2}{Задачи оптимизации}{优化问题}

\subsubsection*{2. Экзогенный прогресс в задаче оптимизации: модель Рамсея}
Если вместо заданной нормы накопления $s$ выбирается потребление $C(t)$ или доля накопления $s(t)$, то возникает задача оптимального управления. Например,
\[
\max_{C(\cdot)}\int_0^\infty e^{-\rho t}U(c(t))L(t)\,dt,
\]
где
\[
c(t)=\frac{C(t)}{L(t)}
\]
-- потребление на душу населения, $\rho>0$ -- коэффициент дисконтирования.

\zh{如果不再把积累率固定，而是选择消费或积累份额，就得到最优控制问题。Ramsey 型模型以贴现效用总和为目标，其中人均消费和贴现系数进入目标函数。}

Ограничение:
\[
\dot K(t)=F(K(t),A(t)L(t))-C(t)-\delta K(t),
\]
\[
K(0)=K_0,\qquad C(t)\ge 0,\qquad K(t)\ge 0.
\]
Здесь технический прогресс по-прежнему экзогенен:
\[
A(t)=A_0e^{gt}.
\]
Переходя к переменным на единицу эффективного труда, можно записать:
\[
\dot k(t)=f(k(t))-c_e(t)-(\delta+n+g)k(t),
\]
где $c_e(t)=C(t)/(A(t)L(t))$.

\zh{约束由资本积累方程、初始资本以及消费和资本的非负条件构成。技术进步仍是外生给定的。转到单位有效劳动变量后，可以写出含消费和有效折旧项的状态方程。}

Гамильтониан в current-value форме:
\[
\mathcal H=U(c)+\lambda\left[f(k)-c-(\delta+n+g)k\right].
\]
Условие максимума:
\[
\frac{\partial \mathcal H}{\partial c}=0
\quad\Longrightarrow\quad
U'(c)=\lambda.
\]
Сопряженное уравнение:
\[
\dot\lambda=\lambda\left(\rho+\delta+n+g-f'(k)\right),
\]
в current-value записи, с учетом выбранной нормировки. Условие трансверсальности на бесконечности обычно записывается в виде
\[
\lim_{t\to\infty}e^{-\rho t}\lambda(t)k(t)=0,
\]
или в эквивалентной форме для present-value сопряженной переменной.

\zh{用现值形式写 Hamilton 函数后，极大值条件给出边际效用等于伴随变量；伴随方程描述影子价格的动态。无限期问题还需要横截性条件，常写成贴现后的影子价值乘状态在无穷远处趋于零。}

Таким образом, в модели Рамсея технический прогресс может быть экзогенным, но задача уже является оптимизационной: выбирается траектория потребления и накопления.

\zh{因此，在 Ramsey 模型中技术进步可以是外生的，但问题已经是优化问题：需要选择消费和积累的时间轨迹。}

\subsubsection*{3. Эндогенный технический прогресс}
В моделях эндогенного прогресса уровень технологии, человеческий капитал или запас знаний становятся внутренними переменными модели. Их динамика зависит от решений экономических агентов. Общая идея:
\[
Y(t)=F(K(t),L(t),A(t),H(t)),
\]
где $A(t)$ -- запас знаний или технологический уровень, $H(t)$ -- человеческий капитал. В отличие от экзогенной модели, здесь $A(t)$ и $H(t)$ не задаются заранее, а удовлетворяют собственным динамическим уравнениям.

\zh{在内生技术进步模型中，技术水平、人力资本或知识存量成为模型内部变量。它们的动态由经济主体的决策决定，而不是预先给定。}

\subsubsection*{4. Модель с человеческим капиталом}
Один из стандартных вариантов -- модель Узавы--Лукаса. Пусть
\[
u(t)\in[0,1]
\]
-- доля человеческого капитала, используемая в производстве. Тогда $1-u(t)$ -- доля, направляемая на образование и накопление человеческого капитала.

\zh{一个标准例子是 Uzawa-Lucas 模型。控制 $u(t)$ 表示用于生产的人力资本份额，$1-u(t)$ 则用于教育和人力资本积累。}

Производственная функция:
\[
Y(t)=K^\alpha(t)(u(t)H(t))^{1-\alpha},\qquad 0<\alpha<1.
\]
Динамика физического капитала:
\[
\dot K(t)=Y(t)-C(t)-\delta K(t).
\]
Динамика человеческого капитала:
\[
\dot H(t)=B(1-u(t))H(t),\qquad B>0.
\]
Задача оптимизации:
\[
\max_{C(\cdot),u(\cdot)}\int_0^\infty e^{-\rho t}U(C(t))\,dt
\]
при ограничениях
\[
C(t)\ge 0,\qquad 0\le u(t)\le 1,\qquad K(t)\ge 0,\quad H(t)\ge 0.
\]
Здесь управление $u(t)$ распределяет человеческий капитал между текущим производством и образованием. Если $u(t)$ велико, экономика больше производит сегодня; если $u(t)$ мало, больше ресурсов идет на накопление знаний и будущий рост.

\zh{模型给出生产函数、物质资本动态和人力资本动态，并以贴现效用最大化为目标。控制 $u(t)$ 在当前生产和教育之间分配人力资本：$u(t)$ 大则当前产出高，$u(t)$ 小则更多资源用于知识积累和未来增长。}

\subsubsection*{5. Модели с НИОКР и накоплением знаний}
Другой класс моделей эндогенного прогресса описывает технологический уровень $A(t)$ как результат НИОКР:
\[
\dot A(t)=\Phi(A(t),L_A(t),A_F(t)),
\]
где $L_A(t)$ -- трудовые ресурсы, направленные в научно-исследовательский сектор, а $A_F(t)$ -- внешний технологический уровень, например уровень страны-лидера.

\zh{另一类内生进步模型把技术水平 $A(t)$ 描述为研发活动的结果，其中 $L_A(t)$ 是投入科研部门的劳动，$A_F(t)$ 可以表示外部或领先国家的技术水平。}

Простые варианты:
\[
\dot A(t)=\beta L_A(t)A(t),
\]
или с внешними знаниями:
\[
\dot A(t)=\beta L_A(t)(A_F(t)-A(t)).
\]
Разность $A_F(t)-A(t)$ характеризует технологический разрыв между лидером и последователем. Тогда задача оптимизации имеет вид
\[
\max \int_0^\infty e^{-\rho t}U(C(t))\,dt
\]
при ограничениях
\[
Y(t)=F(K(t),A(t)L_Y(t)),
\]
\[
\dot K(t)=Y(t)-C(t)-\delta K(t),
\]
\[
\dot A(t)=\Phi(A(t),L_A(t),A_F(t)),
\]
\[
L_Y(t)+L_A(t)=L(t).
\]
Управления -- $C(t)$, $L_A(t)$, $L_Y(t)$ или доля труда в НИОКР
\[
v(t)=\frac{L_A(t)}{L(t)}.
\]

\zh{简单形式可以使技术增长取决于科研劳动和已有知识，也可以取决于与外部前沿技术之间的差距。优化问题通常选择消费、科研劳动、生产劳动或科研劳动份额，并满足产出、资本积累、技术积累和劳动分配约束。}

\qpart{3}{Подходы к решению}{求解方法}

\subsubsection*{6. Подходы к решению задач оптимизации}
Для моделей технического прогресса используются три основных подхода.

\zh{技术进步模型的优化问题主要用三类方法处理：Pontryagin 极大值原理、Bellman 方法，以及稳态或平衡增长路径分析。}

\paragraph{Принцип максимума Понтрягина.}
Для задачи
\[
\max_{u(\cdot)}\int_0^\infty e^{-\rho t}F_0(x(t),u(t))\,dt,
\qquad \dot x=f(x,u),
\]
строится гамильтониан в current-value форме
\[
H(x,u,\psi)=F_0(x,u)+\psi^\top f(x,u).
\]
Необходимые условия:
\[
\dot x=\frac{\partial H}{\partial \psi},\qquad
\dot\psi=\rho\psi-\frac{\partial H}{\partial x},
\]
\[
H(x^*,u^*,\psi)=\max_{u\in U}H(x^*,u,\psi).
\]
Для бесконечного горизонта обязательно требуется условие трансверсальности, например
\[
\lim_{t\to\infty}e^{-\rho t}\psi(t)x(t)=0,
\]
или более точная форма, зависящая от модели и состояния.

\zh{使用 Pontryagin 极大值原理时，先建立现值形式的 Hamilton 函数，再写出状态方程、伴随方程和极大值条件。对于无限期问题还必须加入横截性条件，其具体形式依赖于模型和状态变量。}

\paragraph{Метод Беллмана.}
Если модель стационарна, можно искать функцию ценности
\[
V(x)=\sup_{u(\cdot)}\int_0^\infty e^{-\rho t}F_0(x(t),u(t))\,dt.
\]
Тогда формально она удовлетворяет уравнению Гамильтона--Якоби--Беллмана:
\[
\rho V(x)=\max_{u\in U}\{F_0(x,u)+V_x(x)f(x,u)\}.
\]
Этот метод удобен для построения стационарных обратных связей $u^*=u^*(x)$.

\zh{Bellman 方法在模型平稳时寻找价值函数，并把问题写为 Hamilton-Jacobi-Bellman 方程。该方法适合构造稳态反馈控制。}

\paragraph{Стационарные траектории и сбалансированный рост.}
В моделях роста часто ищут steady state или balanced growth path. Для переменных на единицу эффективного труда стационарное состояние определяется условиями
\[
\dot k=0,\qquad \dot c=0,
\]
или, в более сложных моделях,
\[
\dot k=0,\qquad \dot h=0,\qquad \dot A/A=\mathrm{const}.
\]
В эндогенных моделях особый интерес представляет возможность положительного долгосрочного темпа роста без экзогенно заданного $g$.

\zh{在增长模型中，经常研究稳态或平衡增长路径。对单位有效劳动变量，稳态由相关变量导数为零刻画；在内生增长模型中，特别重要的是是否能在没有外生给定增长率的情况下产生正的长期增长率。}

\ticket{25}{Математические модели оптимальной добычи невозобновляемых ресурсов. Задачи оптимизации, подходы к решению.}{不可再生资源最优开采的数学模型。优化问题与求解方法。}

\qpart{1}{Модели оптимальной добычи невозобновляемых ресурсов}{不可再生资源的最优开采模型}

Невозобновляемый ресурс -- это ресурс, запас которого за рассматриваемый экономический период не восстанавливается естественным образом. Примеры: нефть, газ, уголь, руды, редкие металлы. Главная особенность таких моделей состоит в том, что текущая добыча уменьшает будущие возможности. Поэтому запас ресурса является фазовой переменной, а интенсивность добычи -- управлением.

\zh{不可再生资源是在所考察经济时期内不能自然恢复其存量的资源，例如石油、天然气、煤、矿石和稀有金属。这类模型的核心特点是：当前开采会减少未来可用资源，因此资源存量是状态变量，开采强度是控制变量。}

Обозначим:
\[
R(t)\ge 0
\]
-- текущий запас ресурса в недрах,
\[
R(0)=R_0>0
\]
-- начальный запас,
\[
E(t)\ge 0
\]
-- интенсивность добычи,
\[
p(t)>0
\]
-- рыночная цена единицы добытого ресурса,
\[
C(E,R)
\]
-- затраты на добычу,
\[
r>0
\]
-- ставка дисконтирования.

\zh{记 $R(t)$ 为地下当前资源存量，$R_0$ 为初始存量，$E(t)$ 为开采强度，$p(t)$ 为单位资源市场价格，$C(E,R)$ 为开采成本，$r$ 为贴现率。}

Динамика запаса:
\[
\dot R(t)=-E(t).
\]
Обычно предполагается ограничение
\[
0\le E(t)\le E_{\max}.
\]
Экономический смысл: $E(t)>0$ означает добычу и продажу ресурса сейчас; $E(t)=0$ означает сохранение ресурса в недрах для будущей добычи.

\zh{资源存量的动态方程表明存量以开采强度的速度减少。通常还假设开采强度有上下界。经济含义是：正开采表示现在出售资源，零开采表示把资源留在地下供未来开采。}

\qpart{2}{Задачи оптимизации}{优化问题}

\subsubsection*{1. Базовая задача оптимальной добычи}
Типичная задача фирмы состоит в максимизации дисконтированной прибыли:
\[
J(E,T)=\int_0^T e^{-rt}\left[p(t)E(t)-C(E(t),R(t))\right]dt\to \max_{E(\cdot),T}.
\]
Ограничения:
\[
\dot R(t)=-E(t),\qquad R(0)=R_0,
\]
\[
R(t)\ge 0,
\qquad
0\le E(t)\le E_{\max}.
\]
Возможны два варианта горизонта: $T$ задан заранее; либо $T$ свободен, и нужно определить оптимальный момент завершения разработки месторождения. В модели с полным исчерпанием часто задают терминальное условие
\[
R(T)=0.
\]
Если $T$ свободен, задача состоит в нахождении $E^*(t)$, $T^*$, $R^*(t)$ и $J^*$.

\zh{基本最优开采问题是最大化贴现利润，约束包括资源存量动态、初始存量、存量非负和开采强度限制。时间区间可以预先给定，也可以自由选择；在完全耗竭模型中常设终端条件 $R(T)=0$。}

\subsubsection*{2. Экономические предположения об издержках}
Обычно предполагается
\[
C_E(E,R)\ge 0,
\]
то есть чем выше интенсивность добычи, тем больше издержки. Также часто предполагается
\[
C_R(E,R)\le 0.
\]
Это означает: чем больше оставшийся запас $R$, тем ниже издержки добычи. Когда ресурс истощается, добыча становится дороже.

\zh{关于成本通常假设：开采强度越高，成本越大；剩余资源越多，开采成本越低。当资源逐渐枯竭时，开采会变得更昂贵。}

В простых моделях возможны частные случаи:
\[
C(E,R)=0,
\]
\[
C(E,R)=C(E),
\]
\[
C(E,R)=cE.
\]

\zh{在简单模型中，成本函数可以取为零、只依赖开采强度，或为开采强度的线性函数。}

\qpart{3}{Подходы к решению}{求解方法}

\subsubsection*{3. Принцип максимума Понтрягина}
Запишем гамильтониан в present-value форме:
\[
H(t,R,E,\psi)=e^{-rt}\left[p(t)E-C(E,R)\right]-\psi E.
\]
Здесь $\psi(t)$ -- сопряженная переменная. Экономически она является теневой ценой единицы ресурса в недрах в present-value измерении.

\zh{Pontryagin 极大值原理中，先写出现值形式的 Hamilton 函数。伴随变量 $\psi(t)$ 在经济上表示地下单位资源的影子价格。}

Сопряженное уравнение:
\[
\dot\psi(t)=-\frac{\partial H}{\partial R}=e^{-rt}C_R(E(t),R(t)).
\]
Условие максимума:
\[
E^*(t)\in \argmax_{0\le E\le E_{\max}}\left\{e^{-rt}[p(t)E-C(E,R^*(t))]-\psi(t)E\right\}.
\]
Если максимум внутренний, то
\[
\frac{\partial H}{\partial E}=0,
\]
следовательно,
\[
e^{-rt}\left[p(t)-C_E(E^*(t),R^*(t))\right]-\psi(t)=0.
\]
То есть
\[
\psi(t)=e^{-rt}\left[p(t)-C_E(E^*(t),R^*(t))\right].
\]
В current-value форме положим
\[
\lambda(t)=e^{rt}\psi(t).
\]
Тогда внутреннее условие принимает вид
\[
\lambda(t)=p(t)-C_E(E^*(t),R^*(t)).
\]
Сопряженное уравнение в current-value форме:
\[
\dot\lambda(t)=r\lambda(t)+C_R(E^*(t),R^*(t)).
\]
Если $C_R=0$, то
\[
\dot\lambda=r\lambda.
\]
Это форма правила Хотеллинга для чистой ренты.

\zh{伴随方程、极大值条件和内部最优条件给出资源影子价格与边际净收益之间的关系。改用现值变量后，可得到 current-value 形式的伴随方程；当成本不依赖存量时，它化为 Hotelling 规则的形式。}

\subsubsection*{4. Bang-bang структура при линейной прибыли}
Если $C(E,R)=0$, то гамильтониан линеен по $E$:
\[
H=e^{-rt}p(t)E-\psi E=(e^{-rt}p(t)-\psi(t))E.
\]
Оптимальное управление определяется знаком функции переключения
\[
\Delta(t)=e^{-rt}p(t)-\psi(t).
\]
Если $\Delta(t)>0$, то выгодно добывать максимально:
\[
E^*(t)=E_{\max}.
\]
Если $\Delta(t)<0$, то добычу выгодно отложить:
\[
E^*(t)=0.
\]
Если $\Delta(t)=0$, возможен особый режим или промежуточная добыча.

\zh{若成本为零，则 Hamilton 函数关于开采量是线性的，最优控制由切换函数的符号决定。切换函数为正时最大开采，为负时推迟开采，为零时可能出现奇异弧或中间开采。}

\subsubsection*{5. Условия на конце}
Если момент окончания $T$ фиксирован и требуется полное исчерпание ресурса,
\[
R(T)=0,
\]
то задача решается как задача с фиксированным концом по фазовой переменной. Если $T$ свободен и нет терминального выигрыша, дополнительно появляется условие трансверсальности по свободному времени:
\[
H(T)=0.
\]
Если конечный запас не фиксирован, то при отсутствии терминальной ценности остатка возникает условие
\[
\psi(T)=0.
\]
Но в типичной задаче разработки месторождения предполагается полное исчерпание $R(T)=0$, поэтому нельзя одновременно без пояснения писать $R(T)=0$ и $\psi(T)=0$: это разные варианты терминальных условий.

\zh{终端条件取决于问题设定。如果终端时刻固定并要求完全耗竭，则是状态终点固定的问题；若终端时刻自由且无终端收益，则出现自由时间横截性条件。若终端存量不固定且剩余资源无终端价值，则伴随变量终值为零。不能在没有说明的情况下同时写 $R(T)=0$ 和 $\psi(T)=0$，因为它们对应不同终端条件。}

\subsubsection*{6. Правило Хотеллинга}
Классическая модель Хотеллинга рассматривает конкурентный рынок невозобновляемого ресурса. Основная идея: ресурс в недрах является капитальным активом. Поэтому его чистая цена должна приносить тот же темп доходности, что и другие активы.

\zh{经典 Hotelling 模型研究不可再生资源的竞争市场。核心思想是地下资源是一种资本资产，因此其净价格应当带来与其他资产相同的收益率。}

Обозначим чистую цену, или ренту ресурса:
\[
\pi(t)=p(t)-c(t),
\]
где $c(t)$ -- предельные издержки добычи. Правило Хотеллинга:
\[
\frac{\dot\pi(t)}{\pi(t)}=r.
\]
Если издержки нулевые, $c(t)=0$, то
\[
\frac{\dot p(t)}{p(t)}=r,
\qquad
p(t)=p_0e^{rt}.
\]
Смысл: если цена ресурса растет медленнее, чем ставка процента, выгоднее добыть ресурс сейчас и вложить выручку в финансовый актив. Если цена растет быстрее, выгоднее оставить ресурс в недрах.

\zh{把资源净价格或资源租金记为价格减去边际开采成本。Hotelling 规则说明，资源租金的增长率等于利率。若资源价格增长慢于利率，则现在开采并投资收益更有利；若增长快于利率，则把资源留在地下更有利。}

\subsubsection*{7. Модель Хотеллинга с заданным спросом}
Пусть спрос на ресурс имеет изоэластический вид
\[
E(t)=d_0p(t)^{-\alpha},\qquad \alpha>0.
\]
При правиле Хотеллинга
\[
p(t)=p_0e^{rt}.
\]
Тогда
\[
E(t)=d_0p_0^{-\alpha}e^{-\alpha rt}.
\]
Постоянная $p_0$ определяется из условия исчерпания начального запаса:
\[
\int_0^T E(t)\,dt=R_0.
\]
В результате оптимальная добыча убывает со временем:
\[
E(t)=\frac{r\alpha R_0}{1-e^{-r\alpha T}}e^{-r\alpha t}
\]
для конечного $T$. При бесконечном горизонте:
\[
E(t)=r\alpha R_0e^{-r\alpha t}.
\]

\zh{若需求为等弹性形式，并结合 Hotelling 价格规律，就可得到随时间下降的最优开采路径。常数由初始存量耗竭条件确定；有限期和无限期情形分别给出不同形式。}

\subsubsection*{8. Модификации модели Хотеллинга}
Классическая модель Хотеллинга слишком упрощена. Поэтому вводятся модификации: издержки добычи $C(E,R)$; технологический прогресс в добыче; неопределенность спроса и цены; открытие новых запасов и разведка; наличие альтернативной технологии; связь добычи с общим экономическим ростом.

\zh{经典 Hotelling 模型过于简化，因此常加入开采成本、开采技术进步、需求和价格不确定性、新储量发现与勘探、替代技术以及与宏观经济增长的联系等修正。}

\subsubsection*{9. Модель Дасгупты--Хила}
Более общий класс моделей связывает невозобновляемый ресурс с производством и потреблением. Пусть $K(t)$ -- капитал, $R(t)$ -- запас ресурса, $E(t)$ -- добыча ресурса, $C(t)$ -- потребление, $I(t)$ -- инвестиции. Производственная функция:
\[
Y(t)=F(K(t),E(t)).
\]
Баланс выпуска:
\[
Y(t)=C(t)+I(t).
\]
Динамика капитала:
\[
\dot K(t)=I(t)-\delta K(t).
\]
Динамика ресурса:
\[
\dot R(t)=-E(t).
\]
Задача социального планировщика:
\[
\max_{C(\cdot),E(\cdot)}\int_0^\infty e^{-\rho t}U(C(t))\,dt
\]
при ограничениях
\[
C(t)\ge 0,
\quad
E(t)\ge 0,
\quad
K(t)\ge 0,
\quad
R(t)\ge 0.
\]
Экономический смысл: ресурс уже не просто продается как товар, а является фактором производства. Поэтому оптимальная политика должна учитывать компромисс между текущим потреблением, накоплением капитала и сохранением ресурса.

\zh{Dasgupta-Heal 模型把不可再生资源与生产和消费联系起来。资源不只是商品，而是生产要素；因此最优政策必须在当前消费、资本积累和资源保存之间进行权衡。}

\subsubsection*{10. Подходы к решению}
Основные подходы: принцип максимума Понтрягина, метод Беллмана, анализ особых режимов и переключений, численное решение.

\zh{主要求解方法包括 Pontryagin 极大值原理、Bellman 方法、奇异弧与切换分析，以及数值求解。}

Для стационарной бесконечногоризонтной задачи можно ввести функцию ценности
\[
V(R)=\sup_{E(\cdot)}\int_0^\infty e^{-rt}\left[pE-C(E,R)\right]dt.
\]
Тогда формально:
\[
rV(R)=\max_{0\le E\le E_{\max}}\{pE-C(E,R)-V_R(R)E\}.
\]
Здесь $V_R(R)$ -- текущая теневая цена ресурса.

\zh{对平稳的无限期问题，可以引入价值函数，并形式化地写出 Hamilton-Jacobi-Bellman 方程。其中价值函数对资源存量的导数表示资源的当前影子价格。}

При сложных функциях $C(E,R)$, $p(t)$, многомерных состояниях или неопределенности задача решается численно: прямой дискретизацией, динамическим программированием, методами нелинейного программирования или shooting method для системы ПМП.

\zh{当成本函数、价格、状态维度或不确定性较复杂时，通常采用数值方法，例如直接离散化、动态规划、非线性规划，或对 Pontryagin 条件系统使用 shooting 方法。}

\ticket{26}{Интегрированные модели климата и экономики. Модель DICE. Задачи оптимизации, подходы к решению.}{综合气候与经济模型。DICE 模型。优化问题与求解方法。}
\qpart{1}{Интегрированные модели климата и экономики}{综合气候与经济模型}

Интегрированные модели климата и экономики, или \textit{integrated assessment models} (IAM), описывают совместную динамику экономической системы и климатической системы. Их задача -- связать
\[
\text{производство}\to \text{энергопотребление}\to \text{выбросы CO}_2
\to \text{концентрация углерода}
\to \text{температура}
\to \text{климатический ущерб}
\to \text{экономическое благосостояние}.
\]
Модель DICE -- это агрегированная глобальная интегрированная модель климата и экономики. Аббревиатура означает
\[
\text{DICE}=\text{Dynamic Integrated model of Climate and the Economy}.
\]
В отличие от чисто климатических моделей, DICE содержит экономический оптимизационный блок. В отличие от обычных моделей экономического роста, DICE добавляет природный капитал климатической системы: концентрация парниковых газов рассматривается как отрицательный природный капитал, а снижение выбросов -- как инвестиции в сохранение климата.
\zh{与纯气候模型不同，DICE 含有经济最优化模块；与一般经济增长模型不同，它把气候系统中的温室气体浓度看成一种负的自然资本，把减排看成维护气候的投资。}

\subsubsection*{1. Назначение IAM и DICE}
Интегрированные модели бывают двух типов: policy evaluation models и policy optimization models. DICE относится главным образом ко второму типу. В ней выбираются такие траектории сбережений и климатической политики, которые максимизируют функцию общественного благосостояния.

Цель DICE -- найти оптимальный компромисс между текущим потреблением и снижением будущего климатического ущерба. Снижение выбросов требует затрат сегодня, но уменьшает будущие потери от изменения климата.
\zh{IAM 分为政策评估模型和政策优化模型两类，DICE 主要属于后者。它要寻找储蓄率和气候政策的最优轨迹，使社会福利最大化。核心目标是在当前消费和未来气候损失之间找到最优折中：减排今天要付出成本，但能减少未来的气候损失。}
\zh{IAM 通常分为政策评估模型和政策优化模型，DICE 主要属于后者。它要选择储蓄和气候政策的路径，使社会福利最大化。核心目标是在当前消费和未来气候损失之间找到最优折中：减排今天要付出成本，但能减少未来的气候损失。}

\subsubsection*{2. Целевая функция}
В модели DICE общественное благосостояние задается как дисконтированная сумма полезностей потребления на душу населения:
\[
W=\sum_{t=1}^{T_{\max}}U(c(t),L(t))R(t)\to \max.
\]
Здесь
\[
c(t)=\frac{C(t)}{L(t)}
\]
-- потребление на душу населения, $L(t)$ -- население / труд, $R(t)$ -- дисконтирующий множитель.

Функция полезности имеет вид с постоянной эластичностью предельной полезности:
\[
U(c(t),L(t))=L(t)\frac{c(t)^{1-\alpha}}{1-\alpha},\qquad \alpha\ne 1.
\]
При $\alpha=1$ используется логарифмическая полезность:
\[
U(c,L)=L\ln c.
\]
Дисконтирующий множитель:
\[
R(t)=(1+\rho)^{-t},
\]
где $\rho$ -- чистая норма социального межвременного предпочтения.
\zh{在 DICE 中，社会福利被写成按时间折现的人均消费效用之和。这里 \(c(t)=C(t)/L(t)\) 是人均消费，\(L(t)\) 表示人口或劳动，\(R(t)\) 是贴现因子。效用函数采用恒定相对风险厌恶型；当 \(\alpha=1\) 时则退化为对数效用。贴现因子 \(R(t)=(1+\rho)^{-t}\)，其中 \(\rho\) 是社会贴现率。}

\subsubsection*{3. Экономический блок}
Экономический блок основан на неоклассической модели роста. Основная производственная функция имеет кобб-дouglasовский вид:
\[
Y_{\mathrm{gross}}(t)=A(t)K(t)^\gamma L(t)^{1-\gamma}.
\]
Здесь $A(t)$ -- совокупная факторная производительность, $K(t)$ -- капитал, $L(t)$ -- труд / население, $0<\gamma<1$ -- эластичность выпуска по капиталу. В DICE технологический прогресс $A(t)$ и население $L(t)$ являются экзогенными траекториями, а капитал определяется оптимизацией потребления и инвестиций.

Баланс выпуска:
\[
Y(t)=C(t)+I(t),
\]
где $C(t)$ -- потребление, $I(t)$ -- инвестиции. Динамика капитала:
\[
K(t+1)=(1-\delta)^hK(t)+hI(t),
\]
где $h$ -- длина временного шага. В DICE 2013R используется шаг 5 лет. Часто управление задается через норму сбережений:
\[
s(t)=\frac{I(t)}{Y(t)}.
\]
Тогда
\[
C(t)=(1-s(t))Y(t),\qquad I(t)=s(t)Y(t).
\]
\zh{经济模块基于新古典增长模型，生产函数为 Cobb--Douglas 形式。这里 \(A(t)\) 是全要素生产率，\(K(t)\) 是资本，\(L(t)\) 是劳动或人口，\(\gamma\) 是资本产出弹性。DICE 中技术进步和人口路径是外生给定的，而资本则由消费和投资的最优选择决定。产出分为消费和投资，资本按离散时间积累；常用储蓄率 \(s(t)=I(t)/Y(t)\) 作为控制，于是 \(C(t)=(1-s(t))Y(t)\)、\(I(t)=s(t)Y(t)\)。}

\subsubsection*{4. Климатический ущерб и издержки снижения выбросов}
Валовой выпуск уменьшается из-за климатического ущерба и затрат на снижение выбросов. В DICE используется функция ущерба, зависящая от повышения атмосферной температуры:
\[
\Omega(t)=\psi_1T_{\mathrm{AT}}(t)+\psi_2T_{\mathrm{AT}}^2(t).
\]
В DICE 2013R ущерб упрощенно задается квадратичной функцией температуры.
\zh{总产出会因气候损失和减排成本而下降。DICE 用一个依赖大气温度升高的损失函数 \(\Omega(t)\) 来表示气候损害，在 2013R 版本中它被简化成温度的二次函数。}

Функция затрат на снижение выбросов имеет вид
\[
\Lambda(t)=\theta_1(t)\mu(t)^{\theta_2}.
\]
Здесь $\mu(t)\in[0,1]$ -- доля сокращения выбросов, или emissions control rate. Если $\mu(t)=0$, выбросы не сокращаются; если $\mu(t)=1$, промышленные выбросы полностью устраняются, если это допускается модельными ограничениями. Функция abatement cost является выпуклой по $\mu(t)$, поэтому предельные затраты сокращения выбросов растут при увеличении контроля.

С учетом ущерба и издержек чистый выпуск можно записать как
\[
Y(t)=\frac{1-\Lambda(t)}{1+\Omega(t)}Y_{\mathrm{gross}}(t),
\]
или эквивалентно как выпуск после вычета ущерба и затрат на сокращение выбросов.
\zh{减排成本函数写成 \(\Lambda(t)=\theta_1(t)\mu(t)^{\theta_2}\)。其中 \(\mu(t)\in[0,1]\) 是减排率；若 \(\mu=0\)，则不减排；若 \(\mu=1\)，则工业排放可被完全消除，只要模型约束允许。由于减排成本关于 \(\mu\) 是凸的，减排越多，边际成本越高。考虑损失和成本后，净产出可写成总产出的折减形式。}

\subsubsection*{5. Блок выбросов}
Промышленные выбросы CO$_2$ задаются как функция выпуска, углеродной интенсивности и контроля выбросов:
\[
E_{\mathrm{ind}}(t)=\sigma(t)(1-\mu(t))Y_{\mathrm{gross}}(t).
\]
Здесь $\sigma(t)$ -- углеродная интенсивность выпуска, $\mu(t)$ -- доля сокращения выбросов. Полные выбросы включают промышленные выбросы и выбросы от землепользования:
\[
E(t)=E_{\mathrm{ind}}(t)+E_{\mathrm{land}}(t).
\]
\zh{工业 CO$_2$ 排放由总产出、碳强度和减排控制率共同决定；总排放还要加上土地利用产生的排放。}

\subsubsection*{6. Углеродный цикл}
DICE использует трехрезервуарную модель углеродного цикла:
\[
M_{\mathrm{AT}}(t),\quad M_{\mathrm{UP}}(t),\quad M_{\mathrm{LO}}(t),
\]
где $M_{\mathrm{AT}}$ -- углерод в атмосфере, $M_{\mathrm{UP}}$ -- углерод в верхнем океане и биосфере, $M_{\mathrm{LO}}$ -- углерод в глубоком океане.

Уравнения имеют вид линейной дискретной системы:
\[
M_{\mathrm{AT}}(t+1)=\phi_{11}M_{\mathrm{AT}}(t)+\phi_{21}M_{\mathrm{UP}}(t)+E(t),
\]
\[
M_{\mathrm{UP}}(t+1)=\phi_{12}M_{\mathrm{AT}}(t)+\phi_{22}M_{\mathrm{UP}}(t)+\phi_{32}M_{\mathrm{LO}}(t),
\]
\[
M_{\mathrm{LO}}(t+1)=\phi_{23}M_{\mathrm{UP}}(t)+\phi_{33}M_{\mathrm{LO}}(t).
\]
Смысл: выбросы сначала попадают в атмосферу, затем углерод перераспределяется между атмосферой, верхним океаном, биосферой и глубоким океаном.
\zh{DICE 采用三储库碳循环模型：大气、上层海洋/生物圈和深海。离散方程表示碳先进入大气，再在几个储库之间逐步重新分配。}

\subsubsection*{7. Радиационное воздействие}
Концентрация CO$_2$ влияет на радиационный баланс Земли. Радиационное воздействие задается формулой
\[
F(t)=\eta\log_2\left(\frac{M_{\mathrm{AT}}(t)}{M_{\mathrm{AT}}^{1750}}\right)+F_{\mathrm{EX}}(t),
\]
где $M_{\mathrm{AT}}^{1750}$ -- доиндустриальный уровень углерода в атмосфере, $F_{\mathrm{EX}}(t)$ -- внешнее радиационное воздействие от других парниковых газов и факторов.
\zh{二氧化碳浓度会改变地球辐射平衡。辐射强迫用对数式表示，\(M_{\mathrm{AT}}^{1750}\) 是工业化前的大气碳水平，\(F_{\mathrm{EX}}(t)\) 是其他温室气体和外部因素造成的附加辐射强迫。}

\subsubsection*{8. Температурный блок}
Климатическая система описывается двумя температурными переменными:
\[
T_{\mathrm{AT}}(t)
\]
-- средняя температура атмосферы / поверхности,
\[
T_{\mathrm{LO}}(t)
\]
-- температура глубокого океана.

Типичная форма уравнений:
\[
T_{\mathrm{AT}}(t+1)=T_{\mathrm{AT}}(t)+\xi_1\left[F(t)-\xi_2T_{\mathrm{AT}}(t)-\xi_3(T_{\mathrm{AT}}(t)-T_{\mathrm{LO}}(t))\right],
\]
\[
T_{\mathrm{LO}}(t+1)=T_{\mathrm{LO}}(t)+\xi_4(T_{\mathrm{AT}}(t)-T_{\mathrm{LO}}(t)).
\]
Первое уравнение описывает нагрев атмосферы под воздействием radiative forcing. Второе описывает медленный перенос тепла в глубокий океан.

\subsubsection*{9. Управления в модели DICE}
Основные управляющие переменные:
\[
s(t)
\]
-- норма сбережений, или investment rate,
\[
\mu(t)
\]
-- норма сокращения выбросов, или emissions control rate.

Иногда вместо $\mu(t)$ можно использовать углеродный налог или carbon price, поскольку в оптимуме carbon price равен предельным издержкам выбросов / social cost of carbon.

\subsubsection*{10. Задача оптимизации DICE}
В общей форме задача DICE:
\[
\max_{\{s(t),\mu(t)\}_{t=1}^{T}}\sum_{t=1}^{T}R(t)L(t)\frac{c(t)^{1-\alpha}}{1-\alpha}
\]
при ограничениях
\[
K(t+1)=(1-\delta)^hK(t)+hI(t),
\]
\[
Y_{\mathrm{gross}}(t)=A(t)K(t)^\gamma L(t)^{1-\gamma},
\]
\[
Y(t)=\frac{1-\Lambda(t)}{1+\Omega(t)}Y_{\mathrm{gross}}(t),
\]
\[
C(t)=Y(t)-I(t),
\]
\[
E_{\mathrm{ind}}(t)=\sigma(t)(1-\mu(t))Y_{\mathrm{gross}}(t),
\]
\[
E(t)=E_{\mathrm{ind}}(t)+E_{\mathrm{land}}(t),
\]
\[
M(t+1)=\Phi M(t)+BE(t),
\]
\[
F(t)=\eta\log_2\frac{M_{\mathrm{AT}}(t)}{M_{\mathrm{AT}}^{1750}}+F_{\mathrm{EX}}(t),
\]
\[
T(t+1)=\Psi T(t)+\Xi F(t),
\]
\[
0\le s(t)\le 1,\qquad 0\le \mu(t)\le \bar\mu(t).
\]
Здесь $M(t)$ -- вектор запасов углерода, $T(t)$ -- вектор температур.

\subsubsection*{11. Подходы к решению}
DICE является дискретной нелинейной оптимизационной задачей большой размерности. На практике строится конечномерная задача
\[
\max_{s_1,\ldots,s_T,\mu_1,\ldots,\mu_T}W
\]
при системе нелинейных ограничений. Обычно используются GAMS / NLP solver / CONOPT, а также Excel Solver.

Теоретически можно записать дискретный аналог принципа максимума или условия Каруша--Куна--Таккера. Состояния:
\[
K,\ M_{\mathrm{AT}},M_{\mathrm{UP}},M_{\mathrm{LO}},T_{\mathrm{AT}},T_{\mathrm{LO}}.
\]
Управления:
\[
s,\ \mu.
\]
Сопряженные переменные имеют экономическую интерпретацию: теневая цена капитала, теневая цена атмосферного углерода, теневая цена температуры, social cost of carbon.

Формально можно ввести функцию ценности
\[
V(K,M,T,t)
\]
и записать дискретное уравнение Беллмана:
\[
V_t(K,M,T)=\max_{s,\mu}\left\{U(c,L)R(t)+V_{t+1}(K',M',T')\right\}.
\]
Однако из-за высокой размерности состояния этот подход сложен численно. Поэтому в DICE обычно применяется прямое нелинейное программирование.

DICE также используется для сравнения сценариев: baseline, optimal, temperature-limited, low discounting, Copenhagen Accord.

\subsubsection*{12. Экономический смысл решения}
Оптимальное решение DICE определяет $s^*(t)$ -- сколько выпуска направлять на инвестиции; $\mu^*(t)$ -- какую долю выбросов сокращать; $C^*(t)$ -- траекторию потребления; $K^*(t)$ -- траекторию капитала; $E^*(t)$ -- траекторию выбросов; $T_{\mathrm{AT}}^*(t)$ -- траекторию глобальной температуры; SCC$(t)$ -- social cost of carbon.

Экономический компромисс: высокая $\mu(t)$ снижает будущий климатический ущерб, но требует текущих затрат; высокая $s(t)$ снижает текущее потребление, но повышает будущий капитал; высокая текущая эмиссия увеличивает сегодняшнее производство, но ухудшает климат в будущем.

\qpart{2}{Модель DICE}

\subsection{Пояснение на китайском}


\[
Y_{\mathrm{gross}}=AK^\gamma L^{1-\gamma},
\qquad
K(t+1)=(1-\delta)^hK(t)+hI(t).
\]
\[
s(t)=I(t)/Y(t).
\]

\[
E_{\mathrm{ind}}=\sigma(t)(1-\mu(t))Y_{\mathrm{gross}}(t),
\]

\[
\Omega=\psi_1T_{\mathrm{AT}}+\psi_2T_{\mathrm{AT}}^2,
\qquad
\Lambda=\theta_1(t)\mu(t)^{\theta_2}.
\]
\[
W=\sum_t R(t)L(t)\frac{c(t)^{1-\alpha}}{1-\alpha}.
\]

\qpart{3}{Задачи оптимизации и подходы к решению}

\blocktitle{I}{Неантагонистические дифференциальные игры, бесконечный горизонт и конкурентная реклама}{Билеты 27--29}
\zh{I 板块：非对抗性微分博弈、无限时域与竞争广告。第27--29题。}
Замечание об источниках и надежности. Для билетов 27--28 используются стандартные результаты теории оптимального управления и дифференциальных игр: принцип максимума Понтрягина для программных стратегий и уравнения Беллмана--Гамильтона--Якоби для позиционных стратегий. Для бесконечного горизонта используются материалы Асеева и Асеева--Кряжимского по бесконечногоризонтному оптимальному управлению, current-value переменным и условиям трансверсальности на бесконечности. Для билета 29 конкретная модель конкурентной рекламы берется по студенческому конспекту и краткой версии 2026 года; текст Sorger (1989) в текущем наборе материалов отсутствует, поэтому специальные утверждения, зависящие от оригинальной статьи, следует считать требующими последующей сверки. Формулы PMP и HJB ниже выводятся непосредственно из указанной модели.
\zh{来源与可靠性说明。第27--28题使用最优控制与微分博弈中的标准结果：程序型策略对应 Pontryagin 极大值原理，位置型策略对应 Bellman--Hamilton--Jacobi 方程。无限时域部分采用 Aseev 与 Aseev--Krjazhimskii 关于无限时域最优控制、current-value 变量和无穷远横截条件的材料。第29题的具体竞争广告模型取自学生笔记和 2026 年简版材料；由于当前材料中没有 Sorger（1989）的原文，凡依赖原论文的特殊结论都应视为待进一步核对。下文中的 PMP 和 HJB 公式是直接根据该模型推导的。}

\ticket{27}{Неантагонистические дифференциальные игры. Равновесие по Нэшу в программных и позиционных стратегиях.}{非对抗性微分博弈。程序型与位置型策略中的 Nash 均衡。}


\qpart{1}{Неантагонистические дифференциальные игры}{非对抗性微分博弈}


Рассмотрим неантагонистическую дифференциальную игру $N$ лиц на конечном промежутке времени $[t_0,T]$. Состояние системы обозначим
\[
 x(t)\in\R^n,
\]
управление $i$-го игрока:
\[
 u_i(t)\in U_i\subset\R^{m_i},\qquad i=1,\ldots,N.
\]
Динамика задается системой
\[
 \dot x(t)=f(t,x(t),u_1(t),\ldots,u_N(t)),\qquad x(t_0)=x_0.
\]
Функционал выигрыша $i$-го игрока имеет вид
\[
 J_i(u_1,\ldots,u_N)=\int_{t_0}^{T}F_i(t,x(t),u_1(t),\ldots,u_N(t))\,dt+q_i(x(T)).
\]
Здесь $F_i$ -- текущий выигрыш, $q_i$ -- терминальный выигрыш.

Игра называется неантагонистической, если выигрыши игроков не сводятся к одной функции с противоположными знаками. В частности, вообще говоря,
\[
 J_i+J_j\ne0.
\]
Это не означает отсутствия конфликта: интересы игроков могут частично совпадать и частично противоречить друг другу.
\zh{考虑在有限时区间 \([t_0,T]\) 上的 \(N\) 人非对抗性微分博弈。状态为 \(x(t)\)，第 \(i\) 个玩家的控制为 \(u_i(t)\)，系统动力学由 \(\dot x=f(t,x,u_1,\ldots,u_N)\) 给出。第 \(i\) 个玩家的收益泛函由当前收益 \(F_i\) 和终端收益 \(q_i\) 组成。所谓非对抗性，是指各玩家的收益并不简单地互为相反数；也就是说，他们之间可以既有部分共同利益，也有部分冲突。}

\subsubsection{Равновесие по Нэшу}

Ситуация
\[
 u^*=(u_1^*,\ldots,u_N^*)
\]
называется равновесием по Нэшу, если ни один игрок не может увеличить свой выигрыш односторонним отклонением при фиксированных стратегиях остальных игроков:
\[
 J_i(u_i,u_{-i}^*)\le J_i(u_i^*,u_{-i}^*)
 \qquad \forall u_i\in\mathcal U_i,\quad i=1,\ldots,N.
\]
Здесь
\[
 u_{-i}^*=(u_1^*,\ldots,u_{i-1}^*,u_{i+1}^*,\ldots,u_N^*).
\]
\zh{Nash 均衡就是：在其他玩家策略固定时，任何一个玩家都不能通过单边偏离提高自己的收益。}

\subsubsection{Программные стратегии}

Программная стратегия -- это управление, заданное как функция времени:
\[
 u_i=u_i(t).
\]
Она может зависеть от начальной позиции $(t_0,x_0)$, но не зависит от текущего состояния $x(t)$ в процессе игры. Игрок заранее выбирает весь план действий
\[
 u_i(\cdot):[t_0,T]\to U_i.
\]

Равновесие по Нэшу в программных стратегиях определяется неравенствами
\[
 J_i(u_i,u_{-i}^*)\le J_i(u_i^*,u_{-i}^*)
 \qquad \forall u_i(\cdot),\quad i=1,\ldots,N.
\]

Для вывода необходимых условий для каждого игрока фиксируются равновесные стратегии остальных игроков и рассматривается обычная задача оптимального управления по $u_i$.
\zh{程序型策略指的是只依赖时间的控制 \(u_i(t)\)。它可以依赖初始时刻和初始状态，但在博弈进行时不依赖当前状态 \(x(t)\)。程序型 Nash 均衡是指对每个玩家来说，只要其他玩家保持均衡程序，他的最优响应就是自己的均衡程序。推导必要条件时，把其他玩家的均衡策略固定下来，对单个玩家就变成普通的最优控制问题。}

Гамильтониан $i$-го игрока:
\[
 H_i(t,x,u_1,\ldots,u_N,\psi_i)=F_i(t,x,u_1,\ldots,u_N)+\langle\psi_i,f(t,x,u_1,\ldots,u_N)\rangle.
\]
У каждого игрока свой функционал, поэтому у каждого игрока своя сопряженная переменная $\psi_i(t)$.

Если $u^*$ является программным равновесием по Нэшу и выполнены стандартные условия регулярности, то для каждого $i$ существует $\psi_i(t)$ такое, что вдоль равновесной траектории $x^*(t)$:
\[
 \dot x^*(t)=f(t,x^*(t),u_1^*(t),\ldots,u_N^*(t)),
\]
\[
 \dot\psi_i(t)=-\frac{\partial H_i}{\partial x}(t,x^*(t),u_1^*(t),\ldots,u_N^*(t),\psi_i(t)),
\]
\[
 \psi_i(T)=q_{i,x}(x^*(T)).
\]
Условие максимума:
\[
 u_i^*(t)\in\argmax_{v_i\in U_i}H_i(t,x^*(t),u_1^*(t),\ldots,u_{i-1}^*(t),v_i,u_{i+1}^*(t),\ldots,u_N^*(t),\psi_i(t)).
\]
Это условие означает: каждый игрок максимизирует свой гамильтониан только по своему управлению $u_i$, считая управления остальных игроков фиксированными. В данной записи используется конвенция максимизации выигрышей; при минимизации затрат знак условий меняется.
\zh{于是得到程序型均衡的 PMP 条件：每个玩家只对自己的控制最大化自己的 Hamilton 函数，其余玩家的控制被视为固定。这里采用的是收益最大化约定；若改写成成本最小化，符号要相应改变。}

\subsubsection{Позиционные стратегии}

Позиционная стратегия -- это стратегия обратной связи:
\[
 u_i=\varphi_i(t,x).
\]
Она зависит от времени и текущего состояния системы. Если игроки используют позиционные стратегии $\varphi^*=(\varphi_1^*,\ldots,\varphi_N^*)$, то равновесная траектория определяется замкнутой системой
\[
 \dot x^*(t)=f(t,x^*(t),\varphi_1^*(t,x^*(t)),\ldots,\varphi_N^*(t,x^*(t))).
\]

Позиционное равновесие по Нэшу означает, что ни один игрок не может увеличить свой выигрыш, заменив свою обратную связь $\varphi_i^*$ на другую допустимую стратегию $\varphi_i$, если остальные игроки сохраняют $\varphi_{-i}^*$.

Для позиционных стратегий особенно важно понятие состоятельного равновесия. Оно означает, что равновесие сохраняется не только из начальной позиции $(t_0,x_0)$, но и из любой промежуточной позиции $(t,x)$. Поэтому оно естественно описывается методом динамического программирования.
\zh{位置型策略是反馈形式 \(u_i=\varphi_i(t,x)\)，依赖时间和当前状态。若所有玩家采用位置型策略，则闭环系统决定均衡轨线。位置型 Nash 均衡要求：任一玩家都不能仅通过改写自己的反馈而提高收益。对于位置型策略，尤其重要的是“可持续均衡”或“sustainable equilibrium”概念，它要求均衡不仅从初始点成立，而且从任意中间状态继续成立，因此自然地用动态规划来描述。}

Пусть существуют функции выигрыша, или функции Беллмана,
\[
 V_i(t,x),\qquad i=1,\ldots,N.
\]
Терминальные условия:
\[
 V_i(T,x)=q_i(x).
\]
Тогда для состоятельного позиционного равновесия должны выполняться уравнения Гамильтона--Якоби:
\[
 -V_{i,t}(t,x)=\max_{v_i\in U_i}\Bigl\{F_i(t,x,\varphi_{-i}^*(t,x),v_i)+V_{i,x}(t,x)f(t,x,\varphi_{-i}^*(t,x),v_i)\Bigr\}.
\]
Равновесная стратегия $i$-го игрока достигает этот максимум:
\[
 \varphi_i^*(t,x)\in\argmax_{v_i\in U_i}\Bigl\{F_i(t,x,\varphi_{-i}^*(t,x),v_i)+V_{i,x}(t,x)f(t,x,\varphi_{-i}^*(t,x),v_i)\Bigr\}.
\]

\subsubsection{Итог}

В программных стратегиях Nash equilibrium задается заранее выбранными управлениями $u_i(t)$ и проверяется через односторонние отклонения. Необходимые условия имеют вид принципа максимума Понтрягина для каждого игрока. В позиционных стратегиях $u_i=\varphi_i(t,x)$; состоятельное позиционное равновесие описывается системой HJB для всех игроков.

\qpart{2}{Равновесие по Нэшу в программных стратегиях}{程序型策略中的 Nash 均衡}

\qpart{3}{Равновесие по Нэшу в позиционных стратегиях}{位置型策略中的 Nash 均衡}

\newpage
\ticket{28}{Дифференциальные игры с бесконечной продолжительностью. Равновесие по Нэшу в программных и позиционных стратегиях.}{无限时域微分博弈。程序型与位置型策略中的 Nash 均衡。}


\qpart{1}{Дифференциальные игры с бесконечной продолжительностью}{无限时域微分博弈}


Рассмотрим неантагонистическую дифференциальную игру $N$ лиц на бесконечном промежутке времени
\[
 [t_0,+\infty).
\]
Состояние системы:
\[
 x(t)\in X\subset\R^n.
\]
Управление $i$-го игрока:
\[
 u_i(t)\in U_i\subset\R^{m_i},\qquad i=1,\ldots,N.
\]
Динамика:
\[
 \dot x(t)=f(t,x(t),u_1(t),\ldots,u_N(t)),\qquad x(t_0)=x_0.
\]
Выигрыш $i$-го игрока:
\[
 J_i(u_1,\ldots,u_N)=\int_{t_0}^{+\infty}e^{-r(t-t_0)}F_i(t,x(t),u_1(t),\ldots,u_N(t))\,dt,
\]
где $r>0$ -- коэффициент дисконтирования.

Дисконтирование отражает меньшую текущую ценность будущих выигрышей. Однако само по себе условие $r>0$ не гарантирует сходимость функционала при любом $F_i$. Обычно дополнительно предполагается ограниченность или допустимый рост $F_i$, чтобы интеграл был конечным.
\zh{这里研究定义在无穷区间 \([t_0,+\infty)\) 上的非对抗性微分博弈。收益函数带有折现因子 \(e^{-r(t-t_0)}\)，所以未来收益的现值会逐渐变小。折现率 \(r>0\) 反映了“未来收益的当前价值更低”；但仅有 \(r>0\) 并不能保证泛函一定收敛，还要对 \(F_i\) 的有界性或增长速度作额外假设。}

\subsubsection{Равновесие в программных стратегиях}

Программная стратегия -- это управление $u_i=u_i(t)$, зависящее от времени и начальной позиции, но не зависящее от текущего состояния. Набор
\[
 u^*=(u_1^*,\ldots,u_N^*)
\]
называется равновесием по Нэшу в программных стратегиях, если
\[
 J_i(u_i,u_{-i}^*)\le J_i(u_i^*,u_{-i}^*)\qquad \forall u_i(\cdot),\quad i=1,\ldots,N.
\]

Для каждого игрока фиксируются стратегии остальных игроков и получается задача оптимального управления на бесконечном горизонте.
\zh{程序型均衡与有限时域完全类似，只是时域变成无穷长。固定其他玩家的程序后，每个玩家都要解一个无限时域最优控制问题。}

Present-value гамильтониан:
\[
 \mathcal H_i(t,x,u,\psi_i)=e^{-r(t-t_0)}F_i(t,x,u_1,\ldots,u_N)+\langle\psi_i,f(t,x,u_1,\ldots,u_N)\rangle.
\]
Необходимые условия PMP:
\[
 \dot x^*(t)=f(t,x^*(t),u_1^*(t),\ldots,u_N^*(t)),
\]
\[
 \dot\psi_i(t)=-\frac{\partial\mathcal H_i}{\partial x}(t,x^*(t),u^*(t),\psi_i(t)),
\]
\[
 u_i^*(t)\in\argmax_{v_i\in U_i}\mathcal H_i(t,x^*(t),u_1^*(t),\ldots,u_{i-1}^*(t),v_i,u_{i+1}^*(t),\ldots,u_N^*(t),\psi_i(t)).
\]

На бесконечном горизонте нет терминального условия вида $\psi_i(T)=q_{i,x}(x(T))$. Вместо него используются условия трансверсальности на бесконечности. Типичная форма:
\[
 \lim_{t\to+\infty}\langle\psi_i(t),x^*(t)\rangle=0,
\]
или более точная модельная форма, например условие для $\langle \psi_i(t),x(t)-x^*(t)\rangle$.
\zh{现值形式的 Hamilton 函数写成 \(\mathcal H_i=e^{-r(t-t_0)}F_i+\langle\psi_i,f\rangle\)。相应的 PMP 给出状态方程、伴随方程、最大化条件和无穷远横截条件。横截条件表示贴现后的影子价值在无穷远处趋于零。}

\subsubsection{Current-value переменные}

В экономических задачах удобно перейти от present-value сопряженной переменной $\psi_i(t)$ к current-value переменной $p_i(t)$:
\[
 p_i(t)=e^{r(t-t_0)}\psi_i(t),\qquad \psi_i(t)=e^{-r(t-t_0)}p_i(t).
\]
Current-value гамильтониан:
\[
 H_i(t,x,u,p_i)=F_i(t,x,u_1,\ldots,u_N)+\langle p_i,f(t,x,u_1,\ldots,u_N)\rangle.
\]
Тогда условия принимают вид
\[
 u_i^*(t)\in\argmax_{v_i\in U_i}H_i(t,x^*(t),u_{-i}^*(t),v_i,p_i(t)),
\]
\[
 \dot p_i(t)=r p_i(t)-\frac{\partial H_i}{\partial x}(t,x^*(t),u^*(t),p_i(t)).
\]
Член $r p_i(t)$ является характерным отличием current-value записи бесконечногоризонтной задачи.
\zh{在经济模型里，常把现值伴随变量 \(\psi_i\) 换成 current-value 变量 \(p_i\)。这样 Hamilton 函数不再带贴现因子，伴随方程中则多出一项 \(r p_i(t)\)。这是无限时域问题的一个典型特征。}

\subsubsection{Позиционные стратегии}

Позиционная стратегия имеет вид
\[
 u_i=\varphi_i(t,x)
\]
или, в автономной стационарной задаче,
\[
 u_i=\varphi_i(x).
\]
Набор $\varphi^*=(\varphi_1^*,\ldots,\varphi_N^*)$ называется равновесием по Нэшу в позиционных стратегиях, если ни один игрок не может увеличить свой функционал, заменив только свою обратную связь $\varphi_i^*$, при условии что остальные игроки сохраняют $\varphi_{-i}^*$.

Равновесная траектория задается замкнутой системой
\[
 \dot x^*(t)=f(t,x^*(t),\varphi_1^*(t,x^*(t)),\ldots,\varphi_N^*(t,x^*(t))).
\]

Для состоятельного позиционного равновесия вводятся функции ценности игроков
\[
 V_i(t,x),\qquad i=1,\ldots,N.
\]
\zh{位置型策略是反馈形式 \(u_i=\varphi_i(t,x)\)，依赖时间和当前状态。若所有玩家采用位置型策略，则闭环系统决定均衡轨线。位置型 Nash 均衡要求：任一玩家都不能仅通过改写自己的反馈而提高收益。对于位置型策略，尤其重要的是“可持续均衡”概念，它要求均衡不仅从初始点成立，而且从任意中间状态继续成立，因此自然地用动态规划来描述。}
В автономном стационарном случае обычно ищут $V_i(x)$ и $\varphi_i^*(x)$. Стационарная система Беллмана имеет вид
\[
 rV_i(x)=\max_{v_i\in U_i}\Bigl\{F_i(x,\varphi_{-i}^*(x),v_i)+V_{i,x}(x)f(x,\varphi_{-i}^*(x),v_i)\Bigr\}.
\]
Равновесная стратегия достигает максимум:
\[
 \varphi_i^*(x)\in\argmax_{v_i\in U_i}\Bigl\{F_i(x,\varphi_{-i}^*(x),v_i)+V_{i,x}(x)f(x,\varphi_{-i}^*(x),v_i)\Bigr\}.
\]

\subsubsection{Отличия от конечного горизонта}

\begin{enumerate}
\item Нет конечного момента $T$ и терминального выигрыша $q_i(x(T))$ в стандартной формулировке.
\item Требуется сходимость бесконечного интеграла $J_i$.
\item В current-value PMP появляется член $r p_i$.
\item Вместо конечного терминального условия появляются условия трансверсальности на бесконечности.
\item В стационарном HJB возникает член $rV_i$.
\end{enumerate}

\qpart{2}{Равновесие по Нэшу в программных стратегиях}{程序型策略中的 Nash 均衡}

\qpart{3}{Равновесие по Нэшу в позиционных стратегиях}{位置型策略中的 Nash 均衡}

\newpage
\ticket{29}{Модель конкурентной рекламы с двумя участниками.}


\qpart{1}{Модель конкурентной рекламы с двумя участниками}


Рассматривается рынок фиксированного размера, на котором конкурируют две фирмы. Переменная
\[
 x(t)\in[0,1]
\]
обозначает долю рынка первой фирмы в момент времени $t$. Тогда $1-x(t)$ -- доля рынка второй фирмы.

Управления игроков:
\[
 u_1(t)\ge0,
 \qquad
 u_2(t)\ge0
\]
-- расходы фирм на рекламу.

Динамика доли рынка первой фирмы задается уравнением
\[
 \dot x(t)=u_1(t)\sqrt{1-x(t)}-u_2(t)\sqrt{x(t)},\qquad x(0)=x_0\in[0,1].
\]
Смысл уравнения:
\begin{itemize}
\item реклама первой фирмы увеличивает ее долю рынка;
\item чем меньше доля первой фирмы, тем больше оставшийся рынок $1-x$, на который она может воздействовать;
\item реклама второй фирмы уменьшает долю первой фирмы;
\item эффект рекламы второй фирмы зависит от текущей доли первой фирмы $x$.
\end{itemize}

Если $u_1,u_2\ge0$, то отрезок $[0,1]$ является естественным фазовым множеством: при $x=0$ имеем $\dot x=u_1\ge0$, а при $x=1$ имеем $\dot x=-u_2\le0$.

\subsubsection{Функционалы выигрыша}

Первая фирма максимизирует
\[
 J_1=\int_0^T\left(q_1x(t)-\frac{c_1}{2}u_1^2(t)\right)e^{-rt}\,dt+e^{-rT}S_1x(T).
\]
Вторая фирма максимизирует
\[
 J_2=\int_0^T\left(q_2(1-x(t))-\frac{c_2}{2}u_2^2(t)\right)e^{-rt}\,dt+e^{-rT}S_2(1-x(T)).
\]
Здесь
\[
 r,q_i,c_i,S_i>0,
 \qquad i=1,2.
\]
Параметры имеют экономический смысл: $q_i$ -- доходность рыночной доли фирмы $i$, $c_i$ -- коэффициент стоимости рекламы, $S_i$ -- терминальная ценность рыночной доли, $r$ -- ставка дисконтирования.

\subsubsection{Программное равновесие}

Программная стратегия -- это управление $u_i=u_i(t)$, не зависящее от текущего состояния. Набор $(u_1^*,u_2^*)$ называется равновесием по Нэшу в программных стратегиях, если
\[
 J_1(u_1,u_2^*)\le J_1(u_1^*,u_2^*)\qquad \forall u_1,
\]
\[
 J_2(u_1^*,u_2)\le J_2(u_1^*,u_2^*)\qquad \forall u_2.
\]

Гамильтониан первой фирмы:
\[
 H_1=e^{-rt}\left(q_1x-\frac{c_1}{2}u_1^2\right)+\Lambda_1\left(u_1\sqrt{1-x}-u_2\sqrt{x}\right).
\]
Гамильтониан второй фирмы:
\[
 H_2=e^{-rt}\left(q_2(1-x)-\frac{c_2}{2}u_2^2\right)+\Lambda_2\left(u_1\sqrt{1-x}-u_2\sqrt{x}\right).
\]

Условия максимума:
\[
 \frac{\partial H_1}{\partial u_1}=-c_1e^{-rt}u_1+\Lambda_1\sqrt{1-x}=0,
\]
откуда
\[
 u_1^*(t)=\frac{\Lambda_1(t)}{c_1}e^{rt}\sqrt{1-x^*(t)}.
\]
Для второго игрока:
\[
 \frac{\partial H_2}{\partial u_2}=-c_2e^{-rt}u_2-\Lambda_2\sqrt{x}=0,
\]
поэтому
\[
 u_2^*(t)=-\frac{\Lambda_2(t)}{c_2}e^{rt}\sqrt{x^*(t)}.
\]
Если явно накладываются ограничения $u_i(t)\ge0$, корректнее использовать проекционную форму:
\[
 u_1^*(t)=\max\left\{0,\frac{\Lambda_1(t)}{c_1}e^{rt}\sqrt{1-x^*(t)}\right\},
\]
\[
 u_2^*(t)=\max\left\{0,-\frac{\Lambda_2(t)}{c_2}e^{rt}\sqrt{x^*(t)}\right\}.
\]

Сопряженные уравнения:
\[
 \dot\Lambda_1=-q_1e^{-rt}+\Lambda_1\left(\frac{u_1^*}{2\sqrt{1-x^*}}+\frac{u_2^*}{2\sqrt{x^*}}\right),
\]
\[
 \dot\Lambda_2=q_2e^{-rt}+\Lambda_2\left(\frac{u_1^*}{2\sqrt{1-x^*}}+\frac{u_2^*}{2\sqrt{x^*}}\right).
\]
Терминальные условия:
\[
 \Lambda_1(T)=e^{-rT}S_1,
\]
\[
 \Lambda_2(T)=-e^{-rT}S_2.
\]

Вместе с уравнением состояния и условием $x(0)=x_0$ это дает краевую задачу для программного равновесия.

\subsubsection{Позиционное равновесие}

Позиционная стратегия имеет вид
\[
 u_i=\varphi_i(t,x).
\]
Ищем состоятельное позиционное равновесие через функции Беллмана. Удобно использовать линейный по $x$ вид:
\[
 V_1(t,x)=e^{-rt}(A_1(t)x+B_1(t)),
\]
\[
 V_2(t,x)=e^{-rt}(A_2(t)(1-x)+B_2(t)).
\]
Тогда
\[
 V_{1x}=e^{-rt}A_1(t),\qquad V_{2x}=-e^{-rt}A_2(t).
\]

Из HJB-условия максимума по $u_1$ получаем
\[
 \varphi_1^*(t,x)=\frac{A_1(t)}{c_1}\sqrt{1-x}.
\]
Для второго игрока:
\[
 \varphi_2^*(t,x)=\frac{A_2(t)}{c_2}\sqrt{x}.
\]

Подстановка обратных связей в уравнения Беллмана дает систему для коэффициентов:
\[
 \dot A_1=rA_1-q_1+\frac{A_1^2}{2c_1}+\frac{A_1A_2}{c_2},
\]
\[
 \dot B_1=rB_1-\frac{A_1^2}{2c_1},
\]
\[
 \dot A_2=rA_2-q_2+\frac{A_1A_2}{c_1}+\frac{A_2^2}{2c_2},
\]
\[
 \dot B_2=rB_2-\frac{A_2^2}{2c_2}.
\]
Терминальные условия:
\[
 A_1(T)=S_1,\\quad B_1(T)=0,
\]
\[
 A_2(T)=S_2,
\quad B_2(T)=0.
\]

Равновесная динамика при позиционных стратегиях:
\[
 \dot x=\varphi_1^*(t,x)\sqrt{1-x}-\varphi_2^*(t,x)\sqrt{x},
\]
то есть
\[
 \dot x=\frac{A_1(t)}{c_1}(1-x)-\frac{A_2(t)}{c_2}x.
\]
Это линейное уравнение по $x$.

\subsubsection{Экономический смысл}

Если первая фирма увеличивает рекламу $u_1$, ее доля рынка растет. Но эффект рекламы уменьшается при $x\to1$, потому что оставшийся рынок $1-x$ становится мал. Если вторая фирма увеличивает рекламу $u_2$, доля первой фирмы падает. Квадратичные затраты $\frac{c_i}{2}u_i^2$ выражают возрастающую предельную стоимость рекламной активности.

Программное равновесие строится через PMP и зависит от начальной траектории. Позиционное равновесие задает обратную связь $\varphi_i(t,x)$ и потому реагирует на изменение состояния в процессе игры.

\end{document}
